% Amélioration de la production animale par l'exploitation des insectes — SVT Tle D — Cours signé OnBuch+
% Programme officiel MINESEC. Compiler avec : tectonic N18-entomophagie-production-animale.tex
\newcommand{\DOCMATIERE}{SVT}
\newcommand{\DOCNIVEAU}{Tle D}
\newcommand{\DOCMODULE}{Module IV : La Biotechnologie — Entomophagie}
\newcommand{\DOCLECON}{N18}
\newcommand{\DOCTITRE}{Amélioration de la production animale par l'exploitation des insectes}
% OnBuch+ — préambule « cours signés » ENRICHI (compilé avec tectonic / XeLaTeX)
% Système visuel « Pop » d'OnBuch+ : fond crème (jamais blanc), bordures encre
% épaisses, ombres « dures » décalées, titres Archivo Black en capitales.
% Version MATHÉMATIQUES : graphes (pgfplots/tikz), boîtes pédagogiques
% (définition, propriété, méthode, attention, le savais-tu, exercice résolu,
% exercice, TP, bilan), page de garde et signature OnBuch+ ancrée en bas.
%
% Macros attendues dans main.tex AVANT \input{preamble} :
%   \DOCMATIERE  \DOCNIVEAU  \DOCMODULE  \DOCLECON (numéro)  \DOCTITRE
\documentclass[11pt]{article}

\usepackage{fontspec}
\usepackage[french]{babel}
\usepackage[a4paper,margin=2cm,top=3.4cm,bottom=2.8cm,headheight=1.4cm,footskip=1.3cm]{geometry}
\usepackage{amsmath,amssymb,amsfonts}
\usepackage{mathtools} % \xmapsto, \coloneqq... souvent utilisés par les modèles
\usepackage{cancel} % simplifications barrées (bilans de forces)
% \abs{z} (module, valeur absolue) et \norm{u} (norme), utilisés par les modèles
\providecommand{\abs}{}\let\abs\relax\DeclarePairedDelimiter{\abs}{\lvert}{\rvert}
\providecommand{\norm}{}\let\norm\relax\DeclarePairedDelimiter{\norm}{\lVert}{\rVert}
\usepackage[table]{xcolor} % [table] : fournit aussi \rowcolors (alternance de lignes)
\usepackage{pagecolor}
\usepackage{tikz}
\usetikzlibrary{arrows.meta,positioning,calc,shapes.geometric,shapes.misc,shapes.arrows,decorations.pathmorphing,decorations.pathreplacing,decorations.markings,patterns,shadows,mindmap,trees,fit,backgrounds,matrix,intersections,angles,quotes}
\usepackage{pgfplots}
\usepackage[version=4]{mhchem} % équations nucléaires \ce{^{235}_{92}U}
\usepackage[european,siunitx]{circuitikz} % schémas électriques
\pgfplotsset{compat=1.18}
\usepgfplotslibrary{fillbetween,groupplots} % aires entre courbes (calcul intégral)
\usepackage{enumitem}
\usepackage{fancyhdr}
% [most] charge toutes les bibliothèques tcolorbox (listings, minted, raster,
% documentation...) et gonfle inutilement la mémoire TeX sur un document long
% (~300 boîtes) : on ne charge que ce qui est réellement utilisé.
\usepackage[skins,breakable]{tcolorbox}
\usepackage{graphicx}
\usepackage{colortbl}
\usepackage{booktabs}
\usepackage{tabularx}
\usepackage{diagbox} % case d'en-tête diagonale des tableaux à double entrée
% Colonnes extensibles centrée / gauche / droite (C, L, R), souvent utilisées par les modèles
\newcolumntype{C}{>{\centering\arraybackslash}X}
\newcolumntype{L}{>{\raggedright\arraybackslash}X}
\newcolumntype{R}{>{\raggedleft\arraybackslash}X}
\usepackage{array}
\usepackage{multicol}
\usepackage{siunitx}
% parse-numbers=false : accepte aussi \SI{2,5 \times 10^{-3}}{...}
\sisetup{locale=FR,output-decimal-marker={,},group-minimum-digits=4,parse-numbers=false}
\DeclareSIUnit\ohm{\ensuremath{\symup{\Omega}}} % U+2126 absent de la police
\DeclareSIUnit\division{div} % réglages d'oscilloscope (ms/div, V/div)
\DeclareSIUnit\tr{tr} % tours (vitesse de rotation, ex. \SI{1500}{\tr\per\minute})
\DeclareSIUnit\molar{M} % concentration molaire (ex. \SI{3}{\molar}, \SI{1}{\micro\molar})
\DeclareSIUnit\an{an} % année (le modèle confond parfois avec \year, un compteur TeX) — ex. \SI{5}{\kilo\meter\per\an}
\DeclareSIUnit\FCFA{FCFA} % franc CFA (ex. \SI{500}{\FCFA\per\kilo\gram})
\DeclareSIUnit\degreeBrix{\textdegree Brix} % teneur en sucre d'un sirop/jus (réfractométrie)
\DeclareSIUnit\month{mois} % le modèle utilise parfois \month, un compteur TeX, comme unité de durée
\usepackage{qrcode}
% \degree : commande fréquemment utilisée par les modèles pour un angle ou une
% température en dehors de \SI{}{} (non fournie par les paquets chargés).
\providecommand{\degree}{^{\circ}}
% \qed (fin de démonstration), utilisé par les modèles sans amsthm
\providecommand{\qed}{\hfill$\square$}
% \barre : double barre des valeurs interdites dans un tableau de variations (tkz-tab)
\providecommand{\barre}{\ensuremath{\|}}
% environnement proof (amsthm non chargé) utilisé par les modèles
\ifdefined\proof\else\newenvironment{proof}[1][Démonstration]{\par\noindent\textit{#1.}\ }{\qed\par}\fi
\usepackage{lastpage}
\usepackage{hyperref}
\hypersetup{hidelinks}

% ============================================================
% Palette « Pop » OnBuch+ — mêmes hex que la classe P (plus_ui.dart)
% ============================================================
\definecolor{poporange}{HTML}{F59321}
\definecolor{popdark}{HTML}{E07A0C}
\definecolor{popcream}{HTML}{FFF6E8}
\definecolor{popink}{HTML}{1C1714}
\definecolor{poppurple}{HTML}{7A5AE0}
\definecolor{popgreen}{HTML}{2E9E6A}
\definecolor{poppink}{HTML}{E0457B}
\definecolor{popblue}{HTML}{2F7DE1}
\definecolor{popgold}{HTML}{C98A12}
\definecolor{popmuted}{HTML}{8A7F76}
\definecolor{popline}{HTML}{EADFD2}
\definecolor{popgray}{HTML}{8A7F76}
\colorlet{popgrey}{popgray}
% Alias tolérants (le modèle nomme parfois les couleurs différemment)
\colorlet{popwhite}{white}
\colorlet{popblack}{popink}
\colorlet{popred}{poppink}
\colorlet{popyellow}{popgold}
% Teintes claires (fonds de tableaux/encadrés) — le modèle nomme parfois
% les couleurs avec un suffixe L (ex. popblueL) sans les avoir définies.
\colorlet{poporangeL}{poporange!20}
\colorlet{popdarkL}{popdark!20}
\colorlet{poppurpleL}{poppurple!20}
\colorlet{popgreenL}{popgreen!20}
\colorlet{poppinkL}{poppink!20}
\colorlet{popblueL}{popblue!20}
\colorlet{popgoldL}{popgold!20}
\colorlet{popredL}{popred!20}
\colorlet{popgrayL}{popgray!20}
% Teintes claires pour fonds de tableaux / zones
\colorlet{popblueL}{popblue!10!white}
\colorlet{poporangeL}{poporange!14!white}
\colorlet{popgreenL}{popgreen!12!white}
\colorlet{poppurpleL}{poppurple!12!white}
\colorlet{poppinkL}{poppink!12!white}
\colorlet{popgoldL}{popgold!14!white}

\pagecolor{popcream}

% ============================================================
% Typographie
% ============================================================
\defaultfontfeatures{Ligatures=TeX}
\setmainfont{PlusJakartaSans-Medium.ttf}[Path=../../../fonts/,BoldFont=PlusJakartaSans-Bold.ttf]
% Maths en sans-serif (Fira Math) avec chiffres et lettres droites de Plus
% Jakarta, pour que les formules se fondent dans le texte.
\usepackage[math-style=ISO,bold-style=ISO]{unicode-math}
\setmathfont{FiraMath-Regular.otf}[Path=../../../fonts/]
\setmathfont{PlusJakartaSans-Medium.ttf}[Path=../../../fonts/,range={up/num,up/{Latin,latin},\mathup}]
\usepackage{icomma} % virgule décimale sans espace en mode math
% Caractères mathématiques Unicode absents des polices de texte (Plus Jakarta,
% Archivo Black) : rendus par leur équivalent mathématique, en texte comme en
% formule — sinon ils s'affichent en carré vide (ex. « Arithmétique dans ℤ »).
\usepackage{newunicodechar}
\newunicodechar{ℝ}{\ensuremath{\mathbb{R}}}
\newunicodechar{ℤ}{\ensuremath{\mathbb{Z}}}
\newunicodechar{ℂ}{\ensuremath{\mathbb{C}}}
\newunicodechar{ℕ}{\ensuremath{\mathbb{N}}}
\newunicodechar{ℚ}{\ensuremath{\mathbb{Q}}}
\newunicodechar{𝔻}{\ensuremath{\mathbb{D}}}
\newunicodechar{α}{\ensuremath{\alpha}}
\newunicodechar{β}{\ensuremath{\beta}}
\newunicodechar{γ}{\ensuremath{\gamma}}
\newunicodechar{δ}{\ensuremath{\delta}}
\newunicodechar{ε}{\ensuremath{\varepsilon}}
\newunicodechar{θ}{\ensuremath{\theta}}
\newunicodechar{λ}{\ensuremath{\lambda}}
\newunicodechar{μ}{\ensuremath{\mu}}
\newunicodechar{σ}{\ensuremath{\sigma}}
\newunicodechar{τ}{\ensuremath{\tau}}
\newunicodechar{φ}{\ensuremath{\varphi}}
\newunicodechar{ψ}{\ensuremath{\psi}}
\newunicodechar{ω}{\ensuremath{\omega}}
\newunicodechar{Σ}{\ensuremath{\Sigma}}
% majuscules produites par \MakeUppercase dans les titres (α-aminés -> Α)
\newunicodechar{Α}{\ensuremath{\alpha}}
\newunicodechar{Β}{\ensuremath{\beta}}
\newunicodechar{Γ}{\ensuremath{\Gamma}}
\newunicodechar{Δ}{\ensuremath{\Delta}}
\newunicodechar{Ε}{\ensuremath{\varepsilon}}
\newunicodechar{Θ}{\ensuremath{\Theta}}
\newunicodechar{Λ}{\ensuremath{\Lambda}}
\newunicodechar{Μ}{\ensuremath{\mu}}
\newunicodechar{Τ}{\ensuremath{\tau}}
\newunicodechar{Φ}{\ensuremath{\Phi}}
\newunicodechar{Ψ}{\ensuremath{\Psi}}
\newunicodechar{Ω}{\ensuremath{\Omega}}
\newunicodechar{↦}{\ensuremath{\mapsto}}
\newunicodechar{∈}{\ensuremath{\in}}
\newunicodechar{∉}{\ensuremath{\notin}}
\newunicodechar{∧}{\ensuremath{\wedge}}
\newunicodechar{⇔}{\ensuremath{\Leftrightarrow}}
\newunicodechar{⟺}{\ensuremath{\Longleftrightarrow}}
\newunicodechar{⇒}{\ensuremath{\Rightarrow}}
\newunicodechar{⟹}{\ensuremath{\Longrightarrow}}
\newunicodechar{∘}{\ensuremath{\circ}}
\newunicodechar{∪}{\ensuremath{\cup}}
\newunicodechar{∩}{\ensuremath{\cap}}
\newunicodechar{≡}{\ensuremath{\equiv}}
\newunicodechar{⊂}{\ensuremath{\subset}}
\newunicodechar{⊥}{\ensuremath{\perp}}
\newunicodechar{∃}{\ensuremath{\exists}}
\newunicodechar{∀}{\ensuremath{\forall}}
\newunicodechar{∛}{\ensuremath{\sqrt[3]{\,}}}
\newunicodechar{∜}{\ensuremath{\sqrt[4]{\,}}}
\newunicodechar{⃗}{\ensuremath{{}^{\to}}}
\newunicodechar{ⁿ}{\ifmmode^{n}\else\textsuperscript{\ensuremath{n}}\fi}
\newunicodechar{ˣ}{\ifmmode^{x}\else\textsuperscript{\ensuremath{x}}\fi}
\newunicodechar{ᵇ}{\ifmmode^{b}\else\textsuperscript{\ensuremath{b}}\fi}
\newunicodechar{ʳ}{\ifmmode^{r}\else\textsuperscript{\ensuremath{r}}\fi}
\newunicodechar{ᵘ}{\ifmmode^{u}\else\textsuperscript{\ensuremath{u}}\fi}
\newunicodechar{ʸ}{\ifmmode^{y}\else\textsuperscript{\ensuremath{y}}\fi}
\newunicodechar{ᵃ}{\ifmmode^{a}\else\textsuperscript{\ensuremath{a}}\fi}
\newunicodechar{ᵉ}{\ifmmode^{e}\else\textsuperscript{\ensuremath{e}}\fi}
\newunicodechar{ᵏ}{\ifmmode^{k}\else\textsuperscript{\ensuremath{k}}\fi}
\newunicodechar{ᵖ}{\ifmmode^{p}\else\textsuperscript{\ensuremath{p}}\fi}
\newunicodechar{ᴺ}{\ifmmode^{N}\else\textsuperscript{\ensuremath{N}}\fi}
\newunicodechar{ᶜ}{\ifmmode^{c}\else\textsuperscript{\ensuremath{c}}\fi}
\newunicodechar{ʰ}{\ifmmode^{h}\else\textsuperscript{\ensuremath{h}}\fi}
\newunicodechar{ᵠ}{\ifmmode^{q}\else\textsuperscript{\ensuremath{q}}\fi}
\newunicodechar{ₙ}{\ifmmode_{n}\else\textsubscript{\ensuremath{n}}\fi}
\newunicodechar{ₐ}{\ifmmode_{a}\else\textsubscript{\ensuremath{a}}\fi}
\newunicodechar{ᵢ}{\ifmmode_{i}\else\textsubscript{\ensuremath{i}}\fi}
\newunicodechar{ᵣ}{\ifmmode_{r}\else\textsubscript{\ensuremath{r}}\fi}
\newunicodechar{ₖ}{\ifmmode_{k}\else\textsubscript{\ensuremath{k}}\fi}
\newunicodechar{ᵦ}{\ifmmode_{\beta}\else\textsubscript{\ensuremath{\beta}}\fi}
\newunicodechar{ₓ}{\ifmmode_{x}\else\textsubscript{\ensuremath{x}}\fi}
\newfontfamily\popdisplay{ArchivoBlack-Regular.ttf}[Path=../../../fonts/]
\newfontfamily\poplogo{SpaceGrotesk-Bold.ttf}[Path=../../../fonts/]
\newfontfamily\popsemi{PlusJakartaSans-SemiBold.ttf}[Path=../../../fonts/]
\newfontfamily\popxbold{PlusJakartaSans-ExtraBold.ttf}[Path=../../../fonts/]
\newfontfamily\popxboldls{PlusJakartaSans-ExtraBold.ttf}[Path=../../../fonts/,LetterSpace=3.0]

\setlist[enumerate]{leftmargin=1.7em,itemsep=5pt,topsep=4pt}
\setlist[itemize]{leftmargin=1.7em,itemsep=4pt,topsep=4pt,label={\color{poporange}\textbullet}}
\setlength{\parindent}{0pt}
\setlength{\parskip}{5pt}
\renewcommand{\arraystretch}{1.25}

% pgfplots : style Pop par défaut
\pgfplotsset{
  every axis/.append style={axis line style={popink,line width=1pt},tick style={popink},
    grid style={popline,line width=0.5pt},label style={font=\small},tick label style={font=\footnotesize}},
  popcurve/.style={poporange,line width=1.8pt,no markers,smooth},
}
% Même style disponible dans un \draw TikZ simple (hors pgfplots)
\tikzset{popcurve/.style={poporange,line width=1.8pt}}
\tikzset{popgrid/.style={popline,very thin}}
\colorlet{popcurve}{poporange} % le nom du style est parfois utilisé comme couleur

% ============================================================
% Pill / squiggle
% ============================================================
\newcommand{\poppill}[3]{%
  \tikz[baseline=-0.55ex]\node[draw=popink,line width=1.2pt,rounded corners=2.6pt,%
    fill=#2,inner xsep=6.5pt,inner ysep=3pt]{%
    \popxboldls\fontsize{7}{7}\selectfont\color{#3}\MakeUppercase{#1}};%
}
\newcommand{\popsquiggle}[1]{%
  \tikz[baseline=-0.4ex]{
    \draw[#1,line width=2.6pt,line cap=round]
      plot [smooth] coordinates {(0,0.09) (0.34,0.01) (0.68,0.21) (1.02,0.01) (1.36,0.10)};
  }%
}
% Mot-clé surligné (marqueur orange sous le texte)
\newcommand{\cle}[1]{{\popxbold\color{popdark}#1}}

% ============================================================
% En-tête / pied
% ============================================================
\pagestyle{fancy}
\fancyhf{}
\renewcommand{\headrulewidth}{0pt}
\renewcommand{\footrulewidth}{0pt}
\lhead{\raisebox{-0.15ex}{\poplogo\fontsize{13}{13}\selectfont\color{popink}OnBuch\color{poporange}+}}
\rhead{\poppill{\DOCMATIERE\ \textbullet\ \DOCNIVEAU\ \textbullet\ Leçon \DOCLECON}{white}{popink}}
\lfoot{\popxbold\fontsize{7.6}{7.6}\selectfont\color{popmuted}\MakeUppercase{Cours signé OnBuch+}\ \textbullet\ {\popsemi\DOCTITRE}}
\rfoot{\tikz[baseline=-0.6ex]\node[rounded corners=8.5pt,fill=popink,minimum height=17pt,minimum width=17pt,inner xsep=5pt,inner sep=0]{\popxbold\fontsize{8}{8}\selectfont\color{popcream}\thepage\,/\,\pageref*{LastPage}};}

% ============================================================
% Titres
% ============================================================
\newcommand{\chaptitre}[2]{%
  \begin{tcolorbox}[enhanced,colback=poporange,colframe=popink,boxrule=2.6pt,arc=11pt,
      drop shadow=popink,left=15pt,right=15pt,top=12pt,bottom=11pt,before skip=3pt,after skip=18pt]
    \poppill{#2}{popink}{popcream}\par\vspace{9pt}
    {\raggedright\hyphenpenalty=10000\exhyphenpenalty=10000\popdisplay\fontsize{22}{24}\selectfont\color{popcream}\MakeUppercase{#1}\par}
  \end{tcolorbox}
}
% Section numérotée : pastille orange avec le numéro + titre Archivo
\newcounter{popsec}
\newcommand{\coursec}[1]{%
  \stepcounter{popsec}\setcounter{popsub}{0}%
  \par\vspace{16pt}\noindent
  \begin{minipage}{\linewidth}
  \tikz[baseline=-0.7ex]\node[draw=popink,line width=1.6pt,fill=poporange,rounded corners=4pt,minimum size=22pt,inner sep=0]{\popdisplay\fontsize{12}{12}\selectfont\color{popcream}\thepopsec};%
  \hspace{8pt}{\popdisplay\fontsize{15}{17}\selectfont\color{popink}\MakeUppercase{#1}}\par
  \vspace{4pt}\noindent{\color{poporange}\rule{36pt}{3.6pt}}
  \end{minipage}\par\nopagebreak\vspace{8pt}
}
\newcounter{popsub}
\newcommand{\courssub}[1]{%
  \stepcounter{popsub}%
  \par\vspace{9pt}\noindent{\popxbold\fontsize{11.5}{13}\selectfont\color{popink}{\color{poporange}\thepopsec.\thepopsub}\hspace{6pt}#1}\par\nopagebreak\vspace{4pt}
}

% ============================================================
% Boîtes pédagogiques — même recette : fond blanc, bordure épaisse colorée,
% ombre dure encre, badge plein. Titre optionnel [#1].
% ============================================================
\tcbset{popbase/.style={breakable,enhanced,colback=white,boxrule=2.3pt,arc=9pt,
  shadow={3pt}{-3pt}{0pt}{popink},left=14pt,right=14pt,top=11pt,bottom=11pt,before skip=11pt,after skip=15pt,
  fontupper=\popsemi\fontsize{10.6}{14.5}\selectfont\color{popink},
  fontlower=\popsemi\fontsize{10.6}{14.5}\selectfont\color{popink}}}
% Coche dessinée (le glyphe ✓ n'existe pas dans les polices du document)
\newcommand{\popcheck}[1]{\tikz[baseline=-0.5ex]\draw[#1,line width=1.6pt,line cap=round,line join=round](-0.13,0.0)--(-0.04,-0.09)--(0.13,0.1);}
\providecommand{\coche}{\popcheck{popgreen}}
\providecommand{\resultat}[1]{\par\smallskip\noindent\fcolorbox{popgreen}{popgreenL}{\parbox{\dimexpr\linewidth-2\fboxsep-2\fboxrule\relax}{\textbf{Résultat.} #1}}\par\smallskip}
\newcommand{\popboxhead}[3]{\poppill{#1}{#2}{white}\ifx\relax#3\relax\else\hspace{6pt}{\popxbold\fontsize{10.6}{12}\selectfont\color{popink}#3}\fi\par\vspace{7pt}}

\newtcolorbox{aretenir}[1][]{popbase,colframe=popgreen,before upper={\popboxhead{À retenir}{popgreen}{#1}}}
\newtcolorbox{exemplebox}[1][]{popbase,colframe=poporange,before upper={\popboxhead{Exemple}{poporange}{#1}}}
\newtcolorbox{definition}[1][]{popbase,colframe=popblue,before upper={\popboxhead{Définition}{popblue}{#1}}}
\newtcolorbox{propriete}[1][]{popbase,colframe=poppurple,before upper={\popboxhead{Propriété}{poppurple}{#1}}}
\newtcolorbox{methode}[1][]{popbase,colframe=popgold,before upper={\popboxhead{Méthode}{popgold}{#1}}}
\newtcolorbox{attention}[1][]{popbase,colframe=poppink,before upper={\popboxhead{Attention}{poppink}{#1}}}
\newtcolorbox{experience}[1][]{popbase,colframe=popblue,colback=popblueL,before upper={\popboxhead{Expérience}{popblue}{#1}}}
% « Le savais-tu ? » : carte violette pleine (contexte, culture, Cameroun)
\newtcolorbox{savaistu}[1][]{popbase,colback=poppurple,colframe=popink,
  fontupper=\popsemi\fontsize{10.4}{14.2}\selectfont\color{white},
  before upper={\poppill{Le savais-tu\,?}{white}{poppurple}\ifx\relax#1\relax\else\hspace{6pt}{\popxbold\color{white}#1}\fi\par\vspace{7pt}}}
% Exercice résolu : énoncé en haut, \tcblower puis la solution sur fond crème
\newcounter{popexo}\newcounter{popexr}
\newtcolorbox{exoresolu}[1][]{popbase,colframe=poporange,colbacklower=popcream,
  segmentation style={popink,line width=1.2pt,dashed},
  before upper={\stepcounter{popexr}\popboxhead{Exercice résolu \thepopexr}{poporange}{#1}},
  before lower={\poppill{Solution}{popink}{popcream}\par\vspace{6pt}}}
% Exercice (non résolu en place) — la correction est dans la partie Corrigés
\newtcolorbox{exercice}[1][]{popbase,colframe=popink,
  before upper={\stepcounter{popexo}\popboxhead{Exercice \thepopexo}{popink}{#1}}}
% Corrigé
\newtcolorbox{corrige}[1][]{popbase,colframe=popgreen,colback=popgreenL,
  before upper={\popboxhead{Corrigé}{popgreen}{#1}}}
% Objectifs / prérequis / bilan
\newtcolorbox{objectifs}{popbase,colframe=popink,colback=poporangeL,
  before upper={\popboxhead{Objectifs de la leçon}{popink}{}}}
\newtcolorbox{prerequis}{popbase,colframe=popink,colback=popblueL,
  before upper={\popboxhead{Prérequis}{popblue}{}}}
\newtcolorbox{bilan}{popbase,colframe=popink,colback=white,boxrule=2.8pt,
  before upper={\popboxhead{Fiche bilan}{poporange}{L'essentiel à connaître pour l'examen}}}
% Figure encadrée avec légende
\newtcolorbox{popfigure}[1][]{enhanced,breakable=false,colback=white,colframe=popline,boxrule=1.4pt,arc=8pt,
  left=10pt,right=10pt,top=10pt,bottom=8pt,before skip=10pt,after skip=12pt,halign=center,
  lower separated=false,halign lower=center,
  fontlower=\popsemi\fontsize{9}{11.5}\selectfont\color{popmuted},
  #1}
\newcommand{\legende}[1]{\par\vspace{6pt}{\popsemi\fontsize{9}{11.5}\selectfont\color{popmuted}#1}}

% ============================================================
% Signature OnBuch+
% ============================================================
\newcommand{\poplogoinline}[1]{{\poplogo\fontsize{#1}{#1}\selectfont\color{popink}OnBuch\color{poporange}+}}
\newcommand{\signatureonbuch}{%
  \begin{tcolorbox}[enhanced,colback=popink,colframe=popink,boxrule=2.2pt,arc=10pt,
      drop shadow=poporange,left=14pt,right=14pt,top=10pt,bottom=10pt,before skip=0pt,after skip=0pt]
    \tikz[baseline=-0.7ex]\node[circle,fill=popgreen,draw=popcream,line width=1.3pt,minimum size=22pt,inner sep=0]{\popcheck{white}};%
    \hspace{9pt}\begin{minipage}[c]{0.62\linewidth}
      {\popxbold\fontsize{12}{13}\selectfont\color{popcream}Signé\ \textbullet\ {\poplogo OnBuch\color{poporange}+}}\par\vspace{2pt}
      {\raggedright\popsemi\fontsize{8}{10.5}\selectfont\color{popline}\MakeUppercase{Équipe pédagogique OnBuch+}\\\MakeUppercase{Conforme au programme officiel MINESEC}\par}
    \end{minipage}\hfill
    {\poplogo\fontsize{22}{22}\selectfont\color{popcream}OnBuch\color{poporange}+}
  \end{tcolorbox}%
}

% ============================================================
% Page de garde
% #1 titre · #2 matière · #3 niveau · #4 module · #5 n° leçon · #6 durée ·
% #7 description / objectifs (texte ou itemize)
% ============================================================
\newcommand{\pagegarde}[7]{%
  \thispagestyle{empty}
  \newgeometry{a4paper,margin=1.7cm,top=1.6cm,bottom=1.4cm}%
  \begin{tikzpicture}[remember picture,overlay]
    \fill[poporange!18] ([xshift=-3.2cm,yshift=-3.6cm]current page.north east) circle (4.6cm);
    \fill[poppurple!14] ([xshift=2.4cm,yshift=7.6cm]current page.south west) circle (3.8cm);
  \end{tikzpicture}%
  {\poplogo\fontsize{30}{30}\selectfont\color{popink}OnBuch\color{poporange}+}\hfill
  \raisebox{0.6ex}{\poppill{Cours signé \textbullet\ Premium}{popink}{popcream}}\par
  \vspace{1pt}\popsquiggle{poppurple}\par
  \vspace{8pt}
  \begin{tcolorbox}[enhanced,colback=poporange,colframe=popink,boxrule=2.8pt,arc=13pt,
      drop shadow=popink,left=18pt,right=18pt,top=12pt,bottom=13pt,after skip=10pt]
    \poppill{#2 \textbullet\ #3}{popink}{popcream}\hspace{5pt}\poppill{#4}{white}{popink}\par\vspace{8pt}
    {\popdisplay\fontsize{14}{14}\selectfont\color{popink}LEÇON #5}\par\vspace{4pt}
    {\raggedright\hyphenpenalty=10000\exhyphenpenalty=10000\popdisplay\fontsize{25}{27}\selectfont\color{popcream}\MakeUppercase{#1}\par}
    \vspace{8pt}
    {\popxbold\fontsize{9.5}{9.5}\selectfont\color{popink}\MakeUppercase{Durée conseillée : #6}}
  \end{tcolorbox}
  \begin{tcolorbox}[enhanced,colback=white,colframe=popink,boxrule=2.2pt,arc=10pt,
      drop shadow=popink,left=16pt,right=16pt,top=10pt,bottom=10pt,after skip=8pt]
    \poppill{Dans cette leçon}{poppurple}{white}\par\vspace{5pt}
    \popsemi\fontsize{9.3}{12.6}\selectfont\color{popink}#7
  \end{tcolorbox}
  \noindent\begin{minipage}[t]{0.487\linewidth}
    \begin{tcolorbox}[enhanced,colback=popgreen,colframe=popink,boxrule=2.2pt,arc=10pt,
        drop shadow=popink,left=11pt,right=11pt,top=10pt,bottom=10pt,height=3.1cm,valign=center]
      \tcbox[colback=white,colframe=popink,boxrule=1.3pt,arc=5pt,boxsep=0pt,nobeforeafter,
        left=3pt,right=3pt,top=3pt,bottom=3pt,valign=center]{\qrcode[height=1.9cm]{https://play.google.com/store/apps/details?id=com.luvvix.onbuch&hl=fr}}%
      \hspace{8pt}\begin{minipage}[c]{0.5\linewidth}
        {\popdisplay\fontsize{11}{12.5}\selectfont\color{white}\MakeUppercase{Télécharge OnBuch}}\par\vspace{4pt}
        {\raggedright\popsemi\fontsize{8.4}{11}\selectfont\color{white}Gratuit \textbullet\ Google Play\\Tuteur IA, annales, concours, résultats\par}
      \end{minipage}
    \end{tcolorbox}
  \end{minipage}\hfill
  \begin{minipage}[t]{0.487\linewidth}
    \begin{tcolorbox}[enhanced,colback=poppurple,colframe=popink,boxrule=2.2pt,arc=10pt,
        drop shadow=popink,left=11pt,right=11pt,top=10pt,bottom=10pt,height=3.1cm,valign=center]
      \tikz[baseline=-0.7ex]\node[circle,fill=white,draw=popink,line width=1.3pt,minimum size=1.6cm,inner sep=0]{\popdisplay\fontsize{24}{24}\selectfont\color{poppurple}+};%
      \hspace{8pt}\begin{minipage}[c]{0.6\linewidth}
        {\popdisplay\fontsize{11}{12.5}\selectfont\color{white}\MakeUppercase{Passe à OnBuch+}}\par\vspace{4pt}
        {\raggedright\popsemi\fontsize{8.4}{11}\selectfont\color{white}Tous les cours signés, Tuteur IA illimité, fiches PDF — onglet OnBuch+ de l'app\par}
      \end{minipage}
    \end{tcolorbox}
  \end{minipage}
  \vspace{8pt}
  \popcontenu
  \vfill
  \signatureonbuch
  \restoregeometry
  \newpage
}

% Rangée « Ce cours contient » (page de garde)
\newcommand{\poptile}[3]{% #1 couleur #2 gros texte #3 légende
  \begin{tcolorbox}[enhanced,colback=#1,colframe=popink,boxrule=2pt,arc=9pt,shadow={2.5pt}{-2.5pt}{0pt}{popink},
      left=6pt,right=6pt,top=9pt,bottom=9pt,halign=center,valign=center,height=1.75cm,width=\linewidth,nobeforeafter]
    {\popdisplay\fontsize{12.5}{13}\selectfont\color{white}\MakeUppercase{#2}}\par\vspace{3pt}
    {\centering\popsemi\fontsize{7.8}{9.5}\selectfont\color{white}#3\par}
  \end{tcolorbox}}
\newcommand{\popcontenu}{%
  \par\noindent{\popxboldls\fontsize{7.5}{7.5}\selectfont\color{popmuted}CE COURS SIGNÉ CONTIENT}\par\vspace{7pt}
  \noindent\begin{minipage}[t]{0.235\linewidth}\poptile{popblue}{Cours}{complet, illustré, pas à pas}\end{minipage}\hfill
  \begin{minipage}[t]{0.235\linewidth}\poptile{poporange}{Méthodes}{et exercices résolus}\end{minipage}\hfill
  \begin{minipage}[t]{0.235\linewidth}\poptile{poppink}{Exercices}{type examen + corrigés}\end{minipage}\hfill
  \begin{minipage}[t]{0.235\linewidth}\poptile{popgreen}{Fiche bilan}{l'essentiel à retenir}\end{minipage}\par}

% Sommaire
\newtcolorbox{sommaire}{enhanced,colback=white,colframe=popink,boxrule=2.2pt,arc=10pt,
  drop shadow=popink,left=18pt,right=18pt,top=13pt,bottom=13pt,before skip=6pt,after skip=20pt,
  before upper={\poppill{Sommaire}{poppurple}{white}\par\vspace{9pt}\popsemi\fontsize{10.6}{14.5}\selectfont\color{popink}}}

% Signature de fin de cours — poussée tout en bas de la dernière page
\newcommand{\findecours}{%
  \par\vfill\nopagebreak
  \begin{center}\popsquiggle{poporange}\end{center}\vspace{4pt}
  \signatureonbuch
}
\AtBeginDocument{\def\llbracket{\mathopen{[\mkern-3.5mu[}}\def\rrbracket{\mathclose{]\mkern-3.5mu]}}} % intervalles d'entiers (glyphe ⟦ absent de la police)
% Glyphes absents de Fira Math : redéfinis à partir de glyphes disponibles
\AtBeginDocument{%
  \def\setminus{\mathbin{\backslash}}\let\smallsetminus\setminus
  \def\vdots{\mathord{\vbox{\baselineskip3.2pt\lineskiplimit0pt\kern4pt\hbox{.}\hbox{.}\hbox{.}}}}%
  \def\ddots{\mathinner{\mkern1mu\raise6.5pt\hbox{.}\mkern2mu\raise3.6pt\hbox{.}\mkern2mu\raise0.7pt\hbox{.}\mkern1mu}}%
  \def\triangle{\mathord{\scalebox{0.9}{$\symup{\Delta}$}}}%
  \def\nabla{\mathord{\raisebox{\depth}{\scalebox{1}[-1]{$\symup{\Delta}$}}}}%
  \def\checkmark{\popcheck{popdark}}%
  \def\star{\mathbin{\ast}}\def\bigstar{\mathord{\ast}}%
  \def\diamond{\mathbin{\scalebox{0.6}{$\Diamond$}}}%
  \def\bigcup{\mathop{\mathchoice{\raisebox{-0.3em}{\scalebox{1.7}{$\cup$}}}{\scalebox{1.25}{$\cup$}}{\cup}{\cup}}\displaylimits}%
  \def\bigcap{\mathop{\mathchoice{\raisebox{-0.3em}{\scalebox{1.7}{$\cap$}}}{\scalebox{1.25}{$\cap$}}{\cap}{\cap}}\displaylimits}%
  \def\lesseqqgtr{\mathrel{\lessgtr}}\def\bigoplus{\mathop{\oplus}\displaylimits}%
}
\providecommand{\ssi}{si et seulement si} % abréviation fréquente
\AtBeginDocument{\providecommand{\wideparen}{\overparen}} % arc de cercle

%% FIGFIX-BEGIN v5 — correctifs globaux des figures (docs/onbuch-generation/tools/figfix)
\usepackage{adjustbox}\usepackage{etoolbox}\tcbuselibrary{raster}
% toute figure TikZ/pgfplots plus large que la ligne est réduite à la largeur disponible
\BeforeBeginEnvironment{tikzpicture}{\begin{adjustbox}{max width=\linewidth}}
\AfterEndEnvironment{tikzpicture}{\end{adjustbox}}
% légendes pgfplots lisibles par-dessus les courbes
\pgfplotsset{every axis/.append style={legend style={fill=white,fill opacity=0.92,text opacity=1,draw=black!15,inner sep=3pt}}}
% étiquettes d'arêtes (syntaxe edge["..."]) avec fond blanc
\tikzset{every edge quotes/.append style={fill=white,inner sep=1.2pt}}
%% FIGFIX-END
\begin{document}

\pagegarde{Amélioration de la production animale par l'exploitation des insectes}{SVT}{Tle D}{Module IV : La Biotechnologie — Entomophagie}{N18}{4 h}{Cette leçon explore l'entomophagie et l'entomoculture comme réponses biotechnologiques à l'insuffisance des protéines animales classiques. Elle couvre la valeur nutritionnelle et thérapeutique des insectes locaux, l'exploitation de leurs produits (miel, venins), leur rôle en lutte biologique et les techniques de production pour l'alimentation humaine et animale.
\vspace{4pt}
\begin{multicols}{2}\small
\begin{itemize}
\item À la fin de cette leçon, je sais identifier les principaux insectes comestibles et thérapeutiques de l'environnement camerounais (termites, vers blancs, abeilles, fourmis).
\item Je sais comparer la composition nutritionnelle (protéines, lipides, minéraux, vitamines) d'insectes comestibles à celle des viandes conventionnelles à partir de tableaux de composition.
\item Je sais expliquer les vertus thérapeutiques attribuées à certains insectes et produits de la ruche (miel, venin, propolis) en distinguant savoirs traditionnels et données scientifiques validées.
\item Je sais décrire les techniques de collecte, d'élevage (entomoculture) et de transformation des insectes pour la consommation humaine ou l'alimentation animale (volaille, poisson).
\item Je sais analyser le rôle des insectes auxiliaires (coccinelles, libellules) dans la lutte biologique contre les ravageurs des cultures et son impact indirect sur la production animale.
\item Je sais argumenter sur l'apport de l'entomophagie à la sécurité alimentaire et à la couverture des besoins en protéines animales face à la démographie galopante.
\end{itemize}\end{multicols}}

\begin{sommaire}
\begin{multicols}{2}
\begin{enumerate}
\item 1. Contexte : Insuffisance protéique et potentialité de l'entomophagie
\item 2. Valeur nutritionnelle et thérapeutique des insectes comestibles locaux
\item 3. Produits de la ruche et substances bioactives : Miel, Venins, Propolis
\item 4. Lutte biologique : Insectes auxiliaires et protection des cultures fourragères
\item 5. Techniques d'entomoculture : De la collecte à l'élevage industriel pour l'alimentation animale
\item 6. Communication, sensibilisation et acceptabilité sociale
\item Activité d'intégration
\item Méthodes et exercices résolus
\item Exercices
\item Corrigés des exercices
\item Fiche bilan
\end{enumerate}
\end{multicols}
\end{sommaire}

\begin{objectifs}
\begin{itemize}
\item À la fin de cette leçon, je sais identifier les principaux insectes comestibles et thérapeutiques de l'environnement camerounais (termites, vers blancs, abeilles, fourmis).
\item Je sais comparer la composition nutritionnelle (protéines, lipides, minéraux, vitamines) d'insectes comestibles à celle des viandes conventionnelles à partir de tableaux de composition.
\item Je sais expliquer les vertus thérapeutiques attribuées à certains insectes et produits de la ruche (miel, venin, propolis) en distinguant savoirs traditionnels et données scientifiques validées.
\item Je sais décrire les techniques de collecte, d'élevage (entomoculture) et de transformation des insectes pour la consommation humaine ou l'alimentation animale (volaille, poisson).
\item Je sais analyser le rôle des insectes auxiliaires (coccinelles, libellules) dans la lutte biologique contre les ravageurs des cultures et son impact indirect sur la production animale.
\item Je sais argumenter sur l'apport de l'entomophagie à la sécurité alimentaire et à la couverture des besoins en protéines animales face à la démographie galopante.
\item Je sais concevoir un outil de sensibilisation (affiche, dépliant, vidéo courte) valorisant la valeur nutritionnelle et thérapeutique des insectes comestibles auprès d'un public cible.
\end{itemize}
\end{objectifs}

\begin{prerequis}
\begin{itemize}
\item Notions de biochimie : structure et rôle des protéines, acides aminés essentiels, lipides (AG saturés/insaturés), glucides.
\item Besoins nutritionnels humains : apports journaliers recommandés (AJR) en protéines, fer, zinc, vitamines B12.
\item Écologie : relations interspécifiques (prédation, parasitisme), réseaux trophiques, dynamique des populations.
\item Bases de biotechnologie (Module IV) : fermentation, bioconversion, valorisation de la biomasse.
\item Connaissance du milieu local : noms vernaculaires des insectes consommés (mingwel, cocos, ngon, etc.) et contextes de consommation.
\end{itemize}
\end{prerequis}

\newpage

\coursec{Contexte : Insuffisance protéique et potentialité de l'entomophagie}

\begin{experience}[Situation d'accroche : le marché Nkololoun]
Au marché Nkololoun de Yaoundé, le prix du kilogramme de poisson a fortement augmenté en quelques années, rendant l'apport protéique quotidien difficile d'accès pour de nombreuses familles. Pendant ce temps, Madame Ngo vend des sacs de termites (\emph{mingwel}) et de vers de palmier (\emph{cocos}) récoltés localement, très prisés pour leur goût et leur réputation d'aliment « fortifiant » selon la tradition. Un élève de Terminale D se demande : ces insectes peuvent-ils réellement compléter, voire remplacer partiellement, la viande et le poisson sur le plan nutritionnel, et comment transformer cette ressource négligée en une filière d'élevage durable pour nourrir le Cameroun de demain ?
\end{experience}

\textbf{Problématique.} Face à l'insuffisance chronique des ressources animales classiques (viande, poisson, lait), les insectes comestibles constituent-ils une alternative nutritionnelle crédible, économiquement viable et écologiquement durable ? Pour répondre à cette question, nous allons d'abord quantifier le déficit protéique camerounais, définir rigoureusement l'entomophagie et l'entomoculture, inventorier les insectes comestibles de nos régions, puis comparer l'empreinte environnementale des insectes à celle du bétail.

\courssub{Une crise protéique silencieuse : l'analyse de la situation camerounaise}

\textbf{Observation.} Selon les données du MINADER et de la FAO, la population camerounaise dépasse 28 millions d'habitants et croît d'environ \num{2,6} \% par an. Or la production nationale de viande, de poisson et de lait ne suit pas ce rythme : le Cameroun importe chaque année des centaines de milliers de tonnes de poisson (principalement du maquereau congelé) et des quantités croissantes de viande et de lait en poudre pour combler son déficit.

\begin{definition}[Sécurité alimentaire]
La \cle{sécurité alimentaire} désigne la situation dans laquelle tous les êtres humains ont, à tout moment, un accès physique, social et économique à une nourriture suffisante, saine et nutritive leur permettant de satisfaire leurs besoins énergétiques et leurs préférences alimentaires pour mener une vie saine et active (définition FAO, Sommet mondial de l'alimentation, 1996).
\end{definition}

\textbf{Interprétation.} Les besoins d'un adulte sont estimés à environ \SI{0,8}{g} de protéines par kilogramme de masse corporelle et par jour, soit environ \SI{50}{g/j} pour un adulte de \SI{65}{kg}. Les protéines animales sont particulièrement précieuses car elles fournissent les \cle{acides aminés essentiels} (que l'organisme ne sait pas synthétiser) dans des proportions équilibrées, ainsi que le fer héminique (bien mieux absorbé que le fer végétal), le zinc et la vitamine B\textsubscript{12}, quasi absente des végétaux. Lorsque les produits animaux deviennent trop chers, les populations se replient sur les féculents (manioc, maïs, igname) : l'énergie est couverte, mais pas les acides aminés essentiels. C'est la \emph{faim cachée}, responsable du retard de croissance (rabougrissement) qui touche encore environ un enfant camerounais sur trois de moins de cinq ans dans certaines régions.

\begin{aretenir}[Le paradoxe camerounais]
Le Cameroun est « l'Afrique en miniature » : il possède des potentialités agro-pastorales et halieutiques considérables, pourtant il reste importateur net de protéines animales. Le déficit ne vient pas seulement de la quantité produite, mais aussi du \textbf{coût} des protéines classiques, qui exclut les ménages à faible revenu.
\end{aretenir}

\textbf{Conclusion de l'analyse.} Il faut donc chercher des sources de protéines animales \emph{locales, peu coûteuses, rapidement renouvelables et peu gourmandes en intrants}. C'est exactement le profil des insectes comestibles.

\courssub{Entomophagie et entomoculture : des définitions précises}

\begin{definition}[Entomophagie]
L'\cle{entomophagie} (du grec \emph{entomon}, insecte, et \emph{phagein}, manger) désigne la consommation d'insectes par l'être humain, à un ou plusieurs stades de leur développement (œuf, larve, nymphe, adulte), que ce soit à l'état frais, séché, grillé, frit ou transformé (farines, pâtes, huiles).
\end{definition}

\begin{definition}[Entomoculture]
L'\cle{entomoculture} est l'élevage raisonné et contrôlé d'insectes, dans le but de produire de la biomasse comestible (pour l'homme ou le bétail), des produits dérivés (miel, soie, colorants) ou des auxiliaires de lutte biologique. Elle se distingue de la simple \emph{collecte} en milieu naturel par le contrôle du cycle biologique complet : reproduction, alimentation, croissance, récolte.
\end{definition}

\begin{attention}[Ne pas confondre]
Il faut distinguer trois réalités :
\begin{itemize}
\item l'\textbf{entomophagie volontaire} : consommation délibérée d'insectes récoltés ou élevés à cette fin ;
\item l'\textbf{entomophagie accidentelle} : ingestion involontaire de fragments d'insectes présents dans les denrées stockées (larves de charançons dans le maïs, par exemple) — elle n'a aucune visée nutritionnelle et peut signaler une contamination ;
\item la \textbf{consommation d'insectes nuisibles ou toxiques} : toutes les espèces ne sont pas comestibles ! Certains insectes accumulent des toxines végétales (chenilles se nourrissant de plantes cyanogènes), des pesticides (sauterelles capturées sur des champs traités) ou hébergent des parasites. Seules les espèces traditionnellement consommées et correctement préparées (cuisson, séchage, éviscération pour certaines chenilles) sont sûres.
\end{itemize}
\end{attention}

\courssub{Inventaire des insectes comestibles majeurs du Cameroun}

La tradition entomophagique camerounaise est ancienne et régionalement marquée. Le tableau suivant synthétise les cinq groupes majeurs.

\begin{center}
\begin{tabularx}{\linewidth}{|l|X|X|}
\hline
\rowcolor{popblueL} \textbf{Insecte (nom scientifique)} & \textbf{Noms locaux et régions de consommation} & \textbf{Stade consommé et préparation} \\
\hline
Termites (\emph{Macrotermes} spp.) & \emph{Mingwel}, \emph{nguim} ; Centre, Sud, Est, Adamaoua & Adultes ailés (essaims des pluies) ; grillés ou frits \\
\hline
Ver de palmier (\emph{Rhynchophorus phoenicis}) & \emph{Cocos}, \emph{mouton de palme} ; Sud, Littoral, Centre, Ouest & Larve ; crue, grillée ou frite \\
\hline
Chenilles (\emph{Imbrasia} spp.) & \emph{Mbinzo}, chenilles comestibles ; Est, Sud (zone forestière) & Chenille ; bouillie puis séchée ou fumée \\
\hline
Sauterelles et grillons (\emph{Ruspolia} spp., \emph{Brachytrupes} spp.) & \emph{Nsenene}, grillons ; Nord, Adamaoua, saisons sèches & Adulte ; grillé après épilation des ailes \\
\hline
Fourmis comestibles (\emph{Carebara vidua}) & Fourmis ailées ; Ouest, Nord-Ouest & Adultes ailés ; grillés \\
\hline
\end{tabularx}
\end{center}

\begin{popfigure}
\begin{tikzpicture}[scale=0.9, every node/.style={font=\small}]
% Contour schématique du Cameroun (forme simplifiée)
\draw[thick, popdark, fill=popcream] (0,0) -- (1.2,-0.3) -- (2.2,0.4) -- (3.4,0.1) -- (4.4,0.8) -- (5.0,2.0) -- (4.6,3.2) -- (5.2,4.4) -- (4.4,5.6) -- (3.0,6.2) -- (2.6,7.4) -- (2.0,8.6) -- (1.4,7.6) -- (0.8,6.4) -- (-0.2,5.2) -- (-0.6,3.8) -- (-1.0,2.6) -- (-0.4,1.2) -- cycle;
% Régions approximatives et insectes
\node[popblue, font=\small\bfseries] at (1.9,7.6) {Extrême-Nord / Nord};
\node[popmuted] at (1.9,7.1) {Sauterelles, grillons};
\node[popblue, font=\small\bfseries] at (3.6,5.6) {Adamaoua};
\node[popmuted] at (3.6,5.1) {Termites, sauterelles};
\node[popblue, font=\small\bfseries] at (3.9,3.2) {Est};
\node[popmuted] at (3.9,2.7) {Chenilles (\emph{Imbrasia})};
\node[popblue, font=\small\bfseries] at (1.4,3.6) {Centre};
\node[popmuted] at (1.4,3.1) {Termites, vers de palmier};
\node[popblue, font=\small\bfseries] at (-0.2,4.8) {Ouest / N.-Ouest};
\node[popmuted] at (-0.2,4.3) {Fourmis \emph{Carebara}};
\node[popblue, font=\small\bfseries] at (0.6,1.4) {Sud / Littoral};
\node[popmuted] at (0.6,0.9) {Vers de palmier, termites};
% Points de localisation
\fill[poporange] (1.9,6.8) circle (0.08);
\fill[poporange] (3.4,4.9) circle (0.08);
\fill[poporange] (3.7,2.5) circle (0.08);
\fill[poporange] (1.3,2.9) circle (0.08);
\fill[poporange] (0.0,4.1) circle (0.08);
\fill[poporange] (0.5,0.7) circle (0.08);
\end{tikzpicture}
\legende{Carte schématique du Cameroun localisant les zones de consommation préférentielle des cinq groupes d'insectes comestibles majeurs. Les points orange indiquent les zones de forte tradition entomophagique. La répartition suit la végétation : chenilles et vers de palmier en zone forestière humide (Sud, Est), sauterelles et grillons en zone soudano-sahélienne (Nord), termites répartis dans presque toutes les régions.}
\end{popfigure}

\begin{savaistu}[Une biodiversité comestible exceptionnelle]
La FAO a recensé plus de \num{1900} espèces d'insectes comestibles dans le monde. L'Afrique en compte environ 500, et le Cameroun à lui seul près de 150 espèces traditionnellement consommées — l'un des pays les plus riches d'Afrique en diversité entomophagique. Cette ressource est gratuite, renouvelable et déjà culturellement acceptée : c'est un atout que beaucoup de pays nous envient.
\end{savaistu}

\courssub{Cadre réglementaire : de la tradition à la filière encadrée}

Pour transformer la collecte traditionnelle en filière d'élevage, il faut un cadre légal.

\begin{definition}[Nouvel aliment (Novel Food)]
Un \cle{nouvel aliment} (\emph{Novel Food}) est, selon la réglementation européenne (règlement UE 2015/2283), un aliment qui n'a pas été consommé de façon significative dans l'Union européenne avant le 15 mai 1997 et qui relève de catégories définies (aliments contenant des insectes entiers ou leurs parties, par exemple). Sa mise sur le marché exige une évaluation préalable de sécurité par l'EFSA. Depuis 2021, plusieurs insectes (ver de farine, grillon domestique, criquet migrateur) sont ainsi autorisés dans l'UE.
\end{definition}

\textbf{Au niveau international}, le \cle{Codex Alimentarius} (programme conjoint FAO/OMS) fixe des normes d'hygiène applicables aux denrées alimentaires, y compris aux insectes transformés : contrôle microbiologique, traçabilité, étiquetage. La FAO a publié en 2013 un rapport de référence (\emph{Edible insects: future prospects for food and feed security}) qui légitime scientifiquement la filière.

\textbf{Au Cameroun}, l'élevage des espèces animales non conventionnelles (dont les insectes, les escargots, le gibier d'élevage) relève du MINADER et du MINEPIA ; des arrêtés encadrent l'exercice de l'élevage non conventionnel, l'identification des éleveurs et les conditions sanitaires de production. Ce cadre, encore jeune, est indispensable pour sécuriser la commercialisation à grande échelle et l'exportation.

\courssub{Enjeux environnementaux : pourquoi les insectes sont imbattables}

\textbf{Observation fondatrice.} Les insectes sont des animaux \cle{poïkilothermes} (à température corporelle variable, dits « à sang froid ») : ils ne dépensent pas d'énergie à maintenir leur température corporelle constante, contrairement aux vertébrés homéothermes (bovins, porcs, volailles). Une fraction beaucoup plus grande de l'aliment ingéré est donc convertie en biomasse corporelle.

\begin{propriete}[Efficacité de conversion des insectes]
Du fait de leur métabolisme poïkilotherme, les insectes convertissent l'aliment en biomasse avec un taux de conversion alimentaire (TCA) \textbf{2 à 10 fois plus efficace} que celui des vertébrés homéothermes. Conséquences mesurables : pour une même quantité de protéines produites, l'élevage d'insectes requiert beaucoup moins d'aliment, d'eau et de surface, et émet beaucoup moins de gaz à effet de serre (notamment pas de méthane entérique, contrairement aux ruminants). Ce savoir est admis au programme ; il repose sur les bilans énergétiques comparés du métabolisme.
\end{propriete}

\begin{methode}[Calculer et interpréter le Taux de Conversion Alimentaire (TCA)]
\begin{enumerate}
\item \textbf{Mesurer} la masse totale d'aliment sec distribuée pendant la période d'élevage : $M_a$ (en kg).
\item \textbf{Mesurer} le gain de poids vif des animaux sur la même période : $G = M_{\text{finale}} - M_{\text{initiale}}$ (en kg).
\item \textbf{Calculer} : $\mathrm{TCA} = \dfrac{M_a}{G}$.
\item \textbf{Interpréter} : le TCA indique combien de kilogrammes d'aliment sont nécessaires pour produire 1 kg de poids vif. \emph{Plus le TCA est faible, plus la conversion est efficace.}
\end{enumerate}
\end{methode}

\begin{exoresolu}[Comparaison bœuf / grillon]
Énoncé. Un éleveur de bovins distribue \SI{900}{kg} d'aliment sec à un lot dont le poids vif passe de \SI{800}{kg} à \SI{950}{kg}. Une éleveuse de grillons (\emph{Acheta domesticus}) distribue \SI{34}{kg} d'aliment sec à un lot qui gagne \SI{20}{kg} de biomasse. Calculer les deux TCA et conclure.
\tcblower
Solution détaillée.
\begin{align*}
\text{Bœuf : } \mathrm{TCA} &= \frac{900}{950-800} = \frac{900}{150} = 6 \\
\text{Grillon : } \mathrm{TCA} &= \frac{34}{20} = \num{1,7}
\end{align*}
Interprétation : il faut \SI{6}{kg} d'aliment pour produire \SI{1}{kg} de gain de poids vif de bœuf, contre seulement \SI{1,7}{kg} pour \SI{1}{kg} de grillon. Le grillon est environ $6/\num{1,7} \approx \num{3,5}$ fois plus efficace, ce qui illustre l'efficacité de conversion des insectes poïkilothermes.
\end{exoresolu}

\begin{exemplebox}[L'exemple chiffré de référence (ordres de grandeur FAO)]
Pour produire \SI{1}{kg} de protéines :
\begin{itemize}
\item \textbf{Bœuf} : environ \SI{200}{m^2} de surface et plus de \SI{20000}{L} d'eau ;
\item \textbf{Grillon} : environ \SI{15}{m^2} et environ \SI{1500}{L} d'eau.
\end{itemize}
Soit environ 13 fois moins de surface et plus de 10 fois moins d'eau pour le grillon.
\end{exemplebox}

\begin{popfigure}
\begin{center}
\begin{tikzpicture}
\begin{axis}[
    ybar, bar width=0.35cm,
    width=0.95\linewidth, height=6.5cm,
    symbolic x coords={Boeuf, Porc, Poulet, Grillon, Mouche soldat},
    xtick=data, x tick label style={font=\small, rotate=15, anchor=east},
    ylabel={Valeur (unités rapportées à 1 kg de protéines)},
    ymode=log, log basis y=10,
    ymin=1, ymax=100000,
    legend style={at={(0.5,1.02)}, anchor=south, legend columns=3, font=\small},
    ymajorgrids=true, grid style={popline!40},
]
\addplot[fill=popblue!80, draw=popblue] coordinates {(Boeuf,22000) (Porc,6000) (Poulet,3500) (Grillon,1500) (Mouche soldat,1000)};
\addplot[fill=poporange!80, draw=poporange] coordinates {(Boeuf,200) (Porc,60) (Poulet,45) (Grillon,15) (Mouche soldat,10)};
\addplot[fill=popgreen!80, draw=popgreen] coordinates {(Boeuf,100) (Porc,25) (Poulet,15) (Grillon,2) (Mouche soldat,1)};
\legend{Eau (L), Surface (m$^2$), CO$_2$eq (kg)}
\end{axis}
\end{tikzpicture}
\end{center}
\legende{Comparaison de l'empreinte environnementale de cinq sources de protéines animales (échelle logarithmique) : consommation d'eau (L/kg de protéines), surface au sol (m\textsuperscript{2}/kg) et émissions de gaz à effet de serre (kg CO\textsubscript{2} équivalent/kg). Le bœuf est de loin le plus coûteux en ressources ; le grillon (\emph{Acheta domesticus}) et la mouche soldat noire (\emph{Hermetia illucens}) affichent les empreintes les plus faibles. Ordres de grandeur d'après FAO 2013 et van Huis et al.}
\end{popfigure}

\begin{attention}[Pièges fréquents à l'examen]
\begin{itemize}
\item Ne pas écrire que « tous les insectes sont comestibles » : c'est faux et dangereux (toxines, parasites, résidus de pesticides).
\item Ne pas confondre \textbf{TCA faible} et \textbf{mauvaise conversion} : c'est l'inverse, un TCA faible signifie une conversion efficace.
\item Ne pas attribuer aux insectes d'élevage (grillons, mouches) des émissions de méthane entérique : celles-ci caractérisent essentiellement les ruminants ; seuls quelques insectes comme les termites en produisent, en faible quantité.
\item Dans les comparaisons chiffrées, toujours rapporter à la \textbf{même unité} (1 kg de protéines, et non 1 kg de poids vif, car les teneurs en protéines diffèrent).
\item Ne pas oublier que la \textbf{chitine} de la cuticule des insectes est peu digestible : la fraction de protéines réellement assimilable est légèrement inférieure à la teneur brute mesurée.
\end{itemize}
\end{attention}

\begin{aretenir}[Conclusion de la section]
Le Cameroun fait face à un déficit croissant en protéines animales classiques, dont le coût exclut une grande partie de la population. L'\cle{entomophagie}, déjà ancrée dans les traditions régionales (termites, vers de palmier, chenilles, sauterelles, fourmis), offre une ressource riche en protéines de qualité. L'\cle{entomoculture} permettrait de la transformer en filière durable : grâce à leur métabolisme poïkilotherme, les insectes convertissent l'aliment 2 à 10 fois plus efficacement que le bétail, avec des besoins en eau, en surface et des émissions de gaz à effet de serre drastiquement réduits. Reste à encadrer cette filière (Codex Alimentarius, réglementation sur les nouveaux aliments, arrêtés nationaux sur l'élevage non conventionnel) et à vaincre les réticences culturelles — ce que nous étudierons dans les sections suivantes.
\end{aretenir}

\coursec{Valeur nutritionnelle et thérapeutique des insectes comestibles locaux}

Nous avons vu dans la section précédente que le déficit en protéines animales constitue un problème majeur de santé publique au Cameroun : la ration protéique moyenne reste inférieure aux recommandations de la FAO dans de nombreuses zones rurales, tandis que le prix de la viande et du poisson ne cesse d'augmenter. Face à cette insuffisance, une ressource abondante, renouvelable et traditionnellement consommée s'offre à nous : les insectes. Mais au-delà des habitudes culturelles, que dit la science ? Les insectes sont-ils réellement comparables, sur le plan nutritionnel, à la viande de bœuf ou au poisson ? C'est à cette question rigoureuse que nous allons répondre, en exploitant des données expérimentales de composition chimique, puis en distinguant soigneusement les vertus thérapeutiques prouvées des croyances traditionnelles.

\courssub{Composition centésimale des insectes comestibles locaux}

\begin{definition}[Composition centésimale]
La \cle{composition centésimale} d'un aliment est la teneur de ses constituants majeurs (eau, protéines, lipides, glucides, cendres, fibres) exprimée en grammes pour \SI{100}{g} d'aliment. Elle peut être exprimée sur \cle{base matière fraîche} (aliment tel que consommé, avec son eau) ou sur \cle{base matière sèche} (MS, après élimination de l'eau).
\end{definition}

Les analyses menées par l'IRAD (Institut de Recherche Agricole pour le Développement) et les laboratoires de biochimie des universités de Yaoundé~I, Dschang et Ngaoundéré fournissent les valeurs moyennes suivantes (base matière sèche) :

\begin{center}
\begin{tabularx}{\linewidth}{|l|X|X|X|X|X|}
\hline
\rowcolor{popblueL} \textbf{Insecte (MS)} & \textbf{Protéines (g/100 g)} & \textbf{Lipides (g/100 g)} & \textbf{Cendres (g/100 g)} & \textbf{Chitine/fibres (g/100 g)} & \textbf{Glucides (g/100 g)} \\
\hline
Termites (\emph{Macrotermes} spp.) & 35 -- 45 & 30 -- 45 & 3 -- 6 & 5 -- 10 & 2 -- 8 \\
\hline
Vers de palmier (\emph{Rhynchophorus phoenicis}) & 25 -- 35 & 45 -- 60 & 2 -- 4 & 3 -- 6 & 2 -- 5 \\
\hline
Chenilles comestibles & 45 -- 60 & 10 -- 20 & 5 -- 8 & 6 -- 12 & 5 -- 10 \\
\hline
Sauterelles / criquets & 55 -- 70 & 5 -- 15 & 4 -- 8 & 8 -- 15 & 3 -- 8 \\
\hline
\end{tabularx}
\end{center}

\begin{methode}[Comparer des aliments : passer de la matière fraîche à la matière sèche]
Pour comparer rigoureusement deux aliments d'humidités différentes, il faut raisonner à \textbf{poids de matière sèche égal} :
\begin{enumerate}
\item Relever la teneur en eau de chaque aliment (ex. termites frais : $\approx 60\,\%$ d'eau ; bœuf maigre : $\approx 70\,\%$ d'eau).
\item Calculer la matière sèche : $MS = 100 - \text{eau}$ (en g pour \SI{100}{g}).
\item Convertir chaque constituant : $C_{MS} = C_{frais} \times \dfrac{100}{MS}$.
\item Comparer les valeurs $C_{MS}$ entre elles.
\end{enumerate}
\textbf{Exemple :} des termites frais contiennent \SI{14}{g} de protéines pour \SI{100}{g} et $60\,\%$ d'eau, donc $MS = \SI{40}{g}$. Sur base sèche : $14 \times \dfrac{100}{40} = \SI{35}{g}$ de protéines pour \SI{100}{g} MS.
\end{methode}

\begin{attention}[Piège classique de comparaison]
Comparer « \SI{14}{g} de protéines pour les termites » à « \SI{20}{g} pour le bœuf » sur base fraîche n'a pas de sens si les teneurs en eau diffèrent : une partie de l'écart ne reflète que la dilution par l'eau. Toute comparaison honnête se fait sur \textbf{base matière sèche} ou, mieux, en \textbf{densité nutritionnelle par rapport aux apports journaliers recommandés} (AJR).
\end{attention}

\courssub{Qualité protéique : les acides aminés essentiels}

Une protéine n'est pas seulement une quantité : c'est une \textbf{qualité}. Les protéines sont des polymères d'acides aminés ; parmi les 20 acides aminés, neuf sont dits \cle{essentiels} (AAE) car l'organisme humain ne sait pas les synthétiser : histidine, isoleucine, leucine, lysine, méthionine, phénylalanine, thréonine, tryptophane et valine.

\begin{definition}[Score chimique, PDCAAS et DIAAS]
\begin{itemize}
\item Le \cle{score chimique} d'une protéine est le rapport, exprimé en \%, entre la teneur de l'acide aminé le plus déficitaire (\emph{acide aminé limitant}) et la teneur du même acide aminé dans une protéine de référence (œuf entier).
\item Le \cle{PDCAAS} (\emph{Protein Digestibility Corrected Amino Acid Score}) corrige le score chimique par la digestibilité fécale vraie de la protéine.
\item Le \cle{DIAAS} (\emph{Digestible Indispensable Amino Acid Score}), recommandé par la FAO depuis 2013, corrige chaque AAE par sa digestibilité \textbf{iléale} vraie (mesurée en fin d'intestin grêle) : c'est l'indice le plus rigoureux.
\end{itemize}
\end{definition}

\begin{propriete}[Les insectes, source de protéines « complètes » mais à digestibilité limitée]
Les protéines des insectes comestibles contiennent les \textbf{neuf acides aminés essentiels} : ce sont des protéines dites \emph{complètes}, souvent riches en \cle{lysine} et en \cle{leucine} (acides aminés fréquemment déficitaires dans les tubercules et céréales de base de l'alimentation camerounaise : manioc, maïs). En revanche :
\begin{itemize}
\item l'acide aminé limitant est souvent la \textbf{méthionine + cystine} (acides aminés soufrés) ou le \textbf{tryptophane} selon l'espèce ;
\item la \textbf{digestibilité vraie iléale} est réduite par la \cle{chitine} de l'exosquelette, sauf traitement technologique (dégraissage, mouture fine, hydrolyse enzymatique).
\end{itemize}
\end{propriete}

\begin{definition}[Chitine]
La \cle{chitine} est un polysaccharide structural de l'exosquelette des arthropodes, formé d'unités N-acétylglucosamine liées en $\beta$-1,4. C'est une \textbf{fibre insoluble} que les enzymes digestives humaines hydrolysent très mal (la chitinase humaine est peu active) : elle réduit l'accessibilité des protéines aux protéases, mais joue aussi un rôle de fibre alimentaire.
\end{definition}

Pour l'alimentation animale (volailles, poissons d'élevage), la farine de larves de Mouche Soldat Noire (\emph{Hermetia illucens}) est comparée aux sources protéiques classiques :

\begin{popfigure}
\begin{center}
\begin{tikzpicture}
\begin{axis}[
  ybar, bar width=7pt,
  width=\linewidth, height=7cm,
  symbolic x coords={Lysine, Méthionine, Thréonine, Leucine, Valine},
  xtick=data,
  ylabel={Teneur (mg/g de protéine)},
  ymin=0, ymax=90,
  legend style={at={(0.5,-0.18)},anchor=north,legend columns=3,font=\small},
  ymajorgrids=true, grid style={popline!40},
]
\addplot+[ybar,fill=popgreen,draw=popdark] coordinates {(Lysine,60) (Méthionine,20) (Thréonine,40) (Leucine,75) (Valine,55)};
\addplot+[ybar,fill=popblue,draw=popdark] coordinates {(Lysine,75) (Méthionine,28) (Thréonine,42) (Leucine,72) (Valine,50)};
\addplot+[ybar,fill=poporange,draw=popdark] coordinates {(Lysine,62) (Méthionine,13) (Thréonine,38) (Leucine,78) (Valine,48)};
\legend{Larve de Mouche Soldat Noire, Farine de poisson, Tourteau de soja}
\end{axis}
\end{tikzpicture}
\end{center}
\legende{Profils comparés en cinq acides aminés essentiels (mg/g de protéine) de trois sources protéiques utilisées en alimentation animale. La larve de Mouche Soldat Noire rivalise avec la farine de poisson et le soja, avec une limite en méthionine commune aux trois sources (valeurs moyennes d'après données IRAD et littérature FAO).}
\end{popfigure}

\courssub{Lipides : quantité variable, qualité remarquable}

La teneur en lipides des insectes est extrêmement variable selon l'espèce et le stade : de $2\,\%$ de la matière sèche chez certains criquets à plus de $50\,\%$ chez les vers de palmier. Mais c'est le \textbf{profil en acides gras} qui retient l'attention des nutritionnistes :

\begin{itemize}
\item les orthoptères (sauterelles, criquets) et de nombreuses larves sont riches en \cle{acides gras polyinsaturés} (AGPI), notamment oméga-6 (acide linoléique) et oméga-3 (acide $\alpha$-linolénique) ;
\item les insectes \textbf{ne contiennent pas de cholestérol} en quantité notable : leurs stérols membranaires sont des \cle{phytostérols} et de l'ergostérol, ce qui constitue un \textbf{avantage cardiovasculaire} par rapport aux viandes grasses ;
\item les vers de palmier, très lipidiques, apportent surtout des acides gras saturés (acide palmitique) : leur consommation doit donc être modérée, comme celle de l'huile de palme.
\end{itemize}

\courssub{Micronutriments : fer, zinc, vitamines du groupe B}

Les insectes se distinguent par une densité remarquable en micronutriments :

\begin{itemize}
\item \textbf{Fer} : les termites en sont exceptionnellement riches (jusqu'à \SI{35}{mg} pour \SI{100}{g} frais chez \emph{Macrotermes bellicosus}). Le fer alimentaire existe sous deux formes : le fer \cle{hémique} (viande, poisson, très biodisponible : $15$ à $35\,\%$ absorbé) et le fer \cle{non hémique} (végétaux, insectes : $2$ à $10\,\%$ absorbé, mais favorisé par la vitamine C du repas).
\item \textbf{Zinc, cuivre, magnésium} : présents en quantités significatives, le zinc des termites couvrant environ la moitié des AJR pour \SI{100}{g}.
\item \textbf{Vitamines B2 (riboflavine), acide folique et B12} : la vitamine B12, absente des végétaux, est présente chez plusieurs insectes, mais sa teneur est \textbf{très variable} car elle dépend du microbiote intestinal de l'insecte et de son substrat d'élevage.
\end{itemize}

\begin{definition}[Biodisponibilité]
La \cle{biodisponibilité} d'un nutriment est la fraction de celui-ci qui est réellement absorbée par l'intestin puis utilisable par l'organisme. Elle dépend de la forme chimique du nutriment, de la matrice alimentaire (fibres, phytates, tanins) et de l'état physiologique du consommateur.
\end{definition}

\begin{attention}[Erreur fréquente : « les insectes contiennent plus de fer que le bœuf »]
Cette affirmation, souvent reprise dans les médias, est trompeuse si elle ne précise pas la \textbf{forme chimique} du fer ni sa \textbf{biodisponibilité réelle}. Le fer des insectes est majoritairement non hémique, et la chitine peut piéger une partie des minéraux dans le tube digestif. Une teneur élevée dans le tableau de composition ne garantit pas une absorption élevée chez l'homme : seules des études d'absorption (marquage isotopique) permettent de conclure. Rédige toujours : « teneur élevée en fer, de biodisponibilité à préciser ».
\end{attention}

\begin{exemplebox}[Portrait nutritionnel chiffré des termites]
\SI{100}{g} de termites frais (\emph{Macrotermes bellicosus}) apportent environ : \SI{14}{g} de protéines, \SI{35}{mg} de fer (largement supérieur aux AJR d'un homme adulte, $\approx \SI{8}{mg}$), \SI{5}{mg} de zinc ($\approx 50\,\%$ des AJR) et jusqu'à \SI{20}{\micro g} de vitamine B12 (valeur très variable selon les individus et la saison). C'est un aliment particulièrement précieux pour les femmes en âge de procréer et les enfants, populations les plus exposées à l'anémie ferriprive au Cameroun.
\end{exemplebox}

\begin{popfigure}
\begin{center}
\begin{tikzpicture}[scale=1.05]
% Grille radar
\foreach \r in {1,2,3,4} {
  \draw[popline!50] (90:\r) \foreach \a in {30,-30,-90,-150,150} { -- (\a:\r) } -- cycle;
}
\foreach \a/\t in {90/{Protéines},30/{Fer},-30/{Zinc},-90/{Vit. B12},-150/{AGPI},150/{Calcium}} {
  \draw[popline!70] (0,0) -- (\a:4);
  \node[font=\small\bfseries, popdark] at (\a:4.75) {\t};
}
% Termites (vert)
\draw[very thick, popgreen, fill=popgreen!20] (90:2.8) -- (30:4) -- (-30:2) -- (-90:3.2) -- (-150:2.4) -- (150:0.8) -- cycle;
% Vers de palmier (orange)
\draw[very thick, poporange, fill=poporange!15] (90:1.6) -- (30:2.4) -- (-30:1.6) -- (-90:2) -- (-150:2.8) -- (150:0.6) -- cycle;
% Bœuf (bleu)
\draw[very thick, popblue, fill=popblue!12] (90:2.4) -- (30:1) -- (-30:2.4) -- (-90:3.6) -- (-150:0.4) -- (150:0.2) -- cycle;
% Poisson (violet)
\draw[very thick, poppurple, fill=poppurple!10] (90:2.2) -- (30:0.6) -- (-30:1.2) -- (-90:4) -- (-150:1.6) -- (150:1.2) -- cycle;
% Légende placée sous le graphique
\begin{scope}[shift={(0,-5.5)}]
\draw[very thick, popgreen, fill=popgreen!20] (0,0.3) rectangle (0.4,0.6); \node[right, font=\small] at (0.5,0.45) {Termites (soldats)};
\draw[very thick, poporange, fill=poporange!15] (0,-0.3) rectangle (0.4,0); \node[right, font=\small] at (0.5,-0.15) {Vers de palmier};
\draw[very thick, popblue, fill=popblue!12] (0,-0.9) rectangle (0.4,-0.6); \node[right, font=\small] at (0.5,-0.75) {Bœuf maigre};
\draw[very thick, poppurple, fill=poppurple!10] (0,-1.5) rectangle (0.4,-1.2); \node[right, font=\small] at (0.5,-1.35) {Capitaine (poisson)};
\end{scope}
\end{tikzpicture}
\end{center}
\legende{Graphique radar : couverture des apports journaliers recommandés (AJR) par \SI{100}{g} d'aliment frais, pour six nutriments clés. Chaque cercle concentrique représente $25\,\%$ des AJR (le cercle externe $= 100\,\%$). Les termites dominent pour le fer, le poisson pour la vitamine B12, le bœuf pour le zinc : les sources sont \textbf{complémentaires}, non concurrentes (valeurs moyennes indicatives, sources IRAD/FAO).}
\end{popfigure}

\courssub{Vertus thérapeutiques : entre tradition et preuve scientifique}

La médecine traditionnelle camerounaise attribue aux insectes de nombreuses vertus. La démarche scientifique exige de classer ces affirmations selon leur \textbf{niveau de preuve}.

\begin{center}
\begin{tabularx}{\linewidth}{|l|X|X|}
\hline
\rowcolor{popblueL} \textbf{Usage traditionnel} & \textbf{Justification scientifique} & \textbf{Niveau de preuve} \\
\hline
Termites contre l'anémie & Forte teneur en fer + folates & Plausible, études cliniques contrôlées encore insuffisantes \\
\hline
Termites et fertilité féminine & Apport en zinc et protéines & Non démontré chez l'homme \\
\hline
Vers de palmier en convalescence & Aliment très énergétique (lipides) et protéiné & Plausible, cohérent avec la diététique de récupération \\
\hline
Chenilles pour la croissance infantile & Richesse en protéines, fer, zinc & Études d'observation favorables (Afrique centrale) \\
\hline
\end{tabularx}
\end{center}

\begin{attention}[Savoir admis ou preuve clinique ?]
Une \textbf{tradition d'usage} n'est pas une \textbf{preuve scientifique}. Pour affirmer qu'un aliment « soigne », il faut une \emph{étude clinique contrôlée} : groupe traité contre groupe témoin, répartition aléatoire, critère de jugement mesurable. À l'écrit du Baccalauréat, emploie les formulations honnêtes : « usage traditionnel rapporté », « effet démontré \emph{in vitro} », « effet démontré chez l'animal », « effet démontré chez l'homme ».
\end{attention}

\begin{savaistu}[La chitine, un prébiotique prometteur]
La chitine et son dérivé déacétylé, le \cle{chitosan}, ne sont pas que des fibres inertes : des travaux récents ont montré \emph{in vitro} et sur modèles murins (souris) des effets \textbf{prébiotiques} (stimulation sélective des \emph{Bifidobacteria} du microbiote intestinal) et \textbf{immunomodulateurs} (activation de macrophages via les récepteurs de reconnaissance des motifs moléculaires). Le chitosan est d'ailleurs déjà utilisé en pharmacie (pansements hémostatiques, excipients). C'est un savoir admis récent, mais les études chez l'homme restent à conduire.
\end{savaistu}

\begin{exoresolu}[Exploiter un tableau de composition centésimale]
\textbf{Énoncé.} Des sauterelles séchées contiennent $8\,\%$ d'eau et \SI{12}{g} de protéines pour \SI{100}{g} de produit frais (avant séchage, l'humidité était de $70\,\%$). (1) Calculer la teneur en protéines sur base matière sèche du produit frais. (2) Le produit séché contient \SI{62}{g} de protéines pour \SI{100}{g}. Comparer avec le bœuf maigre ($\approx \SI{68}{g}$ de protéines pour \SI{100}{g} MS) et conclure. (3) Pourquoi cette comparaison reste-t-elle insuffisante pour conclure à une équivalence nutritionnelle ?
\tcblower
\textbf{Solution.} (1) Produit frais : $MS = 100 - 70 = \SI{30}{g}$. Teneur sur base sèche : $12 \times \dfrac{100}{30} = \SI{40}{g}$ de protéines pour \SI{100}{g} MS.
(2) Produit séché : \SI{62}{g}/\SI{100}{g} MS contre $\approx \SI{68}{g}$ pour le bœuf : les teneurs sont \textbf{du même ordre de grandeur} ; les sauterelles séchées sont donc une source protéique concentrée comparable à la viande.
(3) La quantité ne suffit pas : il faut examiner la \textbf{qualité} (profil en acides aminés essentiels, acide aminé limitant), la \textbf{digestibilité} (réduite par la chitine) et la \textbf{biodisponibilité} des micronutriments. Seul un indice comme le DIAAS permet une conclusion complète.
\end{exoresolu}

\begin{aretenir}[L'essentiel de la section]
\begin{itemize}
\item Les insectes comestibles locaux (termites, vers de palmier, chenilles, sauterelles) présentent des teneurs en protéines de $25$ à $70\,\%$ de la matière sèche, comparables à la viande et au poisson.
\item Leurs protéines sont \textbf{complètes} (9 AAE), riches en lysine et leucine, mais limitées en méthionine/cystine ou tryptophane, et leur digestibilité est réduite par la \textbf{chitine}.
\item Les lipides sont riches en AGPI oméga-3/oméga-6, sans cholestérol notable : un atout cardiovasculaire.
\item Les insectes sont d'excellentes sources de fer, zinc et vitamines B, mais la \textbf{biodisponibilité} doit toujours être discutée avant toute conclusion.
\item Les vertus thérapeutiques traditionnelles doivent être hiérarchisées selon leur niveau de preuve scientifique.
\end{itemize}
\end{aretenir}

\begin{exercice}[Pour t'entraîner]
Un échantillon de vers de palmier frais contient $55\,\%$ d'eau et \SI{13}{g} de protéines pour \SI{100}{g}. (1) Calcule la teneur en protéines sur base matière sèche. (2) Sachant que le score chimique de cette protéine est de $65$ (méthionine limitante), explique en deux phrases ce que cela signifie et propose une association alimentaire pertinente dans un repas camerounais pour compenser cette limite.
\end{exercice}

Cette valeur nutritionnelle exceptionnelle ne se limite pas aux insectes eux-mêmes : certains produits \emph{issus} des insectes, comme le miel, la propolis ou les venins, possèdent des propriétés comestibles et thérapeutiques propres, que nous étudierons dans la section suivante.

\coursec{Produits de la ruche et substances bioactives : Miel, Venins, Propolis}

Dans la section précédente, nous avons montré que les insectes comestibles de notre environnement (termites, chenilles, vers blancs) constituent une source remarquable de protéines, de lipides et de micronutriments capables de combler le déficit protéique des populations. Mais l'abeille, elle, offre bien plus que sa chair : elle produit, transforme et concentre des substances dont les vertus nutritionnelles et thérapeutiques sont connues depuis l'Antiquité et aujourd'hui validées par la science biomédicale. Au marché de Mokolo comme dans les cliniques de Yaoundé, le miel est utilisé pour soigner les plaies ; dans certaines pratiques traditionnelles, les piqûres d'abeilles sont recherchées contre les douleurs articulaires. \textbf{Problème scientifique :} quelles molécules confèrent au miel, au venin et à la propolis leurs propriétés, et par quels mécanismes agissent-elles ? C'est l'objet de cette section, au cœur de l'\cle{apithérapie}.

\begin{definition}[Apithérapie]
L'\cle{apithérapie} est l'ensemble des pratiques thérapeutiques utilisant les produits de la ruche : miel, venin d'abeille (apitoxine), propolis, gelée royale, pollen et cire. Elle va de l'usage traditionnel empirique à des applications médicales contrôlées (miels de grade médical, essais cliniques sur le venin).
\end{definition}

\courssub{Le miel : composition et propriétés physico-chimiques}

\textbf{Observation.} Un miel pur se conserve des années sans réfrigération, sans moisir, alors qu'un sirop de sucre de même concentration fermente en quelques jours. De plus, des pansements au miel accélèrent la cicatrisation de plaies chroniques rebelles aux antibiotiques. \textbf{Hypothèse :} le miel contient des substances antibactériennes actives.

\begin{definition}[Miel]
Le miel est la substance sucrée produite par l'abeille (\emph{Apis mellifera}) à partir du nectar des fleurs ou du miellat, transformé par des enzymes salivaires, concentré par évaporation (dessiccation dans la ruche) et stocké dans les rayons operculés.
\end{definition}

Sa composition moyenne est la suivante :

\begin{center}
\begin{tabularx}{\linewidth}{|l|X|}\hline
\rowcolor{popblueL} \textbf{Constituant} & \textbf{Proportion et rôle} \\ \hline
Glucides & \SI{75}{\percent} à \SI{80}{\percent} : fructose (environ \SI{38}{\percent}), glucose (environ \SI{31}{\percent}), saccharose, maltose \\ \hline
Eau & \SI{17}{\percent} à \SI{20}{\percent} (activité de l'eau $a_w \approx 0,55$ à $0,60$) \\ \hline
Enzymes & \cle{glucose-oxydase} (GOD), \cle{invertase} (saccharase), diastase, catalase \\ \hline
Peptides et polyphénols & \cle{défensine-1} (peptide antimicrobien), flavonoïdes, acides phénoliques \\ \hline
Composés spécifiques & \cle{MGO} (méthylglyoxal) dans le miel de Manuka \\ \hline
Minéraux, acides organiques & K, Ca, Mg ; acide gluconique (responsable du pH acide) \\ \hline
\end{tabularx}
\end{center}

Trois propriétés physico-chimiques expliquent déjà une partie de l'activité biologique :

\begin{itemize}
\item \textbf{pH acide} : compris entre 3,2 et 4,5 (acide gluconique), inhospitalier à la plupart des bactéries pathogènes dont l'optimum de croissance se situe vers pH 7.
\item \textbf{Activité de l'eau ($a_w$) très faible} : la forte concentration en sucres crée une pression osmotique élevée qui déshydrate les cellules bactériennes (plasmolyse).
\item \textbf{HMF} : indicateur de qualité.
\end{itemize}

\begin{definition}[HMF (Hydroxyméthylfurfural)]
L'\cle{HMF} (5-hydroxyméthylfurfural) est un composé formé par déshydratation des sucres en milieu acide, lors du chauffage ou du vieillissement du miel. Un miel frais en contient moins de \SI{15}{mg/kg} ; la norme internationale (Codex Alimentarius) fixe la limite à \SI{40}{mg/kg}. Un taux élevé d'HMF signale un miel chauffé ou trop vieux, dont les enzymes ont été détruites.
\end{definition}

\courssub{Activité antibactérienne du miel : mécanismes peroxyde et non-peroxyde}

\textbf{Mécanisme peroxyde-dépendant.} L'enzyme clé est la \cle{glucose-oxydase} (GOD), sécrétée par les glandes hypopharyngiennes de l'abeille. Dans le miel concentré, elle est inactive ; mais dès que le miel est \emph{dilué} (par les exsudats d'une plaie, par exemple), elle catalyse l'oxydation du glucose :

\[ \ce{Glucose + O2 + H2O ->[GOD] \text{acide gluconique} + H2O2} \]

Le peroxyde d'hydrogène \ce{H2O2} produit en faible quantité continue est un oxydant bactéricide qui détruit les membranes et les protéines bactériennes, sans toxicité pour les tissus humains à ces concentrations.

\begin{propriete}[Activité antibactérienne du miel (savoir admis)]
L'activité antibactérienne d'un miel est corrélée à sa capacité à produire du \ce{H2O2} (activité peroxyde) et/ou à sa teneur en composés non-peroxydiques (MGO, polyphénols, défensine-1). La glucose-oxydase étant \textbf{thermolabile}, un chauffage au-delà de \SI{40}{\celsius} la dénature : le miel perd alors son activité peroxyde. La lumière et la catalase des tissus dégradent aussi le \ce{H2O2}.
\end{propriete}

\textbf{Mécanisme non-peroxyde.} Certains miels, comme le miel de Manuka (Nouvelle-Zélande), conservent une activité antibactérienne même après destruction du \ce{H2O2} par la catalase. Cette activité est due au \cle{MGO} (méthylglyoxal), formé à partir du dihydroxyacétone du nectar, qui altère les protéines et l'ADN bactériens. On la mesure par l'indice UMF ou la teneur en MGO (ex. MGO 400+ = au moins \SI{400}{mg/kg}).

\begin{popfigure}
\begin{center}
\begin{tikzpicture}[font=\small, >=Stealth]
% Schéma réactionnel GOD
\node[draw=popblue, fill=popblueL, rounded corners, minimum width=2.6cm, minimum height=0.9cm] (glu) at (0,0) {Glucose + \ce{O2}};
\node[draw=poporange, fill=poporangeL, rounded corners, minimum width=2.6cm, minimum height=0.9cm, right=2.6cm of glu] (h2o2) {\ce{H2O2} + acide gluconique};
\node[draw=popgreen, fill=popgreenL, rounded corners, minimum width=3.2cm, minimum height=0.9cm, right=2.6cm of h2o2] (effet) {Lyse bactérienne};
\draw[->, very thick, popdark] (glu) -- node[above, popdark]{GOD} node[below, popmuted]{miel dilué} (h2o2);
\draw[->, very thick, popdark] (h2o2) -- node[above, popdark]{oxydation} node[below, popmuted]{membranes, protéines} (effet);
% Inhibition par chauffage
\node[draw=poppink, fill=poppinkL, rounded corners, below=1.1cm of glu, minimum width=3.4cm, minimum height=0.8cm] (chauf) {Chauffage $>$ \SI{40}{\celsius}};
\draw[->, very thick, poppink, dashed] (chauf.east) -- ++(1.2,0) |- node[pos=0.25, right, poppink]{GOD dénaturée} ($(glu.east)!0.5!(h2o2.west)$);
\end{tikzpicture}
\end{center}
\legende{Mécanisme peroxyde-dépendant de l'activité antibactérienne du miel. La glucose-oxydase (GOD), activée par la dilution, produit du peroxyde d'hydrogène \ce{H2O2} bactéricide. Le chauffage au-delà de \SI{40}{\celsius} dénature l'enzyme et abolit cette voie.}
\end{popfigure}

\begin{experience}[Cinétique de production de \ce{H2O2} par différents miels]
\textbf{Protocole :} on dilue trois miels (Manuka MGO 400+, miel camerounais multiflore frais, même miel chauffé à \SI{60}{\celsius} pendant \SI{30}{min}) à \SI{25}{\percent} (m/v) dans de l'eau stérile ; on dose le \ce{H2O2} accumulé (méthode colorimétrique) toutes les heures pendant \SI{6}{h} à \SI{37}{\celsius}. \textbf{Résultats typiques :}
\end{experience}

\begin{popfigure}
\begin{center}
\begin{tikzpicture}
\begin{axis}[
  width=0.85\linewidth, height=6cm,
  xlabel={Temps (h)}, ylabel={\ce{H2O2} accumulé ($\mu$mol/L)},
  xmin=0, xmax=6, ymin=0, ymax=100,
  xtick={0,1,2,3,4,5,6},
  legend pos=north west, legend style={font=\scriptsize},
  grid=both, grid style={popline!40},
]
\addplot[popcurve, popblue, domain=0:6, samples=60]{85*(1-exp(-0.55*x))};
\addlegendentry{Manuka MGO 400+ (peroxyde + MGO)}
\addplot[popcurve, popgreen, domain=0:6, samples=60]{60*(1-exp(-0.45*x))};
\addlegendentry{Miel multiflore frais (Cameroun)}
\addplot[popcurve, poppink, domain=0:6, samples=60]{4*(1-exp(-0.3*x))};
\addlegendentry{Miel multiflore chauffé \SI{60}{\celsius}}
\end{axis}
\end{tikzpicture}
\end{center}
\legende{Cinétique de production de \ce{H2O2} après dilution de trois miels. Le miel frais produit du peroxyde de façon continue ; le chauffage, en dénaturant la glucose-oxydase, réduit la production à presque zéro. Le Manuka conserve en outre une activité non-peroxyde (MGO) non mesurée ici.}
\end{popfigure}

\textbf{Interprétation :} l'activité peroxyde exige une enzyme intacte (miel brut, non chauffé) et une dilution modérée. \textbf{Conclusion :} seuls les miels bruts ou de grade médical (stérilisés par irradiation, qui préserve les enzymes) conviennent à un usage clinique.

\textbf{Utilisation clinique.} Les miels de grade médical sont indiqués dans la cicatrisation des plaies chroniques (ulcères de jambe, escarres, pied diabétique) et des brûlures : ils nettoient la plaie (osmolarité), inhibent les bactéries y compris multirésistantes (SARM), réduisent l'inflammation et maintiennent un milieu humide favorable à la granulation.

\begin{exemplebox}[Concentrations minimales inhibitrices (CMI) mesurées]
Sur \emph{Staphylococcus aureus}, la CMI du miel de Manuka MGO 400+ est d'environ \num{3} à \SI{5}{\percent} (v/v), tandis que celle des miels locaux camerounais (eucalyptus, multiflore) testés dans les laboratoires universitaires se situe autour de \num{10} à \SI{20}{\percent} (v/v). Ces valeurs montrent que les miels locaux sont actifs, mais à des concentrations plus élevées, car leur activité repose surtout sur la voie peroxyde.
\end{exemplebox}

\begin{methode}[Réaliser un antibiogramme au miel (méthode des puits)]
\begin{enumerate}
\item Préparer une gélose Mueller-Hinton coulée en boîtes de Pétri ; laisser solidifier.
\item Ensemencer à la surface une suspension calibrée (\num{0,5} McFarland) de \emph{S. aureus} sur une boîte et de \emph{Pseudomonas aeruginosa} sur une autre.
\item Percer des puits de \SI{6}{mm} de diamètre dans la gélose à l'aide d'un emporte-pièce stérile.
\item Déposer dans chaque puits \SI{100}{\mu L} de miel à tester (brut, dilué, chauffé) ; prévoir un témoin positif (antibiotique) et un témoin négatif (eau stérile).
\item Incuber \SI{24}{h} à \SI{37}{\celsius}, puis mesurer le diamètre des zones d'inhibition autour des puits.
\item Comparer : plus le diamètre est grand, plus l'activité antibactérienne est forte ; comparer aussi les deux souches (\emph{P. aeruginosa} est souvent moins sensible).
\end{enumerate}
\end{methode}

\begin{attention}[Miel commercial et plaies : une erreur dangereuse]
Il ne faut \textbf{jamais} appliquer un miel commercial pasteurisé sur une plaie : le chauffage a détruit la glucose-oxydase (perte de l'activité peroxyde) et le miel n'est pas stérile. De plus, le miel peut véhiculer des spores de \emph{Clostridium botulinum} : il est \textbf{contre-indiqué chez le nourrisson de moins de 12 mois} (risque de botulisme infantile, le tube digestif immature ne pouvant empêcher la germination des spores). En clinique, on utilise exclusivement des miels de grade médical, irradiés et contrôlés.
\end{attention}

\courssub{Le venin d'abeille (apitoxine) : composition et propriétés}

\textbf{Observation.} Les apiculteurs exposés aux piqûres répétées rapportent souvent une amélioration de leurs rhumatismes. \textbf{Hypothèse :} le venin contient des molécules anti-inflammatoires.

\begin{definition}[Venin d'abeille et mellitine]
Le venin d'abeille (\cle{apitoxine}) est un liquide clair, acide (pH $\approx 5$), sécrété par les glandes à venin de l'ouvrière et injecté par l'aiguillon. Son constituant majeur est la \cle{mellitine}, un peptide de 26 acides aminés, de nature \emph{amphipathique} (une face hydrophile, une face hydrophobe), qui représente environ \SI{50}{\percent} du poids sec du venin.
\end{definition}

\begin{popfigure}
\begin{center}
\begin{tikzpicture}[font=\scriptsize]
% Diagramme en secteurs
\def\r{2.2}
% Mellitine 50%
\fill[popblue] (0,0) -- (90:\r) arc (90:-90:\r) -- cycle;
% PLA2 12%
\fill[poporange] (0,0) -- (-90:\r) arc (-90:-133.2:\r) -- cycle;
% Hyaluronidase 3%
\fill[popgreen] (0,0) -- (-133.2:\r) arc (-133.2:-144:\r) -- cycle;
% Apamine 2%
\fill[poppurple] (0,0) -- (-144:\r) arc (-144:-151.2:\r) -- cycle;
% Autres 33%
\fill[popmuted!50] (0,0) -- (-151.2:\r) arc (-151.2:-270:\r) -- cycle;
\draw[popdark] (0,0) circle (\r);
% Légendes
\node[right, align=left] at (3.0,1.6) {\textcolor{popblue}{\rule{0.35cm}{0.35cm}} Mellitine \SI{\sim 50}{\percent}};
\node[right, align=left] at (3.0,1.0) {\textcolor{poporange}{\rule{0.35cm}{0.35cm}} Phospholipase A2 \SI{\sim 12}{\percent}};
\node[right, align=left] at (3.0,0.4) {\textcolor{popgreen}{\rule{0.35cm}{0.35cm}} Hyaluronidase \SI{\sim 3}{\percent}};
\node[right, align=left] at (3.0,-0.2) {\textcolor{poppurple}{\rule{0.35cm}{0.35cm}} Apamine \SI{\sim 2}{\percent}};
\node[right, align=left] at (3.0,-0.8) {\textcolor{popmuted}{\rule{0.35cm}{0.35cm}} Autres (peptide 401, histamine,} ;
\node[right, align=left] at (3.35,-1.2) {dopamine, minimine...) \SI{\sim 33}{\percent}};
\end{tikzpicture}
\end{center}
\legende{Composition moyenne du venin d'abeille (\emph{Apis mellifera}), en pourcentage du poids sec. La mellitine domine largement ; la phospholipase A2 (PLA2) est le principal allergène.}
\end{popfigure}

\textbf{Rôles des principaux constituants :}

\begin{itemize}
\item \textbf{Mellitine} : s'insère dans les bicouches lipidiques et y forme des pores, provoquant la lyse cellulaire (douleur, hémolyse) ; à doses contrôlées, elle stimule la sécrétion de cortisol (effet anti-inflammatoire indirect) et inhibe la voie $NF-\kappa B$.
\item \textbf{Phospholipase A2 (PLA2)} : enzyme qui hydrolyse les phospholipides membranaires ; synergique avec la mellitine ; principal allergène responsable des chocs anaphylactiques.
\item \textbf{Hyaluronidase} : hydrolyse l'acide hyaluronique de la matrice extracellulaire, facilitant la diffusion du venin (« facteur de spreading »).
\item \textbf{Apamine} : neurotoxine de 18 acides aminés bloquant les canaux potassiques \ce{Ca^2+}-dépendants ; à faible dose, effet \textbf{anti-inflammatoire} et neuroprotecteur étudié.
\item \textbf{Peptide 401 (MCD-peptide)} : dégranule les mastocytes mais possède, à doses contrôlées, une activité anti-inflammatoire cent fois supérieure à celle de l'hydrocortisone dans certains modèles.
\end{itemize}

\begin{popfigure}
\begin{center}
\begin{tikzpicture}[font=\scriptsize, >=Stealth]
% Bicouche lipidique
\foreach \x in {0,0.7,...,6.3}{
  \fill[popgold] (\x,0.45) circle (0.14);
  \draw[popgold!70!black] (\x,0.32) -- (\x-0.05,0.0) (\x,0.32) -- (\x+0.05,0.0);
  \fill[popgold] (\x,-0.45) circle (0.14);
  \draw[popgold!70!black] (\x,-0.32) -- (\x-0.05,0.0) (\x,-0.32) -- (\x+0.05,0.0);
}
\node[left, popgold!60!black] at (-0.2,0.45) {têtes hydrophiles};
\node[left, popgold!60!black] at (-0.2,0) {queues hydrophobes};
% Hélice mellitine (représentation cylindrique)
\draw[popblue, very thick, decorate, decoration={coil, aspect=0.4, segment length=4pt, amplitude=5pt}] (1.2,0.9) -- (3.4,0.1);
\draw[popblue, very thick, decorate, decoration={coil, aspect=0.4, segment length=4pt, amplitude=5pt}] (3.6,-0.1) -- (5.2,-0.7);
\node[popblue, above] at (2.2,1.05) {hélice $\alpha$ amphipathique (26 AA)};
% Pore
\draw[->, poppink, very thick] (4.9,1.3) -- node[right, poppink]{insertion + pore $\Rightarrow$ lyse} (4.7,0.35);
\end{tikzpicture}
\end{center}
\legende{Insertion de la mellitine dans une bicouche lipidique. Le peptide amphipathique (face hydrophile / face hydrophobe) s'oriente dans la membrane et s'agrège en pores transmembranaires, entraînant la lyse cellulaire. Cette propriété inspire la nanomédecine anticancéreuse.}
\end{popfigure}

\textbf{Propriétés pharmacologiques et apithérapie.} À doses maîtrisées, le venin présente des propriétés \textbf{anti-inflammatoires} (apamine, mellitine, peptide 401), \textbf{analgésiques} et \textbf{immunomodulatrices}. L'apithérapie par venin (piqûres contrôlées ou injections de venin purifié) est utilisée empiriquement dans l'arthrose et la polyarthrite rhumatoïde ; des essais cliniques explorent son intérêt dans la sclérose en plaques et la maladie de Parkinson, sans preuve définitive à ce jour.

\begin{attention}[Risque anaphylactique]
La PLA2 et d'autres allergènes du venin peuvent déclencher un \textbf{choc anaphylactique} potentiellement mortel chez les personnes sensibilisées. Toute apithérapie par venin exige un test d'allergie préalable et la disponibilité immédiate d'adrénaline. C'est un acte médical, jamais une automédication.
\end{attention}

\begin{savaistu}[De la mellitine à la nanomédecine]
La capacité de la mellitine à perforer sélectivement les membranes inspire la conception de \textbf{nanoparticules vectrices} : des chercheurs ont greffé la mellitine sur des nanovecteurs (« nanobees ») qui ciblent et percent préférentiellement les membranes des cellules tumorales, épargnant les cellules saines. Cette stratégie de médecine ciblée, issue du biomimétisme, illustre comment une toxine d'insecte devient un outil thérapeutique d'avenir.
\end{savaistu}

\courssub{La propolis : la « colle d'abeille » aux vertus antivirales}

\begin{definition}[Propolis et CAPE]
La \cle{propolis} est un mélange résineux collecté par les abeilles sur les bourgeons et écorces (peupliers, conifères, eucalyptus), malaxé avec de la cire et des enzymes salivaires, utilisé pour colmater et assainir la ruche. Elle contient environ \SI{50}{\percent} de résines, \SI{30}{\percent} de cire, \SI{10}{\percent} d'huiles essentielles et \SI{5}{\percent} de pollen. Sa richesse en \textbf{flavonoïdes} (galangine, chrysine, pinocembrine) et en \textbf{acides phénoliques}, dont le \cle{CAPE} (phényléthylester de l'acide caféique), explique ses activités biologiques.
\end{definition}

Les études \emph{in vitro} et \emph{in vivo} (modèles animaux) attribuent à la propolis :

\begin{itemize}
\item une activité \textbf{antivirale} (inhibition de la réplication des virus de la grippe et de l'herpès) ;
\item une activité \textbf{antifongique} (\emph{Candida albicans}) et antibactérienne ;
\item une activité \textbf{anticancéreuse expérimentale} : le CAPE inhibe la prolifération de lignées tumorales et la voie $NF-\kappa B$, sans qu'aucune indication clinique humaine ne soit encore validée ;
\item une activité \textbf{anti-inflammatoire et cicatrisante} (usage traditionnel camerounais contre les maux de gorge et les plaies buccales).
\end{itemize}

\begin{attention}[Prudence d'interprétation]
Une activité observée \emph{in vitro} (sur cellules en culture) ne garantit pas une efficacité chez l'Homme : biodisponibilité, doses et métabolisme diffèrent. Au Baccalauréat comme en santé publique, il faut distinguer \emph{preuve expérimentale} et \emph{allégation thérapeutique}.
\end{attention}

\courssub{Gelée royale et pollen : entre allégations et preuves}

La \textbf{gelée royale}, sécrétée par les glandes hypopharyngiennes des nourrices, est la nourriture exclusive de la reine. Elle contient \SI{60}{\percent} à \SI{70}{\percent} d'eau, \SI{12}{\percent} à \SI{15}{\percent} de protéines (dont les major royal jelly proteins et la royalactine), \SI{10}{\percent} à \SI{16}{\percent} de sucres, des lipides dont l'acide 10-hydroxy-2-décénoïque (10-HDA), et des vitamines du groupe B. Le \textbf{pollen} récolté apporte \SI{20}{\percent} à \SI{25}{\percent} de protéines, des flavonoïdes et des caroténoïdes.

Les allégations courantes (stimulation de l'immunité, régulation métabolique, « revitalisation ») ont été évaluées par l'\textbf{EFSA} (Autorité européenne de sécurité des aliments) : à ce jour, \textbf{aucune allégation santé} relative à la gelée royale ou au pollen n'a été validée faute de preuves cliniques suffisantes. Ces produits restent des compléments alimentaires, non des médicaments.

\begin{exoresolu}[Analyser des données d'antibiogramme au miel]
\textbf{Énoncé.} Un étudiant de Terminale D réalise un antibiogramme par la méthode des puits sur \emph{Staphylococcus aureus}. Il obtient les diamètres d'inhibition suivants : miel de Manuka MGO 400+ : \SI{22}{mm} ; miel multiflore camerounais frais : \SI{14}{mm} ; même miel chauffé à \SI{70}{\celsius} pendant \SI{20}{min} : \SI{0}{mm} ; eau stérile : \SI{0}{mm}. (1) Interpréter chaque résultat. (2) Expliquer la différence entre le miel frais et le miel chauffé. (3) Pourquoi le Manuka resterait-il actif même après ajout de catalase ?
\tcblower
\textbf{Solution.} (1) Le diamètre d'inhibition est proportionnel à l'activité antibactérienne : le Manuka est le plus actif (\SI{22}{mm}), le miel local frais est actif mais plus faiblement (\SI{14}{mm}), le miel chauffé et l'eau sont inactifs (\SI{0}{mm}). Le témoin négatif (eau) valide l'expérience : l'inhibition est bien due au miel. (2) Le chauffage à \SI{70}{\celsius} dénature la glucose-oxydase, enzyme thermolabile : la production de \ce{H2O2} après dilution est abolie, d'où la perte totale d'activité peroxyde du miel local, qui ne possède pas de composé non-peroxyde majeur. (3) La catalase détruit le \ce{H2O2} ; si le Manuka reste actif malgré elle, c'est que son activité repose en grande partie sur le \textbf{MGO} (méthylglyoxal), mécanisme non-peroxyde insensible à la catalase et à la chaleur modérée. \textbf{Conclusion :} l'activité antibactérienne d'un miel dépend de sa nature (peroxyde et/ou MGO) et de son mode de conservation.
\end{exoresolu}

\begin{aretenir}[L'essentiel de la section]
\begin{itemize}
\item Le miel doit son activité antibactérienne à son pH acide, son osmolarité, la production de \ce{H2O2} par la \cle{glucose-oxydase} (voie peroxyde, thermolabile) et, pour certains miels, au \cle{MGO} (voie non-peroxyde).
\item L'\cle{HMF} est le marqueur de chauffage/vieillissement ; un miel thérapeutique doit être brut ou de grade médical, jamais pasteurisé, et interdit au nourrisson.
\item Le venin d'abeille contient \SI{\sim 50}{\percent} de \cle{mellitine} (peptide amphipathique formant des pores), de la PLA2 (allergène majeur), de l'apamine (anti-inflammatoire) ; son usage médical exige un contrôle du risque anaphylactique.
\item La \cle{propolis}, riche en flavonoïdes et en \cle{CAPE}, présente des activités antivirale, antifongique et anticancéreuse expérimentale.
\item Gelée royale et pollen : compléments nutritionnels dont les allégations santé ne sont pas validées par l'EFSA.
\end{itemize}
\end{aretenir}

Ainsi, la ruche apparaît comme une véritable \emph{pharmacie naturelle} dont les molécules (enzymes, peptides, polyphénols) sont aujourd'hui décryptées par la biochimie. Mais l'apport des insectes à la production animale ne s'arrête pas aux produits thérapeutiques : certains insectes protègent directement les cultures fourragères et vivrières. C'est l'objet de la section suivante, consacrée à la \textbf{lutte biologique} contre les ravageurs.

\coursec{Lutte biologique : Insectes auxiliaires et protection des cultures fourragères}

Dans la section précédente, nous avons vu que les insectes nous offrent directement des produits précieux : le miel, la propolis, les venins aux vertus pharmacologiques. Mais les insectes peuvent aussi travailler \emph{pour} nous, de manière indirecte, dans les champs. Imagine un champ de niébé ou de maïs fourrager envahi de pucerons : l'agriculteur peut pulvériser un insecticide de synthèse... ou laisser faire les coccinelles, les chrysopes et les guêpes parasitoïdes qui patrouillent déjà gratuitement dans sa parcelle. Cette seconde stratégie, qui consiste à utiliser des \cle{ennemis naturels} des ravageurs, s'appelle la \cle{lutte biologique}. Elle protège les cultures fourragères (maïs, soja, légumineuses) qui nourrissent le bétail, tout en préservant la qualité de l'environnement et la santé des animaux d'élevage.

\courssub{Ravageurs, auxiliaires et stratégies de lutte biologique}

\textbf{Observation introductive.} Dans une parcelle de maïs fourrager à Obala (région du Centre), on observe deux situations : dans la parcelle A, traitée chaque semaine avec un insecticide à large spectre, les pucerons réapparaissent massivement dix jours après chaque traitement ; dans la parcelle B, non traitée mais bordée de bandes fleuries, la population de pucerons reste faible et stable. Hypothèse : la parcelle B abrite des ennemis naturels des pucerons, éliminés par l'insecticide dans la parcelle A. Le comptage des coccinelles confirme l'hypothèse : \num{12} coccinelles par m² en B, moins de \num{1} par m² en A.

\begin{definition}[Ravageur et auxiliaire de culture]
Un \cle{ravageur} est un organisme (insecte, acarien, nématode...) qui s'attaque aux organes d'une plante cultivée et en diminue le rendement ou la qualité. Exemples : le puceron du niébé, le foreur des tiges du maïs (\emph{Busseola fusca}), la chenille légionnaire (\emph{Spodoptera frugiperda}).
\par
Un \cle{auxiliaire de culture} est un organisme qui limite naturellement les populations de ravageurs. On distingue :
\begin{itemize}
\item le \cle{prédateur} : il tue et consomme directement de nombreuses proies au cours de sa vie (coccinelle, chrysope, syrphe, libellule) ;
\item le \cle{parasitoïde} : il pond dans ou sur un hôte ; sa larve se développe aux dépens de cet hôte unique, qu'elle finit par tuer (guêpes parasitoïdes \emph{Trichogramma}, \emph{Cotesia}). C'est un intermédiaire entre prédateur et parasite.
\end{itemize}
\end{definition}

\begin{definition}[Les trois stratégies de lutte biologique]
La \cle{lutte biologique} est l'ensemble des méthodes utilisant des ennemis naturels pour maintenir les populations de ravageurs sous le seuil de nuisance. On distingue :
\begin{enumerate}
\item la lutte biologique par \cle{conservation} : on protège et on favorise les auxiliaires déjà présents (bandes fleuries, réduction des pesticides, haies) ;
\item la lutte biologique par \cle{acclimatation} (ou importation) : on introduit un auxiliaire exotique, après étude d'impact, contre un ravageur lui-même introduit ;
\item la lutte biologique par \cle{inondation} (lâchers massifs) : on élève des auxiliaires en laboratoire et on les relâche en grand nombre au moment critique, comme un « biopesticide vivant ». Exemple : lâchers de \emph{Trichogramma} contre les foreurs du maïs.
\end{enumerate}
\end{definition}

\begin{propriete}[Conditions d'efficacité de la lutte biologique]
L'efficacité d'un programme de lutte biologique dépend de quatre facteurs :
\begin{enumerate}
\item la \cle{spécificité} de l'auxiliaire vis-à-vis du ravageur cible (il ne doit pas s'attaquer aux espèces utiles) ;
\item sa \cle{capacité de dispersion} et de recherche de la proie dans la parcelle ;
\item la \cle{synchronisation phénologique} : l'auxiliaire doit être présent et actif au moment où le ravageur apparaît ;
\item l'\cle{absence de pesticides à large spectre}, qui tuent les auxiliaires plus durablement que les ravageurs (les ravageurs, à cycle court, recolonisent plus vite : c'est le phénomène de \emph{résurgence}).
\end{enumerate}
\end{propriete}

\courssub{La coccinelle, prédateur emblématique des pucerons}

\textbf{Observation et données expérimentales.} Au laboratoire, on place une larve de coccinelle à sept points (\emph{Coccinella septempunctata}) dans une boîte contenant \num{100} pucerons du niébé, renouvelés chaque jour. Résultat moyen : la larve consomme \num{50} à \num{100} pucerons par jour, soit environ \num{400} pucerons au cours de son développement larvaire complet (environ \num{3} semaines). L'adulte, lui aussi prédateur, en consomme encore \num{50} à \num{100} par jour.

\begin{exemplebox}[Voracité de la coccinelle à sept points]
Une seule larve de \emph{Coccinella septempunctata} dévore environ \textbf{\num{400} pucerons} durant son développement ; un adulte en consomme jusqu'à \textbf{\num{100} par jour}. Une colonie de quelques dizaines de coccinelles suffit donc à juguler une pullulation de pucerons sur une parcelle familiale de niébé ou de maïs fourrager.
\end{exemplebox}

\textbf{Cycle de vie.} La coccinelle est un insecte à métamorphose complète (holométabole) : œuf (pondu en grappes jaunes près des colonies de pucerons) $\rightarrow$ larve (noire et orange, allongée, très vorace — c'est le stade le plus utile) $\rightarrow$ nymphe (fixée sur une feuille) $\rightarrow$ adulte (élytres rouges à points noirs). Tous les stades mobiles, larve et adulte, sont prédateurs de pucerons, de cochenilles et d'acariens.

\begin{attention}[Le piège de l'espèce exotique : \emph{Harmonia axyridis}]
Erreur fréquente : introduire un auxiliaire exotique sans étude d'impact préalable. La coccinelle asiatique \emph{Harmonia axyridis}, introduite dans plusieurs pays pour lutter contre les pucerons, est devenue \textbf{invasive} : elle entre en compétition avec les coccinelles locales (comme \emph{Cheilomenes} spp. en Afrique), consomme leurs œufs, et envahit les habitations en automne, causant des nuisances. Toute introduction doit être précédée d'une évaluation du risque ; la stratégie de \textbf{conservation} des auxiliaires locaux est toujours préférable.
\end{attention}

\courssub{Libellules, chrysopes, syrphes et guêpes parasitoïdes}

\textbf{Les libellules (Odonates).} Prédateurs à double vie : l'adulte, redoutable chasseur aérien, capture en vol moustiques, mouches et petits insectes ; sa larve (naïade), aquatique, dévore les larves de moustiques dans les mares et rizières. Un couple de libellules peut capturer environ \textbf{\num{500} moustiques par jour}. Intérêt indirect pour l'élevage : en réduisant les populations de moustiques et de mouches autour des pâturages et des points d'eau, les libellules limitent la pression des vecteurs de maladies animales (et humaines, comme le paludisme).

\textbf{Les chrysopes (lions des pucerons).} Leur larve, armée de mandibules crochues, aspire le contenu des pucerons, cochenilles et œufs de chenilles : jusqu'à \num{600} pucerons durant son développement. L'adulte, délicat, aux ailes dentelées vertes, se nourrit de nectar et de pollen — d'où l'intérêt des bandes fleuries.

\textbf{Les syrphes (mouches « déguisées en guêpes »).} L'adulte, butineur inoffensif mimant la guêpe, pond près des colonies de pucerons ; sa larve (petite limace verdâtre) en consomme \num{200} à \num{800} durant sa vie.

\textbf{Les guêpes parasitoïdes.} Minuscules (souvent moins de \SI{1}{mm}), ce sont les auxiliaires les plus utilisés en lâchers massifs :
\begin{itemize}
\item \emph{Trichogramma} spp. : parasitoïde des œufs du foreur des tiges du maïs ; la larve se développe dans l'œuf du ravageur, qui n'éclot jamais. Lâchers de \num{100000} à \num{300000} individus par hectare au moment de la ponte des foreurs ;
\item \emph{Cotesia} spp. : parasitoïde des chenilles (dont la légionnaire \emph{Spodoptera frugiperda}, ravageur majeur du maïs au Cameroun depuis 2016) ; la chenille parasitée porte à terme de petits cocons blancs de soie fixés sur son corps.
\end{itemize}

\begin{popfigure}
\begin{center}
\begin{tikzpicture}
\begin{axis}[
  width=0.95\linewidth, height=6.5cm,
  xlabel={Temps (semaines)}, ylabel={Effectif (individus par m$^2$)},
  xmin=0, xmax=20, ymin=0, ymax=60,
  legend style={at={(0.98,0.97)}, anchor=north east, font=\small},
  grid=both, grid style={popline!40},
  axis line style={popdark}
]
\addplot[popcurve, poporange, very thick, domain=0:20, samples=200]
  {30 + 25*sin(deg(3.14159*x/6.5))};
\addlegendentry{Pucerons (proies)}
\addplot[popcurve, popblue, very thick, domain=0:20, samples=200]
  {18 + 14*sin(deg(3.14159*(x-2)/6.5))};
\addlegendentry{Coccinelles (prédateurs)}
\addplot[poporange, dashed, domain=0:20, samples=60]
  {30 + 22*sin(deg(3.14159*x/6.5)) + 3*sin(deg(x*2))};
\addplot[popblue, dashed, domain=0:20, samples=60]
  {18 + 12*sin(deg(3.14159*(x-2)/6.5)) + 2*sin(deg(x*2))};
\end{axis}
\end{tikzpicture}
\end{center}
\legende{Dynamique proie--prédateur dans un champ de légumineuses fourragères (modèle de Lotka--Volterra simplifié, traits pleins ; données de comptage au champ, pointillés). La population de pucerons (orange) culmine, puis celle des coccinelles (bleu) suit avec un \textbf{décalage temporel} d'environ 2 semaines : l'augmentation des proies permet la reproduction des prédateurs, dont la pression fait ensuite chuter les pucerons. Ce déphasage illustre la \textbf{synchronisation phénologique}, condition d'efficacité de la lutte biologique.}
\end{popfigure}

\courssub{Reconnaître les auxiliaires au champ : clé de détermination}

Confondre une larve de coccinelle (auxiliaire précieuse) avec une chenille ravageuse conduit à détruire ses propres alliés. Voici une clé simplifiée pour les stades larvaires rencontrés sur les cultures fourragères.

\begin{popfigure}
\begin{center}
\begin{tikzpicture}[
  font=\small,
  question/.style={rectangle, rounded corners, draw=popblue, fill=popblueL, text width=4.6cm, align=center, inner sep=4pt},
  result/.style={rectangle, rounded corners, draw=popgreen, fill=popgreenL, text width=3.4cm, align=center, inner sep=4pt},
  rav/.style={rectangle, rounded corners, draw=poporange, fill=poporangeL, text width=3.4cm, align=center, inner sep=4pt},
  lbl/.style={font=\footnotesize\itshape, midway, above},
  node distance=1.1cm and 1.6cm
]
\node[question] (q1) {Larve mobile sur feuilles, munie de pattes visibles ?};
\node[question, below left=of q1] (q2) {Corps allongé, noir et orange, mandibules saillantes ?};
\node[question, below right=of q1] (q3) {Corps mou, sans pattes (ou fausses pattes), vert ou jaunâtre ?};
\node[result, below left=of q2] (coc) {\textbf{Larve de coccinelle} : auxiliaire, à protéger};
\node[result, below right=of q2] (chr) {\textbf{Larve de chrysope} : corps grisâtre, mandibules en crochets : auxiliaire};
\node[result, below left=of q3] (syr) {\textbf{Larve de syrphe} : limace verdâtre sans pattes, dans les colonies de pucerons : auxiliaire};
\node[rav, below right=of q3] (che) {\textbf{Chenille} (foreur, légionnaire) : 5 paires de fausses pattes, galeries dans les tiges : ravageur};
\draw[->, popdark] (q1) -- node[lbl]{oui} (q2);
\draw[->, popdark] (q1) -- node[lbl]{non} (q3);
\draw[->, popdark] (q2) -- node[lbl]{oui} (coc);
\draw[->, popdark] (q2) -- node[lbl]{non} (chr);
\draw[->, popdark] (q3) -- node[lbl]{sans pattes} (syr);
\draw[->, popdark] (q3) -- node[lbl]{fausses pattes} (che);
\end{tikzpicture}
\end{center}
\legende{Clé de détermination simplifiée des stades larvaires rencontrés sur maïs et légumineuses fourragères. Vert : auxiliaires à conserver ; orange : ravageur à combattre. Règle pratique : toute larve trouvée \emph{au milieu d'une colonie de pucerons} est presque toujours un auxiliaire.}
\end{popfigure}

\courssub{Mettre en œuvre la lutte biologique par conservation}

\begin{methode}[Installer des bandes fleuries en bordure de champ fourrager]
\begin{enumerate}
\item Choisir une bordure de \SI{1}{m} à \SI{3}{m} de large, ensoleillée, le long de la parcelle de maïs ou de niébé.
\item Semer un mélange de plantes à fleurs locales et à floraison échelonnée (souci, cosmos, tournesol, sésame, légumineuses spontanées) : le nectar et le pollen nourrissent les adultes de chrysopes, syrphes et parasitoïdes.
\item Éviter tout traitement insecticide sur la bande fleurie et raisonner les traitements sur la culture (seuil de nuisance, produits sélectifs).
\item Faucher tardivement et par moitiés alternées pour conserver un refuge permanent aux auxiliaires.
\item Suivre l'efficacité : compter chaque semaine pucerons et auxiliaires sur \num{20} plantes au hasard.
\end{enumerate}
\end{methode}

\begin{methode}[Fabriquer un abri à coccinelles (hôtel à insectes)]
\begin{enumerate}
\item Assembler une caisse en bois non traité (environ \SI{30}{cm} $\times$ \SI{30}{cm} $\times$ \SI{15}{cm}), ouverte sur une face, toit légèrement incliné.
\item Remplir de matériaux variés : paille, pommes de pin, bûches percées de trous de \SI{3}{mm} à \SI{8}{mm}, tiges creuses (bambou, tiges de maïs).
\item Installer l'abri à \SI{30}{cm}--\SI{50}{cm} du sol, face ouverte orientée au sud-est, à l'abri des pluies dominantes, près de la bande fleurie.
\item Ne jamais déplacer l'abri en hiver : les coccinelles y hibernent en groupes et recoloniseront la parcelle au printemps.
\end{enumerate}
\end{methode}

\courssub{Intérêt pour la production animale}

La lutte biologique boucle la logique de cette leçon : en protégeant les cultures fourragères (maïs, soja, niébé) sans pesticides à large spectre, elle garantit un \textbf{aliment abondant et sain pour le bétail}, réduit les résidus chimiques dans le lait et la viande, et préserve les pollinisateurs et les insectes comestibles étudiés dans les sections précédentes. C'est un pilier de l'intensification écologique de l'élevage camerounais.

\begin{savaistu}[La mouche soldat noire, auxiliaire inattendue des élevages]
La mouche soldat noire (\emph{Hermetia illucens}) adulte ne pique pas et ne se nourrit pas : elle ne transmet donc aucune maladie. Mieux, sa présence dans les déjections et les déchets organiques des fermes fait \textbf{fuir les mouches domestiques} (muscidés), vectrices de pathogènes, par un effet répulsif olfactif. Ses larves, qui dégradent les déchets, sont en outre une source de protéines pour l'alimentation des volailles et des poissons : un insecte à la fois auxiliaire sanitaire et ressource fourragère !
\end{savaistu}

\begin{exoresolu}[Calculer un lâcher inondatif de \emph{Trichogramma}]
Énoncé. Un éleveur de Ngaoundéré cultive \SI{2}{ha} de maïs fourrager. Pour lutter contre le foreur des tiges, on recommande des lâchers de \emph{Trichogramma} à raison de \num{150000} individus par hectare, renouvelés \num{3} fois à une semaine d'intervalle. Une trichocarte (carton porteur) contient \num{1000} parasitoïdes et coûte \SI{500}{FCFA}. Calcule le nombre de trichocartes et le coût total du traitement.
\tcblower
Solution détaillée.
\begin{align*}
\text{Besoin par lâcher} &= \num{150000} \times 2 = \num{300000} \text{ individus}\\
\text{Besoin total (3 lâchers)} &= \num{300000} \times 3 = \num{900000} \text{ individus}\\
\text{Nombre de trichocartes} &= \frac{\num{900000}}{\num{1000}} = \num{900} \text{ cartes}\\
\text{Coût total} &= \num{900} \times \SI{500}{FCFA} = \SI{450000}{FCFA}
\end{align*}
Soit \SI{225000}{FCFA} par hectare pour trois lâchers, sans résidu chimique dans le fourrage — un investissement à comparer au coût et aux risques sanitaires d'un insecticide de synthèse.
\end{exoresolu}

\begin{aretenir}[L'essentiel de la section]
\begin{itemize}
\item La lutte biologique utilise les ennemis naturels des ravageurs : \textbf{conservation} (bandes fleuries, abris), \textbf{acclimatation} (introduction encadrée) ou \textbf{inondation} (lâchers massifs, ex. \emph{Trichogramma}).
\item Les auxiliaires sont des \textbf{prédateurs} (coccinelles : \num{400} pucerons par larve ; chrysopes ; syrphes ; libellules : \num{500} moustiques par jour et par couple) ou des \textbf{parasitoïdes} (guêpes \emph{Trichogramma}, \emph{Cotesia}).
\item Efficacité = spécificité + dispersion + synchronisation + absence de pesticides à large spectre.
\item Jamais d'introduction d'espèce exotique (ex. \emph{Harmonia axyridis}) sans étude d'impact.
\item Enjeu pour l'élevage : fourrage protégé, sain et sans résidus.
\end{itemize}
\end{aretenir}

\begin{exercice}[Pour t'entraîner]
Dans une parcelle de niébé, on compte en moyenne \num{40} pucerons par plante et \num{0,5} larve de coccinelle par plante. Sachant qu'une larve consomme \num{80} pucerons par jour, la population d'auxiliaires est-elle suffisante pour endiguer la pullulation en une journée ? Quelle stratégie (conservation ou inondation) préconises-tu, et pourquoi ?
\end{exercice}

\coursec{Techniques d'entomoculture : De la collecte à l'élevage industriel pour l'alimentation animale}

Dans la section précédente, nous avons vu comment les insectes auxiliaires (coccinelles, libellules) protègent les cultures fourragères et vivrières en limitant les ravageurs, sans pesticides chimiques. Mais les insectes peuvent faire plus pour l'éleveur : au lieu de seulement \emph{protéger} la production animale, ils peuvent la \emph{nourrir}. Comment passer de la cueillette saisonnière des termites ailés à une production contrôlée, toute l'année, de protéines d'insectes destinées aux volailles, aux poissons et aux porcs ? C'est l'objet de l'\cle{entomoculture}, c'est-à-dire l'élevage raisonné des insectes.

\courssub{La collecte traditionnelle : une ressource précieuse mais limitée}

\textbf{Observation.} Aux premières pluies de la saison humide, dans les villages de l'Adamaoua ou du Centre, les enfants placent des récipients d'eau sous les lampadaires ou les feux : les termites ailés, attirés par la lumière, s'y abattent en masse. En une nuit, une famille peut récolter plusieurs kilogrammes de termites. Pourquoi alors ne pas nourrir les poulets toute l'année avec cette manne gratuite ?

\begin{experience}[Techniques traditionnelles de collecte des insectes comestibles]
Trois techniques dominent au Cameroun :
\begin{enumerate}
\item \textbf{Piégeage lumineux} : les insectes à phototropisme positif (termites ailés, sauterelles) sont attirés la nuit par une source lumineuse placée au-dessus d'une bassine d'eau ou d'un piège à entonnoir. Rendement : jusqu'à \SI{5}{kg} par nuit de vol nuptial.
\item \textbf{Écorçage des palmiers} : les vers de palmier (larves du charançon \emph{Rhynchophorus phoenicis}) se développent dans les troncs de palmiers abattus ou blessés. On incise l'écorce et on récolte les larves après 4 à 6 semaines.
\item \textbf{Cueillette manuelle} : les chenilles comestibles (imbrasis, \emph{Imbrasia} spp.) sont ramassées une à une sur leurs arbres hôtes en saison sèche.
\end{enumerate}
\end{experience}

\textbf{Interprétation et limites.} Ces techniques sont efficaces mais présentent trois limites majeures :
\begin{itemize}
\item la \cle{saisonnalité} : les vols de termites ne durent que quelques nuits par an ; impossible de garantir un approvisionnement régulier d'un élevage de volailles ;
\item la \cle{pression sur les populations sauvages} : la récolte intensive non gérée appauvrit les colonies et menace la ressource (et les services écosystémiques associés, comme la pollinisation ou la décomposition) ;
\item les \cle{risques de contamination} : les insectes collectés en bordure de champs traités peuvent accumuler des pesticides ; ceux collectés près des axes routiers ou des zones minières peuvent concentrer des métaux lourds (plomb, cadmium).
\end{itemize}

\begin{attention}[La collecte sauvage n'est pas une filière contrôlée]
Contrairement à un élevage, la collecte ne permet \textbf{aucun contrôle du régime alimentaire} de l'insecte. Or un insecte est ce qu'il mange : les contaminants de son substrat se retrouvent dans son corps, puis dans l'animal qui le consomme, puis dans l'assiette du consommateur. C'est le phénomène de \cle{bioaccumulation}.
\end{attention}

La conclusion s'impose : pour intégrer les insectes à l'alimentation animale de façon fiable et sûre, il faut les \emph{produire} en milieu contrôlé. Trois espèces dominent aujourd'hui l'entomoculture mondiale et africaine : la mouche soldat noire, le grillon et le ver de farine.

\courssub{L'élevage de la Mouche Soldat Noire (\emph{Hermetia illucens}) : la bioconversion des déchets}

\textbf{Situation-problème.} Un marché de Yaoundé produit chaque jour des tonnes de déchets organiques (fruits avariés, fanes, fientes de volailles des poulaillers voisins) coûteux à évacuer et sources de miasmes. Et si ces déchets devenaient la matière première d'une usine à protéines ? C'est exactement ce que réalise la larve de la mouche soldat noire (MSN).

\begin{definition}[Bioconversion]
La \cle{bioconversion} est la transformation, par un organisme vivant (ici la larve d'insecte), d'une matière organique de faible valeur (déchets) en biomasse de haute valeur (protéines et lipides) et en sous-produits utiles (engrais organique).
\end{definition}

\textbf{Principe de l'élevage.} Les adultes pondent près du substrat ; les œufs éclosent et les larves se nourrissent voracement des déchets organiques : fientes de poules, déchets de marchés et de cuisines, drêches de brasserie (résidus de grains après brassage). Le cycle complet, de l'œuf à la larve mature (pré-pupe), dure \textbf{14 à 18 jours} à \SIrange{27}{30}{\celsius}. Les larves récoltées contiennent \textbf{40 à 45 \% de protéines} et environ \textbf{30 \% de lipides} sur matière sèche (MS) : une composition comparable à la farine de poisson.

\begin{definition}[Frass, farine d'insectes et hydrolysat]
\begin{itemize}
\item Le \cle{frass} est le mélange des déjections larvaires, des exuvies (mues) et du substrat résiduel non consommé. Stabilisé, c'est un excellent \textbf{engrais organique} riche en azote, phosphore et potassium.
\item La \cle{farine d'insectes} est obtenue par séchage, dégraissage partiel et broyage des larves : c'est un concentré protéique stable et transportable.
\item Un \cle{hydrolysat protéique} est obtenu par hydrolyse enzymatique des protéines d'insectes en peptides courts : très digestible, il est utilisé dans les aliments de démarrage (alevins, porcelets sevrés).
\end{itemize}
\end{definition}

\begin{popfigure}
\begin{tikzpicture}[font=\small, node distance=0.9cm,
  box/.style={draw=popblue, fill=popblueL, rounded corners, minimum height=1cm, text width=3.1cm, align=center},
  outbox/.style={draw=popgreen, fill=popgreenL, rounded corners, minimum height=0.9cm, text width=3.3cm, align=center},
  fleche/.style={-{Stealth[length=3mm]}, thick, popdark}]
\node[box, align=center] (dechet){\textbf{Entrées}\\Déchets organiques\\(fientes, marchés,\\drêches de brasserie)};
\node[box, right=2.2cm of dechet, fill=poporangeL, draw=poporange, align=center] (bio){\textbf{Bioréacteur}\\Larves de MSN\\(14--18 jours, \SIrange{27}{30}{\celsius})};
\node[outbox, right=2.2cm of bio, yshift=1.5cm, align=center] (larves){\textbf{Larves}\\vivantes ou séchées\\$\rightarrow$ farine (volaille, poisson)};
\node[outbox, right=2.2cm of bio, align=center] (frass){\textbf{Frass}\\engrais organique\\(cultures fourragères)};
\node[outbox, right=2.2cm of bio, yshift=-1.5cm, align=center] (lix){\textbf{Lixiviat}\\à traiter / arroser\\compost};
\draw[fleche] (dechet) -- (bio);
\draw[fleche] (bio) -- (larves);
\draw[fleche] (bio) -- (frass);
\draw[fleche] (bio) -- (lix);
\node[below=0.35cm of bio, align=center, text=popmuted] (indic) {Indicateurs : taux de réduction des déchets \num{50}--\num{70} \% ;\\FCR larvaire $\approx$ \num{1,5}--\num{2,0} ; rendement $\approx$ \SI{150}{g} larves fraîches/kg déchets};
\end{tikzpicture}
\legende{Schéma de flux d'une unité de bioconversion par la mouche soldat noire (MSN). Les déchets organiques entrent dans le bioréacteur où les larves les transforment ; les sorties sont la biomasse larvaire (protéines), le frass (engrais) et un lixiviat à gérer.}
\end{popfigure}

\begin{propriete}[Assainissement du substrat par les larves de MSN (savoir admis)]
La larve de \emph{Hermetia illucens} sécrète dans son milieu des \cle{peptides antimicrobiens} (AMP) qui réduisent fortement les populations de bactéries pathogènes (\emph{Salmonella} sp., \emph{Escherichia coli}) dans le substrat. Cette propriété contribue à la \textbf{sécurité sanitaire} de la farine produite, à condition que le substrat soit par ailleurs contrôlé.
\end{propriete}

\textbf{Quel substrat choisir ?} L'expérience suivante compare la croissance des larves sur trois substrats disponibles localement.

\begin{popfigure}
\begin{center}
\begin{tikzpicture}
\begin{axis}[width=0.9\linewidth, height=6.5cm,
  xlabel={Durée d'élevage (\si{\day})}, ylabel={Poids moyen d'une larve (\si{\milli\gram})},
  xmin=0, xmax=14, ymin=0, ymax=230,
  xtick={0,2,4,6,8,10,12,14}, ytick={0,50,100,150,200},
  legend pos=north west, legend style={font=\scriptsize}, grid=both, grid style={popline!40}]
\addplot[color=popblue, very thick, mark=*, smooth] coordinates {(0,2) (2,15) (4,45) (6,95) (8,150) (10,190) (12,215) (14,220)};
\addlegendentry{Déchets de cuisine}
\addplot[color=poporange, very thick, mark=square*, smooth] coordinates {(0,2) (2,12) (4,38) (6,80) (8,125) (10,160) (12,180) (14,185)};
\addlegendentry{Drêches de brasserie}
\addplot[color=popgreen, very thick, mark=triangle*, smooth] coordinates {(0,2) (2,8) (4,25) (6,55) (8,90) (10,120) (12,140) (14,148)};
\addlegendentry{Fientes de poules}
\end{axis}
\end{tikzpicture}
\end{center}
\legende{Courbes de croissance des larves de MSN sur trois substrats (données expérimentales types). Les déchets de cuisine donnent la croissance la plus rapide et le poids final le plus élevé (environ \SI{220}{mg} à J14) ; les fientes de poules, plus pauvres en nutriments digestibles, plafonnent vers \SI{148}{mg}.}
\end{popfigure}

\textbf{Interprétation.} Les trois courbes ont une allure sigmoïde : phase de latence (J0--J2), croissance exponentielle (J2--J10), puis ralentissement à l'approche de la pré-pupe. Le substrat optimal est celui qui combine richesse nutritionnelle et humidité convenable : ici les \textbf{déchets de cuisine}. La récolte doit se faire vers J14, au maximum de biomasse, avant la nymphose.

\begin{attention}[Le piège des déchets non contrôlés]
Erreur fréquente : nourrir les larves avec n'importe quel déchet. Les substrats contaminés par des \textbf{métaux lourds}, des \textbf{mycotoxines} (aflatoxines des arachides moisies) ou des \textbf{résidus de médicaments vétérinaires} conduisent à une bioaccumulation dans la farine, puis à un risque pour toute la chaîne alimentaire. Règle d'or : \textbf{un substrat tracé et contrôlé}, comme pour n'importe quel aliment du bétail.
\end{attention}

\courssub{L'élevage du grillon (\emph{Acheta domesticus}, \emph{Gryllus bimaculatus})}

\textbf{Protocole.} Les grillons sont élevés en bacs plastiques ventilés, garnis de \textbf{cartons alvéolés} (plaques d'œufs) qui multiplient la surface de refuge et limitent le cannibalisme. L'alimentation associe farine de céréales, son de blé et légumes frais (source d'eau). La température optimale est de \SIrange{28}{32}{\celsius}. La récolte intervient à maturité, au bout de \textbf{6 à 8 semaines} : les grillons sont tués par \textbf{congélation} (euthanasie éthique par abaissement progressif de la température, l'insecte étant poïkilotherme), puis séchés et broyés en farine.

\courssub{L'élevage du ver de farine (\emph{Tenebrio molitor})}

\textbf{Protocole.} Les vers de farine (larves du ténébrion) sont élevés sur un substrat de \textbf{son ou de farine}, complété par une source d'humidité (rondelles de carotte ou morceaux de pomme, renouvelés pour éviter les moisissures). La conduite de l'élevage exige de \textbf{gérer les stades} : les adultes (coléoptères) pondent dans le son ; on sépare régulièrement les œufs et jeunes larves des adultes (qui peuvent les dévorer), et on retire les nymphes, immobiles et vulnérables. Ordre de grandeur du rendement : \textbf{\SI{2}{kg} de substrat fournissent environ \SI{1}{kg} de larves fraîches} (FCR $\approx$ 2).

\begin{methode}[Calculer le rendement massique et le FCR d'un élevage de grillons]
À partir du suivi hebdomadaire (pesées, mortalités) :
\begin{enumerate}
\item \textbf{Biomasse récoltée} $B$ : peser la totalité des insectes récoltés (kg).
\item \textbf{Aliment consommé} $A$ : aliment distribué $-$ refus et restes (kg).
\item \textbf{Rendement massique} : $R = \dfrac{B}{A_{\text{distribué}}}$, exprimé en kg d'insectes par kg d'aliment distribué (ou en \%).
\item \textbf{Indice de conversion (FCR)} : $FCR = \dfrac{A}{B - B_0}$, où $B_0$ est la biomasse initiale (jeunes mis en élevage). Plus le FCR est faible, plus l'élevage est efficace.
\item \textbf{Taux de mortalité} : $M = \dfrac{N_0 - N_f}{N_0} \times 100$ (\%), à suivre pour détecter un problème (température, humidité, densité).
\end{enumerate}
\end{methode}

\begin{exoresolu}[FCR d'un élevage de grillons]
Un éleveur de Bafoussam met en bac \SI{200}{g} de jeunes grillons. En 7 semaines, il distribue \SI{9,0}{kg} d'aliment (dont \SI{0,5}{kg} de restes) et récolte \SI{4,6}{kg} de grillons, avec une mortalité de 8 \%. Calcule le FCR et le rendement massique.
\tcblower
Aliment réellement consommé : $A = 9{,}0 - 0{,}5 = 8{,}5$ kg.\\
Gain de biomasse : $\Delta B = 4{,}6 - 0{,}2 = 4{,}4$ kg.\\
\[ FCR = \frac{A}{\Delta B} = \frac{8{,}5}{4{,}4} \approx 1{,}9 \qquad R = \frac{4{,}6}{9{,}0} \approx 0{,}51 \text{ kg/kg} \]
Interprétation : il faut environ \num{1,9} kg d'aliment pour produire \SI{1}{kg} de grillon, soit une efficacité bien supérieure à celle du bovin (FCR 6 à 8). La mortalité de 8 \% reste acceptable (seuil d'alerte : 15 \%).
\end{exoresolu}

\courssub{Utilisation des farines d'insectes en alimentation animale}

Les farines d'insectes remplacent partiellement ou totalement la \textbf{farine de poisson} et le \textbf{tourteau de soja}, dont les prix flambent et dont la disponibilité est limitée au Cameroun. Les essais menés sur poulets de chair, tilapias et porcelets montrent, à substitution raisonnée (25 à 100 \% de la farine de poisson), des \textbf{résultats zootechniques comparables} : gain moyen quotidien (GMQ) maintenu, indice de consommation (IC) stable, rendement en carcasse inchangé, voire meilleure santé intestinale grâce à la chitine et aux peptides antimicrobiens.

\begin{exemplebox}[Substitution chez le tilapia : un exemple chiffré]
Dans des essais conduits dans la filière IRAD/WorldFish, la substitution de \textbf{50 \% de la farine de poisson par de la farine de MSN} dans l'aliment du tilapia a donné : un \textbf{GMQ identique} au témoin, une baisse du \textbf{coût de l'aliment d'environ 15 \%}, et une réduction de l'\textbf{empreinte carbone de l'aliment d'environ 40 \%} (les déchets valorisés évitent à la fois l'importation de farine de poisson et la mise en décharge).
\end{exemplebox}

\begin{savaistu}[Un cadre réglementaire en pleine construction]
L'Union européenne autorise les protéines animales transformées (PAT) d'insectes en \textbf{aquaculture depuis 2017} (Règlement (UE) 2017/893) et pour la \textbf{volaille et le porc depuis 2021}. Sept espèces sont autorisées, dont \emph{Hermetia illucens}, \emph{Tenebrio molitor} et \emph{Acheta domesticus}. Au Cameroun et en Afrique centrale, des fermes pilotes de MSN se multiplient (Yaoundé, Buea, Dschang), souvent adossées à des brasseries ou à des marchés : l'entomoculture devient un \textbf{métier émergent} pour les jeunes diplômés en productions animales.
\end{savaistu}

\begin{aretenir}[L'essentiel de la section]
\begin{itemize}
\item La collecte traditionnelle (piégeage lumineux, écorçage, cueillette) est saisonnière, non contrôlée et risquée (contaminants) : elle ne peut suffire à l'alimentation animale.
\item L'entomoculture contrôlée repose sur trois espèces clés : la \textbf{MSN} (bioconversion des déchets, cycle 14--18 jours, 40--45 \% de protéines MS), le \textbf{grillon} (6--8 semaines) et le \textbf{ver de farine} (\SI{2}{kg} de substrat $\rightarrow$ \SI{1}{kg} de larves).
\item Les coproduits (frass = engrais organique) et les peptides antimicrobiens des larves renforcent l'intérêt de la filière, à condition d'un \textbf{substrat tracé}.
\item Les farines d'insectes substituent efficacement la farine de poisson, avec des performances zootechniques équivalentes et un coût et une empreinte carbone réduits.
\end{itemize}
\end{aretenir}

\begin{exercice}[Pour s'entraîner]
Un jeune entrepreneur de Douala récupère chaque semaine \SI{500}{kg} de déchets de marché. Le taux de réduction des déchets par les larves de MSN est de 60 \% et le rendement en larves fraîches de \SI{150}{g} par kg de déchets. 1) Calcule la masse de larves produite par semaine. 2) Sachant que les larves fraîches contiennent 35 \% de matière sèche, calcule la masse de matière sèche larvaire hebdomadaire. 3) Que devient la fraction des déchets non convertie en larves ?
\end{exercice}

\begin{corrige}[Exercice 1]
1) $500 \times 0{,}150 = \SI{75}{kg}$ de larves fraîches par semaine. 2) $75 \times 0{,}35 = \SI{26,25}{kg}$ de matière sèche (soit environ \SIrange{10,5}{11,8}{kg} de protéines à 40--45 \% de la MS). 3) Le taux de réduction de 60 \% signifie que \SI{300}{kg} de déchets ont été réduits (convertis en biomasse larvaire, CO$_2$, eau), le reste (\SI{200}{kg}) constituant le \textbf{frass} récoltable (engrais organique), auquel s'ajoute un lixiviat à gérer.
\end{corrige}

\coursec{Communication, sensibilisation et acceptabilité sociale}

Dans la section précédente, nous avons vu comment passer de la simple collecte saisonnière des termites ailés à une véritable \cle{entomoculture} maîtrisée, capable de produire des tonnes de larves de mouche soldat noire (MSN) ou de grillons pour l'alimentation animale et humaine. Mais produire ne suffit pas : une protéine qui n'est pas consommée ne résout aucune insuffisance alimentaire. Toute la question devient alors : comment faire \emph{accepter} l'insecte dans l'assiette camerounaise de demain ? C'est l'objet de cette dernière section : comprendre les freins psychologiques et culturels, construire des messages efficaces, garantir la sécurité sanitaire, et inscrire l'entomophagie dans une économie circulaire vertueuse.

\courssub{Les barrières à l'entomophagie : pourquoi dit-on « non » aux insectes ?}

\textbf{Une observation de terrain.} Dans un lycée de Yaoundé, on propose à 200 élèves des beignets de termites grillés (mets traditionnel très apprécié en saison des pluies dans l'Ouest et le Centre, appelés \emph{mingwel}) et des barres protéinées à la farine de grillon. Résultat : les beignets de termites partent vite chez les élèves qui en mangent au village, mais une majorité d'élèves urbains refuse de « manger des bestioles », alors que les mêmes élèves avalent sans hésiter les barres protéinées... dont ils ignorent la composition. Que nous apprend cette scène ? Que le refus n'est pas \emph{gustatif} mais \emph{psychologique et culturel}.

\begin{definition}[Néophobie alimentaire]
La \cle{néophobie alimentaire} est la peur ou le dégoût instinctif à l'égard d'un aliment nouveau ou perçu comme inhabituel. C'est un comportement adaptatif hérité de l'évolution (se méfier de l'inconnu protège de l'empoisonnement), mais qui devient un frein lorsqu'il s'oppose à des aliments sûrs et nutritifs comme les insectes comestibles. On parle parfois, par abus de langage, d'\textbf{entomophobie} pour désigner le dégoût culturel spécifique des insectes (ne pas confondre avec « néphobie » qui désigne la peur des nuages ou du rein).
\end{definition}

On distingue quatre barrières principales, qu'il faut connaître pour pouvoir les lever :

\begin{itemize}
\item \textbf{Le dégoût culturel} : dans les sociétés urbanisées, l'insecte est associé à la saleté, à la maladie (moustique, cafard) ou à la nuisance, alors que les insectes \emph{comestibles} (termite, chenille, grillon) n'ont rien à voir avec ces espèces.
\item \textbf{La néophobie alimentaire} : peur de l'aliment nouveau, surtout marquée chez l'adulte ; les enfants et adolescents y sont beaucoup moins sensibles, d'où l'intérêt de cibler les jeunes.
\item \textbf{La méconnaissance nutritionnelle} : beaucoup ignorent que \SI{100}{\gram} de grillons séchés (poids sec) apportent environ \SI{60}{\gram} de protéines, soit davantage que la viande de bœuf (environ \SI{20}{\gram} pour \SI{100}{\gram} poids frais), avec du fer et du zinc hautement biodisponibles.
\item \textbf{Le statut social dévalorisé} : l'insecte est perçu comme une « nourriture de pauvre » ou de brousse, alors que c'est un mets de fête dans de nombreuses cultures camerounaises (termite ailé, chenille \emph{minkoul}).
\end{itemize}

\begin{propriete}[Effet de la forme de présentation sur l'acceptabilité]
En psychologie alimentaire, il est établi que l'\cle{acceptabilité} d'un aliment à base d'insectes augmente significativement lorsque l'insecte est \textbf{invisible} (incorporé sous forme de farine dans du pain, des pâtes, des beignets ou des barres protéinées) plutôt que présenté \textbf{entier} (grillé, visible). La transformation technologique est donc un levier majeur d'adoption : on passe du « je vois la bestiole » au « je mange un produit bon et nutritif ».
\end{propriete}

\courssub{Stratégies de communication : que dire, à qui, et comment ?}

Deux grandes stratégies de communication s'opposent, et l'expérience des campagnes internationales montre qu'elles n'ont pas la même efficacité :

\begin{itemize}
\item \textbf{L'approche « Food First »} : on met en avant le \emph{goût}, la gastronomie, la recette, le plaisir. Message type : « Goûte ce beignet de mingwel, il est délicieux ! » C'est l'approche la plus efficace pour convertir durablement, car le plaisir gustatif crée l'habitude.
\item \textbf{L'approche « Health/Planet »} : on met en avant la \emph{nutrition} (protéines, fer, zinc) et l'\emph{écologie} (faible émission de gaz à effet de serre, faible consommation d'eau). Elle convainc les décideurs, les éleveurs et les consommateurs déjà sensibilisés, mais ne suffit pas à faire goûter un réfractaire.
\end{itemize}

La bonne stratégie consiste à \emph{combiner} les deux et à adapter le message à la cible : les \textbf{jeunes urbains} (réseaux sociaux, recettes tendance), les \textbf{décideurs} (chiffres, sécurité alimentaire nationale), les \textbf{éleveurs} (rentabilité, disponibilité des protéines pour la volaille et les porcs).

\begin{methode}[Concevoir un message de sensibilisation efficace]
Pour construire un message qui change réellement les comportements, suivre quatre étapes :
\begin{enumerate}
\item \textbf{Définir la cible} : qui veut-on convaincre ? (ex. : élèves de lycée urbain, femmes enceintes, éleveurs de volaille de la périphérie de Douala).
\item \textbf{Formuler le message clé sous forme de bénéfice} : pas « mangez des insectes », mais « des protéines de qualité à petit prix pour toute la famille ».
\item \textbf{Apporter une preuve chiffrée et locale} : « \SI{100}{\gram} de grillons séchés = autant de protéines que \SI{300}{\gram} de bœuf, pour trois fois moins cher » ; « élevé ici, à Yaoundé ».
\item \textbf{Conclure par un appel à l'action simple} : goûter (dégustation), élever (kit grillons), ou partager (recette sur les réseaux sociaux).
\end{enumerate}
\end{methode}

\courssub{Conception d'outils pédagogiques de sensibilisation}

Le savoir-faire exigé au programme est de \emph{concevoir} des outils. Voici quatre exemples directement réalisables par un élève de Terminale :

\begin{exemplebox}[Quatre outils de sensibilisation]
\begin{itemize}
\item \textbf{Affiche « Assiette protéique camerounaise 2035 »} : une assiette divisée montrant bœuf, poisson, poulet \emph{et} grillons/MSN côte à côte, avec une infographie comparative (eau, CO$_2$, prix par kg de protéine) et un code QR renvoyant vers des recettes.
\item \textbf{Dépliant « Élever ses grillons à la maison »} : recto = matériel (bac, son, chaleur, eau) ; verso = cycle de vie du grillon, calendrier de récolte (6 semaines), recette de base.
\item \textbf{Vidéo courte (format TikTok/Reels) « Recette mingwel sauce arachide »} : 60 secondes, plan serré sur la sauce, l'insecte apparaît déjà cuisiné (forme invisible), commentaire joyeux, hashtags locaux.
\item \textbf{Jeu de cartes « Insectes vs Viandes »} : chaque carte = un aliment avec ses valeurs (protéines, fer, prix, eau consommée) ; on joue à la bataille, et le grillon gagne souvent : le message passe par le jeu.
\end{itemize}
\end{exemplebox}

\begin{popfigure}
\begin{center}
\begin{tikzpicture}[scale=1]
% Cadre affiche
\fill[popcream] (0,0) rectangle (12,8.4);
\draw[popdark,line width=1.2pt] (0,0) rectangle (12,8.4);
% Bandeau titre
\fill[popgreen] (0,7.2) rectangle (12,8.4);
\node[white,font=\bfseries\large] at (6,7.8) {PROT\'EINES POUR TOUS --- L'ASSIETTE CAMEROUNAISE 2035};
% Colonnes comparatives
\foreach \x/\nom/\col in {1.6/{B\oe uf}/popdark, 4.1/{Poisson}/popblue, 6.6/{Poulet}/poporange, 9.1/{Grillon}/popgreen} {
  \fill[\col!25] (\x-1.05,4.4) rectangle (\x+1.05,6.9);
  \draw[\col,line width=0.8pt] (\x-1.05,4.4) rectangle (\x+1.05,6.9);
  \node[font=\bfseries\small,text=\col] at (\x,6.6) {\nom};
}
% Barres eau (litres/kg proteine, echelle relative)
\node[font=\footnotesize\bfseries,text=popink] at (6,4.15) {Eau consomm\'ee (L/kg de prot\'eine)};
\fill[popdark]  (0.9,3.0) rectangle (2.3,3.8);  \node[font=\tiny] at (1.6,2.8) {$\approx$ 15\,000};
\fill[popblue]  (3.4,3.25) rectangle (4.8,3.8); \node[font=\tiny] at (4.1,3.05) {$\approx$ 5\,000};
\fill[poporange](5.9,3.35) rectangle (7.3,3.8); \node[font=\tiny] at (6.6,3.15) {$\approx$ 3\,500};
\fill[popgreen] (8.4,3.65) rectangle (9.8,3.8); \node[font=\tiny] at (9.1,3.45) {$\approx$ 1\,000};
% Prix
\node[font=\footnotesize\bfseries,text=popink] at (6,2.35) {Prix du kg de prot\'eine (FCFA, indicatif)};
\fill[popdark]  (0.9,1.3) rectangle (2.3,2.0);  \node[font=\tiny] at (1.6,1.1) {cher};
\fill[popblue]  (3.4,1.45) rectangle (4.8,2.0); \node[font=\tiny] at (4.1,1.25) {moyen};
\fill[poporange](5.9,1.55) rectangle (7.3,2.0); \node[font=\tiny] at (6.6,1.35) {moyen};
\fill[popgreen] (8.4,1.8) rectangle (9.8,2.0);  \node[font=\tiny] at (9.1,1.6) {bas};
% QR code fictif
\fill[white] (10.5,0.4) rectangle (11.8,1.7);
\draw[popdark] (10.5,0.4) rectangle (11.8,1.7);
\foreach \i in {0,...,5} \foreach \j in {0,...,5} {
  \pgfmathparse{mod(\i*7+\j*3+2,3)==0 ? 1 : 0}
  \ifnum\pgfmathresult=1 \fill[popdark] (10.55+\i*0.21,0.45+\j*0.21) rectangle (10.72+\i*0.21,0.62+\j*0.21);\fi
}
\node[font=\tiny,text=popink,align=center] at (11.15,0.15) {Scanne : recettes};
% Slogan bas
\node[font=\small\itshape,text=popink] at (5.6,0.7) {Grillons et MSN : riches en fer et zinc, produits localement, bons pour la sant\'e et la plan\`ete !};
\end{tikzpicture}
\end{center}
\legende{Maquette d'affiche de sensibilisation « Protéines pour tous ». L'infographie compare quatre sources de protéines animales selon deux critères objectifs (eau consommée, prix du kg de protéine) ; le grillon apparaît systématiquement avantageux. Le code QR renvoie vers des recettes : c'est l'appel à l'action.}
\end{popfigure}

\courssub{Mesurer l'acceptabilité : une enquête comme outil d'investigation}

Avant de lancer une campagne, il faut \emph{mesurer} l'état des opinions : c'est une démarche d'investigation. On interroge par questionnaire un échantillon de 200 étudiants de l'Université de Yaoundé I à la question « Seriez-vous prêt à manger des insectes ? ».

\begin{popfigure}
\begin{center}
\begin{tikzpicture}[scale=1.1]
% Diagramme circulaire manuel
\fill[popgreen]  (0,0) -- (0:2.2) arc (0:126:2.2) -- cycle;      % Oui 35%
\fill[popblue]   (0,0) -- (126:2.2) arc (126:241.2:2.2) -- cycle; % Si transformé 32%
\fill[poporange] (0,0) -- (241.2:2.2) arc (241.2:313.2:2.2) -- cycle; % Si invisible 20%
\fill[poppink]   (0,0) -- (313.2:2.2) arc (313.2:360:2.2) -- cycle;   % Non 13%
\draw[white,line width=1.5pt] (0:2.2) -- (0,0) -- (126:2.2) (126:2.2) -- (0,0) -- (241.2:2.2) (241.2:2.2) -- (0,0) -- (313.2:2.2) (313.2:2.2) -- (0,0) -- (360:2.2);
\draw[popdark] (0,0) circle (2.2);
% Etiquettes
\node[font=\footnotesize\bfseries,white] at (63:1.4) {Oui};
\node[font=\footnotesize,white] at (63:1.0) {35 \%};
\node[font=\footnotesize\bfseries,white] at (183:1.4) {Si transform\'e};
\node[font=\footnotesize,white] at (183:1.0) {32 \%};
\node[font=\footnotesize\bfseries,white] at (277:1.4) {Si invisible};
\node[font=\footnotesize,white] at (277:1.0) {20 \%};
\node[font=\footnotesize\bfseries,white] at (336:1.35) {Non};
\node[font=\footnotesize,white] at (336:1.0) {13 \%};
% Legende raisons
\node[anchor=west,font=\footnotesize,text=popink,align=left] at (2.8,1.2) {\textbf{Raisons du refus (13 \,\%)} :};
\node[anchor=west,font=\footnotesize,text=popink,align=left] at (2.8,0.6) {-- d\'ego\^ut (60 \,\% des refus)};
\node[anchor=west,font=\footnotesize,text=popink,align=left] at (2.8,0.1) {-- peur sanitaire (25 \,\%)};
\node[anchor=west,font=\footnotesize,text=popink,align=left] at (2.8,-0.4) {-- habitude / image sociale (15 \,\%)};
\node[anchor=west,font=\footnotesize\itshape,text=popmuted] at (2.8,-1.2) {Enqu\^ete fictive, $n = 200$ \'etudiants, Yaound\'e};
\end{tikzpicture}
\end{center}
\legende{Diagramme circulaire d'une enquête d'acceptabilité. Lecture : 87 \% des répondants acceptent l'insecte sous au moins une forme, et 52 \% l'acceptent \emph{uniquement} s'il est transformé ou invisible — ce qui illustre la propriété de psychologie alimentaire vue plus haut et justifie l'investissement dans les farines d'insectes.}
\end{popfigure}

\begin{exoresolu}[Exploiter les résultats d'une enquête]
\textbf{Énoncé.} Sur les 200 étudiants interrogés, 35 \% acceptent l'insecte entier, 32 \% seulement transformé, 20 \% seulement invisible, 13 \% refusent. Une entreprise veut lancer un seul produit touchant le maximum de consommateurs. Que conseilles-tu, et combien d'étudiants de l'échantillon seraient concernés ?
\tcblower
\textbf{Solution.} Le produit qui maximise l'acceptabilité est celui où l'insecte est \emph{invisible} (farine incorporée) : il convient à ceux qui acceptent l'insecte entier (qui n'ont aucune raison de refuser la farine), à ceux qui l'acceptent transformé et à ceux qui l'acceptent invisible. Soit :
\[ 35\,\% + 32\,\% + 20\,\% = 87\,\% \quad\Longrightarrow\quad 0{,}87 \times 200 = 174 \text{ \'etudiants.} \]
Un produit « insecte entier grillé » ne toucherait que 70 étudiants. Conclusion : lancer d'abord une farine ou une barre protéinée, puis proposer les produits entiers aux convaincus.
\end{exoresolu}

\courssub{Cadre réglementaire et sécurité sanitaire : la condition de la confiance}

Aucune sensibilisation ne peut réussir si le produit n'est pas \emph{sûr}. L'entomoculture doit donc appliquer la méthode \cle{HACCP} (\emph{Hazard Analysis Critical Control Points}) : identifier les dangers (biologiques : parasites, salmonelles ; chimiques : pesticides, métaux lourds accumulés par les larves ; allergènes), définir les points critiques (qualité du substrat d'élevage, température de cuisson/séchage, stockage) et fixer des seuils de surveillance.

\begin{definition}[Allergène croisé : la tropomyosine]
La \cle{tropomyosine} est une protéine musculaire très conservée au cours de l'évolution, présente chez les crustacés, les acariens \emph{et} les insectes. Chez une personne allergique aux crustacés, les anticorps IgE dirigés contre la tropomyosine des crevettes reconnaissent aussi celle des insectes : c'est une \textbf{réactivité croisée}. Conséquence réglementaire : l'étiquetage doit mentionner « Contient des insectes (allergène : réactivité croisée possible avec les crustacés et les acariens) ».
\end{definition}

\begin{exemplebox}[Un exemple chiffré d'étiquetage]
Le seuil de détection de la tropomyosine par les méthodes immunologiques (ELISA) se situe entre \SI{1}{\milli\gram\per\kilo\gram} et \SI{10}{\milli\gram\per\kilo\gram} de produit. En dessous, la mention « Peut contenir des traces d'insectes » est utilisée en cas de contamination croisée accidentelle ; au-dessus, la mention « Contient des insectes » est obligatoire, au même titre que « Contient des crustacés ».
\end{exemplebox}

\begin{definition}[Traçabilité « de la fourche à la fourchette »]
La \cle{traçabilité} est la capacité à suivre un produit alimentaire à toutes les étapes, de la production (la « fourche » du producteur) jusqu'au consommateur (la « fourchette ») : origine du substrat d'élevage, lot, date de récolte, traitement thermique, conditions de transport. En cas de problème sanitaire, elle permet de retirer rapidement le lot incriminé et protège ainsi la confiance des consommateurs.
\end{definition}

\begin{attention}[Ne jamais promouvoir sans garantir la sécurité]
Erreur fréquente et grave : faire la promotion de la consommation d'insectes \emph{sans} garantir la sécurité sanitaire. Des insectes collectés dans la nature peuvent véhiculer des parasites, avoir accumulé des pesticides (champs traités) ou des métaux lourds (zones industrielles), et provoquer des allergies non signalées. Un seul incident médiatisé détruit durablement la confiance du public. Règle d'or : \textbf{la sécurité sanitaire précède toujours la communication}.
\end{attention}

\courssub{Économie circulaire : l'insecte au cœur de la ferme intégrée}

\begin{definition}[Économie circulaire]
L'\cle{économie circulaire} est un modèle de production dans lequel les déchets d'une activité deviennent les ressources d'une autre, en boucle fermée, au lieu d'être rejetés. Elle s'oppose au modèle linéaire « produire -- consommer -- jeter ».
\end{definition}

L'entomoculture s'intègre parfaitement dans une ferme circulaire camerounaise : les déchets de maraîchage (fanions, fruits invendus) nourrissent les larves de MSN ; les larves nourrissent les volailles et les poissons ; le \cle{frass} (déjections et restes de substrat des larves) fertilise le potager ; le potager produit à nouveau des déchets. Rien ne se perd.

\begin{savaistu}[Le frass, un engrais biostimulant]
Le frass de mouche soldat noire n'est pas un simple engrais : riche en \textbf{chitine} (exosquelettes) et en micro-organismes bénéfiques, il stimule les défenses naturelles des plantes et le microbiote du sol. Des essais maraîchers rapportent des rendements de tomate supérieurs de 20 à 30 \% par rapport à un engrais chimique NPK apporté à dose équivalente. Pour convaincre un éleveur-maraîcher, c'est un argument économique décisif : l'élevage d'insectes produit à la fois la protéine (pour la volaille) et l'engrais (pour le champ).
\end{savaistu}

\begin{exercice}[Concevoir ton outil de sensibilisation]
Ta commune veut lancer la consommation de farine de grillon dans les cantines scolaires. En appliquant la méthode en quatre étapes, rédige le message de sensibilisation complet (cible, message clé, preuve chiffrée locale, appel à l'action), puis indique les deux mentions de sécurité sanitaire qui doivent figurer sur l'emballage.
\end{exercice}

\begin{corrige}[Exercice]
\textbf{Cible} : parents d'élèves et cantinières. \textbf{Message clé (bénéfice)} : « Une farine locale qui donne à vos enfants autant de protéines et de fer que la viande, pour un prix abordable ». \textbf{Preuve} : « \SI{100}{\gram} de farine de grillon $\approx$ \SI{60}{\gram} de protéines, soit l'équivalent de \SI{300}{\gram} de bœuf ; produite ici, dans notre région ». \textbf{Appel à l'action} : « Venez goûter les beignets à la farine de grillon à la fête de l'école ». \textbf{Mentions sanitaires} : (1) « Contient des insectes (allergène : réactivité croisée possible avec les crustacés et les acariens) » ; (2) les informations de traçabilité (numéro de lot, date de production, origine de l'élevage).
\end{corrige}

\begin{aretenir}[L'essentiel de la section]
\begin{itemize}
\item Les freins à l'entomophagie sont psychologiques (néophobie, dégoût) et sociaux (image dévalorisée), non gustatifs ; l'insecte \emph{invisible} (farine) est bien mieux accepté que l'insecte entier.
\item Un message efficace suit la chaîne : cible $\rightarrow$ bénéfice $\rightarrow$ preuve chiffrée locale $\rightarrow$ appel à l'action ; l'approche « Food First » (goût) convertit mieux que les seuls arguments santé/planète.
\item La sécurité sanitaire (HACCP, traçabilité, étiquetage de l'allergène tropomyosine) est la condition préalable de toute communication.
\item L'économie circulaire (déchets $\rightarrow$ insectes $\rightarrow$ volaille + frass $\rightarrow$ maraîchage) fait de l'entomoculture un pilier de l'amélioration durable de la production animale.
\end{itemize}
\end{aretenir}

\newpage

\coursec{Activité d'intégration}

\courssub{Objectifs de l'activité}

À l'issue de cette activité d'intégration, tu seras capable de :

\begin{itemize}
\item mobiliser les savoirs de la leçon (valeur nutritionnelle des insectes, entomoculture, lutte biologique, produits issus des insectes) pour résoudre un problème authentique de production animale ;
\item exploiter des documents (tableaux de composition, courbes de croissance, fiches techniques, témoignages) pour extraire, organiser et traiter des données chiffrées ;
\item effectuer des calculs de dimensionnement d'un élevage d'insectes (besoins protéiques, surface, rendement, bilan matière) ;
\item argumenter l'intérêt zootechnique, sanitaire et environnemental de la mouche soldat noire (\emph{Hermetia illucens}, MSN) ;
\item concevoir un outil de sensibilisation (affiche ou script vidéo) destiné à un public non scientifique, en respectant les critères d'une communication efficace ;
\item développer des savoir-être : travail en équipe, esprit critique, sens de l'initiative, responsabilité environnementale.
\end{itemize}

\courssub{Situation}

\textbf{Projet intégré : « Ferme intégrée durable — De la poubelle à l'assiette ».}\par

À Dschang (région de l'Ouest, Cameroun), M.~Kamga exploite une ferme avicole de \num{5000} poulets de chair. Il fait face à trois problèmes majeurs :

\begin{enumerate}
\item le \textbf{coût de l'aliment} : la farine de poisson et le tourteau de soja, sources de protéines de la ration, sont importés et leur prix ne cesse d'augmenter (\SI{900}{FCFA/kg} pour la farine de poisson contre \SI{350}{FCFA/kg} pour le maïs) ;
\item la \textbf{gestion des fientes} : la ferme produit environ \SI{250}{kg} de fientes fraîches par jour, stockées à l'air libre, source d'odeurs, de pollution des sols et de contamination des eaux de puits ;
\item la \textbf{pullulation des mouches domestiques} (\emph{Musca domestica}) qui se reproduisent dans les fientes et constituent une nuisance ainsi qu'un vecteur mécanique d'agents pathogènes (salmonelles, \emph{Escherichia coli}).
\end{enumerate}

Un technicien de la Délégation régionale de l'élevage lui propose d'installer une unité d'élevage de la \cle{mouche soldat noire} (MSN), dont les larves se développent sur les fientes, en réduisent la masse, empêchent le développement des mouches domestiques et fournissent une \cle{farine de larves} riche en protéines capable de remplacer partiellement la farine de poisson. Le résidu de l'élevage, le \cle{frass}, est un excellent engrais organique. Précision utile : la MSN adulte ne se nourrit pas, ne pique pas et ne fréquente pas les habitations ; ce n'est donc pas un insecte nuisible, contrairement à la mouche domestique.

\textbf{Documents fournis au groupe :}

\textbf{Document 1 — Composition moyenne des matières premières (en \% de la matière sèche) :}

\begin{center}
\begin{tabularx}{\linewidth}{|l|X|X|X|}
\hline
\rowcolor{popblueL} \textbf{Composant} & \textbf{Farine de poisson} & \textbf{Tourteau de soja} & \textbf{Farine de larves MSN} \\
\hline
Protéines brutes (\%) & 60 & 44 & 42 \\
\hline
Lipides (\%) & 8 & 2 & 30 \\
\hline
Calcium (\%) & 5 & 0,3 & 6 \\
\hline
Énergie métabolisable (kcal/kg) & 2900 & 2400 & 3900 \\
\hline
Prix (FCFA/kg) & 900 & 650 & 300 (coût de production estimé) \\
\hline
\end{tabularx}
\end{center}

\textbf{Document 2 — Croissance des larves de MSN} sur substrat fientes de volaille, comparée à un substrat témoin (son de blé) : la masse moyenne d'une larve augmente régulièrement puis plafonne vers le jour 14, moment optimal de la récolte (avant la nymphose). On retient : masse fraîche d'une larve à récolte (jour 14) : \SI{0,20}{g} ; taux de survie : 90 \% ; densité d'élevage recommandée : 5 larves/cm$^2$ de surface de substrat.

\textbf{Document 3 — Fiche technique d'élevage de la MSN :}
\begin{itemize}
\item cycle œuf $\rightarrow$ larve à récolte : 14 jours ;
\item 1 kg de larves fraîches produit consomme environ 2 kg de fientes fraîches et la masse du substrat est réduite de 50 \% ;
\item le frass récolté représente environ 40 \% de la masse de substrat mis en œuvre ;
\item 1 kg de larves fraîches donne environ \SI{0,30}{kg} de farine séchée (taux de matière sèche des larves : 30 \%).
\end{itemize}

\textbf{Document 4 — Témoignages d'éleveurs de la région de Bafoussam} ayant adopté la technologie : baisse de 25 \% du coût alimentaire, disparition quasi totale des mouches domestiques, vente du frass aux maraîchers à \SI{100}{FCFA/kg}.

\textbf{Données zootechniques du lot :} un poulet de chair consomme en moyenne \SI{115}{g} d'aliment par jour ; la ration doit contenir 20 \% de protéines brutes ; la farine de poisson apporte actuellement 25 \% des protéines de la ration.

\courssub{Consignes}

Par groupes de 4 à 5 élèves, et à l'aide des documents :

\begin{enumerate}
\item \textbf{Tâche A — Dimensionnement.} Calculer le besoin journalier total en protéines du lot de \num{5000} poulets, puis la part de protéines apportée par la farine de poisson. En déduire la masse journalière de farine de larves MSN nécessaire pour substituer 30 \% de cette farine de poisson (à quantité de protéines équivalente), puis la masse de larves fraîches correspondante.
\item \textbf{Tâche A (suite).} En utilisant la densité d'élevage et la durée du cycle (14 jours), estimer la surface totale des bacs d'élevage nécessaire, sachant que la production est continue (un lot de larves démarre chaque jour).
\item \textbf{Tâche B — Bilan matière.} Estimer la quantité de fientes valorisée par jour et la comparer à la production journalière de la ferme (\SI{250}{kg}). Calculer la masse de frass produite par jour et le revenu mensuel potentiel de sa vente (30 jours).
\item \textbf{Tâche C — Argumentation sanitaire et environnementale.} À partir de tes connaissances sur la lutte biologique et de la fiche technique, rédiger un court texte (10 lignes maximum) expliquant en quoi l'élevage de MSN constitue une méthode de lutte biologique contre les mouches domestiques et un outil d'assainissement, en citant au moins deux avantages par rapport aux insecticides chimiques.
\item \textbf{Tâche D — Communication.} Produire UN outil de sensibilisation au choix : une affiche A4 (titre accrocheur, 3 messages clés chiffrés, un schéma simple du cycle « fientes $\rightarrow$ larves $\rightarrow$ poulet ») ou le script d'une vidéo de 60 secondes (environ 150 mots) destiné à convaincre M.~Kamga d'adopter la technologie.
\item \textbf{Restitution :} présentation orale de 10 minutes par groupe + remise des supports (feuille de calculs et outil de sensibilisation).
\end{enumerate}

\courssub{Ressources à mobiliser}

\begin{itemize}
\item \textbf{Valeur nutritionnelle des insectes :} les larves de MSN contiennent environ 42 \% de protéines brutes (matière sèche), riches en acides aminés essentiels (lysine, méthionine), comparables à la farine de poisson ; elles sont aussi riches en lipides énergétiques et en calcium, minéral indispensable à la solidité du squelette et à la qualité de la coquille des œufs.
\item \textbf{Lutte biologique :} utilisation d'un organisme vivant (ici la MSN) pour limiter une population d'organismes nuisibles (ici la mouche domestique) par \emph{compétition} pour la ressource (les fientes) et modification du milieu ; alternative aux insecticides chimiques, sans résidus ni résistance. C'est le même principe que l'emploi de la coccinelle contre les pucerons ou de la libellule contre les moustiques.
\item \textbf{Économie circulaire :} les déchets (fientes) deviennent une ressource (substrat d'élevage), les larves deviennent aliment, le résidu (frass) devient engrais.
\item \textbf{Calculs utiles :} proportionnalité, pourcentages, bilan matière :
\[ \text{masse de protéines} = \text{masse d'aliment} \times \text{taux de protéines} \; ; \qquad \text{surface} = \frac{\text{nombre de larves}}{\text{densité}}. \]
\end{itemize}

\begin{popfigure}
\begin{center}
\begin{tikzpicture}[node distance=1.1cm, every node/.style={font=\small}]
\node[draw=popblue, fill=popblueL, rounded corners, minimum width=3.2cm, minimum height=1cm, align=center] (fientes) {\textbf{Fientes}\\ \SI{250}{kg/jour} (déchet)};
\node[draw=poporange, fill=poporangeL, rounded corners, minimum width=3.2cm, minimum height=1cm, align=center, right=2.2cm of fientes] (larves) {\textbf{Larves de MSN}\\ \SI{68,4}{kg/jour}};
\node[draw=popgreen, fill=popgreenL, rounded corners, minimum width=3.2cm, minimum height=1cm, align=center, right=2.2cm of larves] (poulet) {\textbf{Poulets}\\ protéines de la ration};
\node[draw=poppurple, fill=poppurpleL, rounded corners, minimum width=3.2cm, minimum height=1cm, align=center, below=1.4cm of larves] (frass) {\textbf{Frass}\\ engrais organique};
\draw[-{Stealth}, very thick, popblue] (fientes) -- node[above]{substrat} (larves);
\draw[-{Stealth}, very thick, poporange] (larves) -- node[above]{farine 42 \% protéines} (poulet);
\draw[-{Stealth}, very thick, poppurple] (larves) -- node[right]{résidu 40 \%} (frass);
\draw[-{Stealth}, very thick, popgreen, dashed] (poulet.south) to[out=-120,in=0] node[below left=0.05cm and -0.6cm]{fientes} (fientes.south east);
\end{tikzpicture}
\end{center}
\legende{Schéma du cycle de l'économie circulaire dans la ferme intégrée : les fientes (déchet) nourrissent les larves de MSN, qui deviennent farine protéique pour les poulets ; le résidu (frass) est valorisé comme engrais, et les poulets produisent à leur tour les fientes qui alimentent le cycle.}
\end{popfigure}

\courssub{Démarche et corrigé commenté}

\begin{exoresolu}[Corrigé commenté du projet « Ferme intégrée durable »]
\textbf{Énoncé.} Substituer 30 \% de la farine de poisson d'un lot de \num{5000} poulets de chair par de la farine de larves de MSN produite sur les fientes de la ferme ; dimensionner l'unité d'élevage, établir le bilan matière et argumenter l'intérêt sanitaire.
\tcblower
\textbf{Étape 1 — Besoin journalier en protéines du lot.}

Consommation totale d'aliment par jour :
\[ \num{5000} \times \SI{115}{g} = \SI{575000}{g} = \SI{575}{kg/jour}. \]
Ration à 20 \% de protéines brutes :
\[ 575 \times 0{,}20 = \SI{115}{kg} \text{ de protéines par jour.} \]

\textbf{Étape 2 — Part de la farine de poisson et objectif de substitution.}

La farine de poisson apporte 25 \% des protéines :
\[ 115 \times 0{,}25 = \SI{28,75}{kg} \text{ de protéines/jour,} \]
soit une masse de farine de poisson de $28{,}75 / 0{,}60 \approx \SI{47,9}{kg/jour}$.

Substituer 30 \% de cet apport protéique revient à remplacer :
\[ 28{,}75 \times 0{,}30 = \SI{8,63}{kg} \text{ de protéines/jour.} \]

\textbf{Étape 3 — Masse de farine de larves MSN et de larves fraîches.}

La farine de larves contient 42 \% de protéines :
\[ \frac{8{,}63}{0{,}42} \approx \SI{20,5}{kg} \text{ de farine de larves par jour.} \]
Avec un rendement de \SI{0,30}{kg} de farine par kg de larves fraîches :
\[ \frac{20{,}5}{0{,}30} \approx \SI{68,4}{kg} \text{ de larves fraîches à produire chaque jour.} \]
\emph{Attention :} on raisonne ici en deux étapes (protéines $\rightarrow$ farine sèche $\rightarrow$ larves fraîches) ; confondre masse fraîche et matière sèche est l'erreur la plus fréquente.

\textbf{Étape 4 — Surface d'élevage (Tâche A).}

Nombre de larves à récolter par jour (masse unitaire \SI{0,20}{g}) :
\[ \frac{\num{68400}}{0{,}20} = \num{342000} \text{ larves récoltées par jour.} \]
Compte tenu du taux de survie de 90 \%, il faut mettre en élevage :
\[ \frac{\num{342000}}{0{,}90} = \num{380000} \text{ larves chaque jour.} \]
Surface journalière à 5 larves/cm$^2$ :
\[ \frac{\num{380000}}{5} = \num{76000} \text{ cm}^2 = \SI{7,6}{m^2} \text{ par lot journalier.} \]
Le cycle durant 14 jours et un nouveau lot démarrant chaque jour, 14 lots coexistent simultanément :
\[ 7{,}6 \times 14 \approx \SI{106}{m^2} \text{ de bacs d'élevage au total.} \]
\emph{Commentaire :} une centaine de mètres carrés, soit un bâtiment modeste d'environ $10 \times \SI{11}{m}$, suffit : la technologie est réaliste à l'échelle de la ferme.

\textbf{Étape 5 — Bilan matière (Tâche B).}

Fientes consommées : $68{,}4 \times 2 \approx \SI{137}{kg/jour}$, soit environ 55 \% des \SI{250}{kg} de fientes produites : plus de la moitié du déchet est valorisée.

Frass produit (40 \% du substrat mis en œuvre) :
\[ 137 \times 0{,}40 \approx \SI{54,8}{kg/jour}. \]
Revenu mensuel de la vente du frass à \SI{100}{FCFA/kg} :
\[ 54{,}8 \times 30 \times 100 \approx \SI{164400}{FCFA/mois.} \]
\emph{Gain complémentaire :} remplacer $8{,}63/0{,}60 \approx \SI{14,4}{kg}$ de farine de poisson (\SI{900}{FCFA/kg}) par \SI{20,5}{kg} de farine MSN (\SI{300}{FCFA/kg}) économise environ :
\[ 14{,}4 \times 900 - 20{,}5 \times 300 = \num{12960} - \num{6150} \approx \SI{6810}{FCFA/jour,} \]
soit plus de \SI{200000}{FCFA/mois} d'économie sur l'aliment, cohérent avec la baisse de 25 \% du coût alimentaire rapportée par les éleveurs de Bafoussam (document 4).

\textbf{Étape 6 — Argumentation sanitaire (Tâche C).}

Les larves de MSN colonisent massivement les fientes fraîches et en consomment rapidement la matière organique, avant que les mouches domestiques ne puissent y pondre et y achever leur développement : c'est une \cle{lutte biologique} par compétition pour la ressource et par modification du milieu (l'activité intense des larves assèche et réchauffe le substrat, le rendant impropre à la survie des œufs et larves de \emph{Musca domestica}). Avantages par rapport aux insecticides chimiques : (1) absence de résidus chimiques dans la viande, les fientes et l'environnement ; (2) pas de phénomène de résistance chez les mouches ; (3) coût de fonctionnement quasi nul une fois l'unité installée ; (4) double bénéfice : assainissement et production de protéines. La réduction de 50 \% de la masse des fientes limite en outre les odeurs, les eaux de ruissellement polluées et les risques de contamination des puits.

\textbf{Étape 7 — Outil de sensibilisation (Tâche D, exemple de script vidéo 60 s).}

« Monsieur Kamga, vos poulets mangent cher, vos fientes polluent, les mouches envahissent votre ferme. La solution vole déjà autour de vous : la mouche soldat noire ! Ses larves transforment chaque jour 137 kg de vos fientes en 68 kg de larves riches en protéines pour vos poulets — une farine presque aussi riche que la farine de poisson, mais trois fois moins chère. Résultat : plus de 200 000 FCFA économisés par mois sur l'aliment, environ 165 000 FCFA de revenu avec l'engrais frass, et les mouches domestiques disparaissent sans un gramme d'insecticide. Cent mètres carrés de bacs suffisent. De la poubelle à l'assiette : faites travailler les insectes pour votre ferme ! »
\end{exoresolu}

\courssub{Bilan de l'activité}

Cette activité a permis de mettre en évidence que :

\begin{itemize}
\item l'\cle{entomophagie} et l'\cle{entomoculture} ne sont pas des pratiques marginales mais une réponse concrète et chiffrable à l'insuffisance des ressources protéiques classiques : une centaine de mètres carrés de larves de MSN peut substituer une part significative des protéines importées d'un élevage industriel ;
\item les insectes assurent une \cle{lutte biologique} efficace : la MSN élimine les mouches domestiques par compétition, sans insecticide, illustrant le rôle des insectes auxiliaires (comme la coccinelle contre les pucerons ou la libellule contre les moustiques) ;
\item la démarche s'inscrit dans une \cle{économie circulaire} : déchet (fientes) $\rightarrow$ ressource (larves) $\rightarrow$ coproduit valorisable (frass), ce qui transforme un problème sanitaire en double source de revenus ;
\item le traitement de données expérimentales (tableaux de composition, courbes de croissance, bilans matière) et l'argumentation chiffrée sont des savoir-faire transversaux indispensables au futur technicien, agronome ou médecin ;
\item la réussite d'une innovation biotechnologique dépend autant de la \cle{communication} et de l'acceptabilité sociale que de la rigueur scientifique : savoir concevoir un message clair, chiffré et adapté à son public est une compétence à part entière.
\end{itemize}

\begin{attention}[Critères d'évaluation de la restitution et erreurs fréquentes]
Barème indicatif : exactitude des calculs et cohérence des unités (4 pts) ; exploitation correcte des documents (4 pts) ; qualité de l'argumentation scientifique sur la lutte biologique (4 pts) ; pertinence et clarté de l'outil de sensibilisation (4 pts) ; qualité de la présentation orale et du travail d'équipe (4 pts). Attention aux erreurs fréquentes : confusion entre masse fraîche et matière sèche ; oubli du taux de survie dans le calcul du nombre de larves mises en élevage ; surface calculée pour un seul lot au lieu des 14 lots simultanés ; division par le taux de protéines oubliée lors du passage « masse de protéines $\rightarrow$ masse de farine ».
\end{attention}

\newpage

\coursec{Méthodes et exercices résolus}

\coursec{Amélioration de la production animale par l'exploitation des insectes}

\courssub{Contexte : Insuffisance protéique et potentialité de l'entomophagie}

\begin{definition}[Entomophagie et entomoculture]
L'\cle{entomophagie} est la consommation d'insectes par l'homme (entomophagie humaine) ou par les animaux d'élevage (entomophagie animale via l'incorporation de farines d'insectes dans les aliments). L'\cle{entomoculture} (ou élevage d'insectes) désigne la production maîtrisée d'insectes dans un but alimentaire, thérapeutique ou agricole (lutte biologique).
\end{definition}

\begin{savaistu}[Le défi protéique au Cameroun]
Au Cameroun, la production animale classique (bovins, porcins, volailles) ne couvre pas les besoins croissants en protéines animales : le déficit annuel est estimé à plus de \SI{200 000}{tonnes} équivalent viande (données MINEPIA). L'élevage conventionnel demande de vastes surfaces, beaucoup d'eau (\SI{15 000}{L} pour \SI{1}{kg} de bœuf) et émet des gaz à effet de serre (méthane). L'entomophagie, pratique traditionnelle dans de nombreuses régions (Adamaoua, Est, Sud, Ouest), apparaît comme une alternative durable : cycle court, fort taux de conversion alimentaire, valorisation de déchets organiques, faible empreinte écologique.
\end{savaistu}

\begin{aretenir}[Données clés pour le Bac]
\begin{itemize}
\item Besoin protéique moyen adulte : \SI{0,8}{g.kg^{-1}.j^{-1}} (OMS) ; adolescent en croissance : \SI{1,2}{g.kg^{-1}.j^{-1}}.
\item Teneurs en protéines (\% Matter Sèche - MS) : Termites ailés \SI{\sim 40}{\percent}, Chenilles (\emph{Imbrasia} spp.) \SI{\sim 50}{\percent}, Vers blancs (\emph{Rhynchophorus phoenicis}) \SI{\sim 45}{\percent}, Farine de larves de mouche soldat noire \SI{\sim 45-50}{\percent}.
\item Viande de bœuf fraîche : \SI{\sim 20}{\percent} de protéines (\SI{\sim 75}{\percent} en MS).
\item Miel : \SI{\sim 80}{\percent} glucides, \SI{\sim 17}{\percent} eau, valeur énergétique \SI{\sim 3200}{kcal.kg^{-1}} (\SI{\sim 13,4}{MJ.kg^{-1}}).
\end{itemize}
\end{aretenir}

\courssub{Valeur nutritionnelle et thérapeutique des insectes comestibles locaux}

\begin{methode}[Démarche pour décrire et justifier une technique d'utilisation d'un insecte comestible ou thérapeutique]
\begin{enumerate}
\item \textbf{Identifier} précisément l'insecte : nom vernaculaire, nom scientifique (genre), stade collecté (adulte ailé, larve, nymphe) et groupe taxonomique (Isoptères, Coléoptères, Lépidoptères, Hyménoptères).
\item \textbf{Décrire la technique} en étapes chronologiques : période de collecte (saison, nycthémère), matériel, mode de capture (piège lumineux, excavation, cueillette manuelle, enfumage), puis préparation culinaire (blanchiment, grillage, friture, séchage solaire, mouture en farine).
\item \textbf{Justifier scientifiquement} chaque étape par la biologie de l'insecte : l'essaimage des termites est déclenché par la première pluie ramollissant le sol et la baisse de luminosité ; le séchage abaisse l'activité de l'eau ($a_w < 0,6$) stoppant le développement microbien ; la cuisson dénature les protéines (améliore la digestibilité) et détruit les parasites (kystes de ténias, bactéries).
\item \textbf{Conclure} par l'apport nutritionnel ou thérapeutique chiffré : teneur en protéines, lipides, acides aminés essentiels (lysine, méthionine), fer hémique, zinc, vitamines B ; ou molécule bioactive (antimicrobienne, anti-inflammatoire).
\end{enumerate}
\textbf{Réflexes de rédaction :} employer le vocabulaire exact (\cle{essaimage}, \cle{matière sèche}, \cle{activité de l'eau}), chiffrer les compositions (ex. : termites \SI{40}{\percent} prot. MS), relier à la famille de situations (insuffisance ressources classiques).
\end{methode}

\begin{exoresolu}[Technique de collecte des termites ailés et intérêt nutritionnel]
\textbf{Énoncé.} Dans la région de l'Adamaoua, les populations collectent les termites ailés au début de la saison des pluies. Décris la technique de collecte et de préparation de ces insectes, puis justifie leur intérêt dans la couverture des besoins en protéines animales.
\tcblower
\textbf{Solution détaillée.}
\begin{enumerate}
\item \textbf{Identification.} Il s'agit des termites ailés (imagos sexués de \emph{Macrotermes bellicosus} ou \emph{M. subhyalinus}), émergeant des termitières lors de l'essaimage nuptial.
\item \textbf{Technique de collecte.} L'essaimage a lieu au crépuscule (18h--20h), après les premières pluies abondantes : l'humidité ramollit la croûte de la termitière et permet la sortie des imagos. Les collecteurs creusent une cavité à la base de la termitière (méthode traditionnelle) ou placent un bassin d'eau sous une lampe tempête / lampe torche (piège lumineux) : attirés par la lumière (phototactisme positif), les termites tombent dans l'eau, perdent leurs ailes (dealation) et sont récupérés au petit matin. La collecte est saisonnière (début saison des pluies), nocturne et peu coûteuse.
\item \textbf{Préparation.} Les termites sont triés, lavés, puis grillés à la poêle sans huile, frits ou séchés au soleil sur des nattes. La cuisson (grillage/friture) détruit les microorganismes contaminants (salmonelles, \emph{E. coli}) et les éventuels parasites, améliore la saveur (réaction de Maillard) et la digestibilité des protéines ; le séchage solaire réduit la teneur en eau de \SI{\sim 70}{\percent} à \SI{< 10}{\percent}, assurant une conservation de plusieurs mois à l'abri de l'humidité.
\item \textbf{Justification nutritionnelle.} En matière sèche, les termites ailés apportent \SI{\sim 40}{\percent} de protéines (avec un bon profil en acides aminés essentiels, score chimique \SI{> 80}{\percent}), \SI{\sim 40}{\percent} de lipides (riches en acides gras insaturés), \SI{\sim 35}{mg.100g^{-1}} de fer (forme hémique, bien absorbée) et \SI{\sim 15}{mg.100g^{-1}} de zinc. À titre de comparaison, la viande de bœuf fraîche n'apporte que \SI{20}{\percent} de protéines (base fraîche) et \SI{2,5}{mg.100g^{-1}} de fer. \SI{100}{g} de termites séchés couvrent \SI{> 100}{\percent} des AJR en fer et zinc d'un enfant.
\end{enumerate}
\textbf{Conclusion.} La collecte des termites ailés, technique simple calée sur le rythme biologique de l'insecte, fournit une source de protéines animales concentrée, riche en micronutriments critiques (fer, zinc), accessible aux ménages ruraux et complémentaire aux ressources classiques insuffisantes.
\end{exoresolu}

\begin{exemplebox}[Autres insectes comestibles majeurs du Cameroun]
\begin{center}
\begin{tabularx}{\linewidth}{|l|X|X|X|}\hline
\rowcolor{popblueL}
\textbf{Insecte (Nom local)} & \textbf{Groupe / Espèce} & \textbf{Stade consommé} & \textbf{Particularités nutritionnelles / thérapeutiques} \\ \hline
\textbf{Vers blancs} (Mpose, Mbinzo) & Coléoptères : \emph{Rhynchophorus phoenicis} (charançon du palmier) & Larve (dernier stade) & \SI{\sim 45}{\percent} prot. MS, \SI{\sim 30}{\percent} lipides ; très riche en zinc et vit. B12. Utilisés en pâtée pour enfants dénutris. \\ \hline
\textbf{Chenilles} (Mangou, Nkui) & Lépidoptères : \emph{Imbrasia oyemensis}, \emph{Cirina forda} & Larve (dernier stade) & \SI{\sim 50-55}{\percent} prot. MS, riches en fer, calcium, phosphore. Séchées, se conservent 1 an. \\ \hline
\textbf{Fourmis} (Ngong, Kèsèk) & Hyménoptères : \emph{Carebara} spp. (fourmis de terre), reines ailées de \emph{Dorylus} spp. & Reines ailées, ouvrières & Reines : \SI{\sim 40}{\percent} prot. MS, \SI{\sim 45}{\percent} lipides. Vertus aphrodisiaques et toniques (médecine traditionnelle). \\ \hline
\textbf{Abeilles} (larves/nymphes) & Hyménoptères : \emph{Apis mellifera adansonii} & Couvain (larves, nymphes) & Sous-produit de l'apiculture : \SI{\sim 40}{\percent} prot. MS, gelée royale endogène. Consommées grillées ou en sauce. \\ \hline
\end{tabularx}
\end{center}
\legende{Principaux insectes comestibles camerounais et leurs atouts nutritionnels (valeurs moyennes en matière sèche).}
\end{exemplebox}

\courssub{Produits de la ruche et substances bioactives : Miel, Venins, Propolis}

\begin{definition}[Apithérapie]
L'\cle{apithérapie} est l'utilisation médicale des produits de la ruche (miel, venin, propolis, gelée royale, pollen, cire) pour prévenir ou soigner des pathologies humaines ou animales.
\end{definition}

\begin{methode}[Démarche pour présenter un produit de la ruche et son utilisation]
\begin{enumerate}
\item \textbf{Nommer} le produit et préciser son origine biologique : sécrétion glandulaire (venin : glande à venin ; gelée royale : glandes hypopharyngiennes/mandibulaires des nourrices) ou transformation de matières végétales récoltées (miel : nectar + invertase + déshydratation ; propolis : résines de bourgeons + cire + enzymes salivaires).
\item \textbf{Donner la composition} chimique dominante avec unités : miel \SI{\sim 80}{\percent} glucides (glucose \SI{31}{\percent}, fructose \SI{38}{\percent}), \SI{\sim 17}{\percent} eau, enzymes (glucose-oxydase $\rightarrow$ $\ce{H2O2}$), flavonoïdes ; venin : peptides (mélittine \SI{50}{\percent} poids sec, apamine, peptide 401), amines (histamine), enzymes (phospholipase A2) ; propolis : \SI{50}{\percent} résines, \SI{30}{\percent} cire, \SI{10}{\percent} huiles essentielles, \SI{5}{\percent} pollen, riches en flavonoïdes (galangine, pinocembrine).
\item \textbf{Expliquer le mode d'utilisation} : miel (aliment énergétique, pansement topique, sirop antitussif) ; venin (apipuncture : piqûres contrôlées sur points d'acupuncture ; injections purifiées pour désensibilisation ou arthrite) ; propolis (teinture-mère \SI{10-30}{\percent} éthanol, pommade, spray buccal) ; gelée royale (cure orale fraîche ou lyophilisée, \SI{300-600}{mg.j^{-1}}).
\item \textbf{Justifier l'effet} par la molécule active : mélittine (anti-inflammatoire, anti-arthrosique, lyse membranaire bactérienne) ; glucose-oxydase du miel $\rightarrow$ $\ce{H2O2}$ (antiseptique) + forte osmolarité (cicatrisation par débridement autolytique) ; flavonoïdes propolis (inhibition synthèse prostaglandines, antibactérien large spectre).
\item \textbf{Mentionner les précautions obligatoires} : risque d'allergie IgE-dépendante (choc anaphylactique) au venin (\SI{1-3}{\percent} population) et aux protéines du miel ; miel \textbf{interdit} aux nourrissons \SI{< 12}{mois} (risque botulisme infantile par spores \emph{Clostridium botulinum} résistantes à l'acidité du miel) ; propolis : risque dermatite de contact.
\end{enumerate}
\textbf{Formule à retenir :} Valeur énergétique miel $\approx \SI{3200}{kcal.kg^{-1}} = \SI{13,4}{MJ.kg^{-1}}$.
\end{methode}

\begin{exoresolu}[Le miel : produit alimentaire et thérapeutique]
\textbf{Énoncé.} Un apiculteur de l'Ouest-Cameroun récolte \SI{25}{kg} de miel par ruche et par campagne. a) Précise l'origine et la composition du miel. b) Calcule l'énergie totale fournie par la récolte d'une ruche, sachant que le miel apporte \SI{3200}{kcal.kg^{-1}}. c) Cite deux utilisations thérapeutiques du miel et une précaution d'emploi.
\tcblower
\textbf{Solution détaillée.}

a) \textbf{Origine et composition.} Le miel est élaboré par les abeilles ouvrières butineuses à partir du nectar des fleurs (sucrose), enrichi de sécrétions enzymatiques (invertase hydrolysant le sucrose en glucose + fructose) puis déshydraté par ventilation dans les alvéoles operculées (taux d'humidité final \SI{< 18}{\percent}). Composition moyenne : \SI{\sim 80}{\percent} glucides (majoritairement fructose et glucose), \SI{\sim 17}{\percent} eau, \SI{< 0,5}{\percent} protéines/acides aminés, traces de sels minéraux (K, Ca, Mg), vitamines (groupe B, C), enzymes (glucose-oxydase, catalase, diastase) et polyphénols (flavonoïdes) variables selon l'origine florale.

b) \textbf{Calcul de l'énergie.} Relation $E = m \times e$ avec $m = \SI{25}{kg}$ et $e = \SI{3200}{kcal.kg^{-1}}$ :
\[ E = 25 \times 3200 = \SI{80 000}{kcal} = \SI{8,0 \times 10^4}{kcal} \]
Conversion en kJ : $1 \text{ kcal} = \SI{4,184}{kJ}$ $\Rightarrow$ $E = 80 000 \times 4,184 = \SI{334 720}{kJ} \approx \SI{335}{MJ}$.
Cette énergie couvre les besoins journaliers d'un adulte (\SI{\sim 2500}{kcal.j^{-1}}) pendant $80 000 / 2500 = 32$ jours.

c) \textbf{Utilisations thérapeutiques.} (1) Cicatrisant et antiseptique des plaies chroniques, brûlures, ulcères diabétiques : application topique de miel médical (stérilisé par rayons $\gamma$) ; l'osmolarité élevée assèche le lit de la plaie (effet débridant) et la glucose-oxydase génère du $\ce{H2O2}$ à faible dose continue (bactéricide sans cytotoxicité). (2) Adoucissant des voies respiratoires : \SI{10}{g} de miel dans eau tiède (\SI{< 40}{\celsius} pour préserver enzymes) soulagent la toux nocturne chez l'enfant \SI{> 1}{an} (efficacité supérieure au dextrométhorphane dans certaines études). \textbf{Précaution absolue :} ne jamais donner de miel à un nourrisson de moins de 12 mois (risque botulisme infantile).
\end{exoresolu}

\begin{attention}[Précautions obligatoires aux produits de la ruche (Bac)]
\begin{itemize}
\item \textbf{Venin :} Test cutané préalable obligatoire (injection intradermique de \SI{0,01}{mL} venin dilué 1/1000, lecture à 20 min). Contre-indiqué : allergie connue, maladies auto-immunes actives, grossesse, prise de bêta-bloquants.
\item \textbf{Miel :} \textbf{Interdit < 1 an}. Stérilisation par rayons $\gamma$ (\SI{10}{kGy}) pour usage médical (pansements) sans dégrader les enzymes.
\item \textbf{Propolis :} Test cutané (pli du coude) 24h avant usage généralisé (risque eczéma de contact \SI{\sim 3}{\percent}).
\end{itemize}
\end{attention}

\courssub{Lutte biologique : Insectes auxiliaires et protection des cultures fourragères}

\begin{definition}[Lutte biologique]
La \cle{lutte biologique} est une méthode de protection des cultures qui utilise des organismes vivants (\cle{auxiliaires} : prédateurs, parasitoïdes, pathogènes) pour maintenir les populations de ravageurs en dessous du seuil de nuisibilité économique, en substitution ou en complément de la lutte chimique.
\end{definition}

\begin{methode}[Démarche pour analyser un dispositif de lutte biologique]
\begin{enumerate}
\item \textbf{Définir} le type de lutte : acclimatation (introduction classique), augmentation (lâchers inundatifs ou inoculatifs), conservation (aménagement d'habitats : haies, bandes fleuries, refuges).
\item \textbf{Identifier le couple ravageur / auxiliaire} : nom scientifique, stade cible (œuf, larve, adulte), spécificité (monophage, oligophage, polyphage).
\item \textbf{Préciser le mode d'action} : \textbf{Prédation} (consommation directe : coccinelle $\rightarrow$ pucerons ; libellule adulte $\rightarrow$ diptères volants ; larve de libellule \textbf{aquatique} $\rightarrow$ larves de moustiques) ou \textbf{Parasitisme} (guêpe parasitoïde pondant dans l'hôte : \emph{Trichogramma} spp. dans œufs de lépidoptères ; \emph{Aphidius} spp. dans pucerons).
\item \textbf{Chiffrer l'efficacité} : taux de consommation journalier, taux de parasitisme, ratio lâcher/ravageur.
\item \textbf{Comparer} avec la lutte chimique : avantages (pas de résidus, pas de résistance, sélectivité, préservation pollinisateurs/auxiliaires, coût marginal nul après installation) et limites (délai d'action, spécificité exige une identification précise, sensibilité aux pesticides résiduels, coût initial d'élevage).
\item \textbf{Conclure} sur l'intérêt agronomique et zootechnique : protection des cultures vivrières (maïs, niébé, manioc) et fourragères (brachiaria, stylosanthès) $\rightarrow$ fourrages sains, productifs, sans résidus $\rightarrow$ meilleure production animale (lait, viande).
\end{enumerate}
\textbf{Réflexe :} Ne jamais écrire « la coccinelle tue les maladies » : elle consomme des \emph{ravageurs} (vecteurs éventuels de virus), pas des agents pathogènes. Distinguer larve de libellule (aquatique, prédateur de larves de moustiques) et adulte (aérien, prédateur de moustiques adultes).
\end{methode}

\begin{exoresolu}[La coccinelle, auxiliaire des cultures de niébé]
\textbf{Énoncé.} Dans une plantation de niébé de la région du Nord, on estime une population de \SI{60 000}{pucerons} (\emph{Aphis craccivora}). Un lâcher de coccinelles (\emph{Cheilomenes sulphurea}) est réalisé ; chaque coccinelle consomme en moyenne \SI{50}{pucerons.jour^{-1}}. a) Explique le principe de cette méthode de protection. b) Calcule le nombre minimal de coccinelles nécessaires pour éliminer les pucerons en 3 jours. c) Donne deux avantages de cette méthode par rapport aux insecticides chimiques.
\tcblower
\textbf{Solution détaillée.}

a) \textbf{Principe.} Il s'agit de la lutte biologique par \textbf{augmentation inoculative} : introduction d'un prédateur indigène, la coccinelle \emph{Cheilomenes sulphurea}, auxiliaire spécifique des pucerons. Les adultes et larves se nourrissent de pucerons qui prélèvent la sève (affaiblissement, galle) et transmettent le virus de la mosaïque du niébé (BCMV). La réduction de la population de ravageurs sous le seuil de nuisibilité protège la culture sans pesticide.

b) \textbf{Calcul.} Consommation d'une coccinelle en 3 jours : $50 \times 3 = 150$ pucerons.
Nombre minimal de coccinelles :
\[ N = \frac{60 000}{150} = 400 \text{ coccinelles.} \]
En pratique, on lâche \SI{1,5}{à 2} fois ce nombre (\SI{600-800}{individus}) pour compenser la dispersion, la mortalité et la reproduction non immédiate.

c) \textbf{Avantages.} (1) Absence de résidus chimiques toxiques sur les gousses (sécurité alimentaire) et dans l'environnement (préservation des pollinisateurs \emph{Apis mellifera}, \emph{Xylocopa} spp. et de la faune auxiliaire) ; (2) Absence de phénomène de résistance génétique chez le ravageur (coévolution prédateur-proie) et coût faible à long terme (les coccinelles s'installent et se reproduisent sur place si l'habitat est favorable).
\end{exoresolu}

\begin{experience}[Suivi d'un lâcher de coccinelles en parcelle de niébé (TP / Projet)]
\textbf{Matériel.} Parcelle de niébé infestée (pucerons \emph{Aphis craccivora}), coccinelles \emph{Cheilomenes sulphurea} (élevées ou achetées), cadres de comptage (\SI{1}{m^2}), pièges jaunes, loupe binoculaire.
\textbf{Protocole.}
\begin{enumerate}
\item Délimiter 3 parcelles témoins (sans lâcher) et 3 parcelles traitées (lâcher de \SI{500}{coccinelles.ha^{-1}}).
\item Compter les pucerons sur 20 plants/parcelle (feuilles apicales, moyennes, basales) à J0, J3, J7, J14.
\item Observer la présence de larves, nymphes et adultes de coccinelles (reproduction ?).
\item Noter les symptômes viraux (mosaïque, enroulement).
\end{enumerate}
\textbf{Observations attendues.} Baisse significative de la population de pucerons dans les parcelles traitées dès J3-J7 ; apparition de larves de coccinelles (stade 2-3) à J14 prouvant la reproduction ; faible incidence virale vs témoins.
\textbf{Interprétation.} La coccinelle s'établit, se reproduit et régule durablement le ravageur. La lutte biologique est validée comme alternative aux néonicotinoïdes (interdits au Cameroun depuis 2019 pour usage agricole).
\end{experience}

\begin{savaistu}[Lutte biologique au Cameroun : le programme national]
L'IRAD (Institut de Recherche Agricole pour le Développement) et le CNV (Centre National de Vulgarisation) promeuvent la lutte biologique contre le foreur du maïs (\emph{Busseola fusca}, \emph{Chilo partellus}) par le parasitoïde \emph{Cotesia sesamiae} (guêpe braconide) et contre la mouche des fruits (\emph{Bactrocera dorsalis}, invasive depuis 2003) par \emph{Fopius arisanus} et \emph{Diachasmimorpha longicaudata}. La conservation des haies vives et des bandes fleuries (\emph{Chromolaena odorata}, \emph{Tithonia diversifolia}) favorise les auxiliaires naturels (syrphes, chrysopes, araignées).
\end{savaistu}

\courssub{Techniques d'entomoculture : De la collecte à l'élevage industriel pour l'alimentation animale}

\begin{methode}[Démarche pour présenter une filière d'élevage d'insectes (entomoculture)]
\begin{enumerate}
\item \textbf{Distinguer} les deux niveaux : \textbf{Collecte saisonnière} en milieu naturel (termites ailés, chenilles, vers blancs) — aléatoire, non durable si surexploitation — et \textbf{Entomoculture maîtrisée} (élevage en milieu contrôlé : grillons \emph{Gryllus bimaculatus}, vers de farine \emph{Tenebrio molitor}, mouche soldat noire \emph{Hermetia illucens}, abeilles).
\item \textbf{Décrire les étapes de l'élevage industriel} (ex. Mouche Soldat Noire - MSN) :
\begin{enumerate}
\item \textbf{Reproduction :} Cages d'adultes (\SI{25-30}{\celsius}, \SI{60-70}{\percent} HR, lumière naturelle ou LED \SI{5000}{K}), substrat de ponte (bois fendu, carton ondulé) au-dessus d'appât odorant (déchets fermentés).
\item \textbf{Éclosion \& Néonates :} Œufs éclosent en \SI{4}{jours} ; néonates transférés sur substrat d'engraissement.
\item \textbf{Engraissement (10-14 j) :} Substrat = déchets organiques (résidus maraîchers, drêches de brasserie, fientes de volaille, déchets de cuisine) humidifiés \SI{60-70}{\percent}. Densité \SI{5-10}{larves.cm^{-2}}. Température optimale \SI{27-30}{\celsius}. Auto-harvesting : pré-nymphes (stade 6, couleur noire) migrent hors du substrat $\rightarrow$ rampe de collecte.
\item \textbf{Récolte \& Transformation :} Tri (tamisage), jeûne \SI{24}{h} (vidange tube digestif), euthanasie (congélation \SI{-20}{\celsius} ou échaudage \SI{90}{\celsius} \SI{1}{min}), séchage (four \SI{60}{\celsius} \SI{12}{h} ou solaire hybride), mouture $\rightarrow$ \textbf{farine de larves} (\SI{45-50}{\percent} prot. MS, \SI{25-30}{\percent} lipides).
\end{enumerate}
\item \textbf{Justifier chaque condition} par la biologie : cycle court (\SI{3}{semaines} œuf-adulte), taux de conversion substrat/biomasse élevé (\cle{TCR} \SI{\sim 3-5}{kg substrat frais / kg larves fraîches}), thermorégulation comportementale (grégarisme).
\item \textbf{Préciser le débouché principal :} \textbf{Alimentation animale} (farine de larves MSN incorporée \SI{5-15}{\percent} dans aliments volailles/poissons, substitution partielle farine de poisson \SI{30-50}{\percent}) ; \textbf{Alimentation humaine} (grillons, vers de farine : farine enrichie pains, biscuits, pâtes — règlementation Novel Food UE, Codex Alimentarius en cours au Cameroun).
\item \textbf{Évaluer la durabilité (ACV simplifiée) :} Empreinte eau \SI{< 100}{L.kg^{-1}} prot. (vs \SI{15 000}{L} bœuf), surface \SI{< 10}{m^2.tonne^{-1}} prot., GES \SI{< 1}{kg eqCO2.kg^{-1}} prot., valorisation déchets (économie circulaire).
\end{enumerate}
\end{methode}

\begin{exoresolu}[Élevage de larves de mouche soldat noire pour l'alimentation des volailles]
\textbf{Énoncé.} Une coopérative de Bafoussam élève des larves de mouche soldat noire (\emph{Hermetia illucens}) sur des résidus de maraîchage. \SI{100}{kg} de résidus produisent \SI{20}{kg} de larves fraîches contenant \SI{35}{\percent} de protéines en matière sèche, les larves fraîches contenant \SI{65}{\percent} d'eau. a) Décris les étapes de cette entomoculture. b) Calcule la masse de protéines obtenue à partir de \SI{500}{kg} de résidus. c) Donne deux arguments de durabilité de cette filière.
\tcblower
\textbf{Solution détaillée.}

a) \textbf{Étapes de l'élevage (MSN).}
\begin{enumerate}
\item \textbf{Reproduction :} Adultes en cage filet (\SI{3 \times 3 \times 2}{m}), \SI{28}{\celsius}, \SI{70}{\percent} HR, photopériode 12h. Ponte sur carton ondulé placé au-dessus de bac d'appât (son de maïs fermenté + levure).
\item \textbf{Éclosion (J4) :} Cartons transférés en nurserie ; néonates (1 mm) tombent dans substrat starter (son humide).
\item \textbf{Engraissement (J4-J18) :} Transfert en bacs d'engraissement (\SI{1}{m^2}, \SI{10}{cm} substrat). Alimentation \emph{ad libitum} avec résidus maraîchers (feuilles de chou, épluchures tomates, carottes) humidifiés \SI{65}{\percent}. Brassage quotidien pour aération et homogénéité.
\item \textbf{Auto-récolte (J18-J20) :} Pré-nymphes (noires, \SI{18-20}{mm}) sortent du substrat, grimpent sur rampe \SI{35}{\degree} $\rightarrow$ tombent dans bac de collecte.
\item \textbf{Transformation :} Jeûne \SI{24}{h}, blanchiment \SI{90}{\celsius} \SI{1}{min}, séchage four \SI{60}{\celsius} \SI{12}{h} (humidité finale \SI{< 8}{\percent}), broyage $\rightarrow$ farine protéique (\SI{48}{\percent} prot. MS). Incorporation à \SI{10}{\percent} dans aliment poulets de chair (remplace \SI{50}{\percent} farine de poisson).
\end{enumerate}

b) \textbf{Calculs.}
Masse larves fraîches issues de \SI{500}{kg} résidus :
\[ m_{\text{larves fraîches}} = 500 \times \frac{20}{100} = \SI{100}{kg} \]
Matière sèche des larves (teneur en eau \SI{65}{\percent} $\Rightarrow$ MS = \SI{35}{\percent}) :
\[ m_{\text{MS}} = 100 \times 0,35 = \SI{35}{kg} \]
Masse de protéines (teneur \SI{35}{\percent} en MS) :
\[ m_{\text{prot}} = 35 \times 0,35 = \SI{12,25}{kg} \]
On obtient donc \SI{12,25}{kg} de protéines brutes à partir de \SI{500}{kg} de résidus végétaux sans valeur alimentaire directe pour les volailles.

c) \textbf{Arguments de durabilité.}
(1) \textbf{Valorisation des déchets :} \SI{500}{kg} de résidus (polluants s'ils sont déversés : lixiviats, $\ce{CH4}$, odeurs) sont convertis en biomasse noble + frass (résidu d'élevage = engrais organique riche en N, P, K, chitine $\rightarrow$ stimulateur défense plantes).
(2) \textbf{Faible empreinte écologique :} Besoin en eau \SI{\sim 50}{L.kg^{-1}} larves fraîches (vs \SI{2500}{L} pour soja, \SI{15 000}{L} pour bœuf) ; surface au sol \SI{10}{fois} moindre ; émissions GES \SI{< 1}{kg eqCO2.kg^{-1}} prot. ; pas de déforestation.
\end{exoresolu}

\begin{popfigure}
\begin{tikzpicture}[scale=0.8, every node/.style={font=\small}]
% Couleurs
\colorlet{eggcol}{poporange!80}
\colorlet{larvcol}{popgreen!70}
\colorlet{prepcol}{poppurple!70}
\colorlet{pupcol}{popblue!70}
\colorlet{adultcol}{poppink!70}
\colorlet{arcol}{popdark}

% Nœuds du cycle
\node[ellipse, draw=arcol, thick, fill=eggcol, minimum width=2.2cm, minimum height=1.2cm, align=center] (egg) at (0,0){\textbf{Œufs}\\(0,5--1 mm)\\\SI{4}{jours}};
\node[ellipse, draw=arcol, thick, fill=larvcol, minimum width=2.5cm, minimum height=1.3cm, align=center] (larv) at (5,0){\textbf{Larves (6 stades)}\\(1 $\rightarrow$ 20 mm)\\\SI{10-14}{jours}\\\textit{Engraissement déchets}};
\node[ellipse, draw=arcol, thick, fill=prepcol, minimum width=2.2cm, minimum height=1.2cm, align=center] (prep) at (10,0){\textbf{Pré-nymphes}\\(Noires, migrantes)\\\SI{2-3}{jours}\\\textit{Auto-récolte}};
\node[ellipse, draw=arcol, thick, fill=pupcol, minimum width=2.2cm, minimum height=1.2cm, align=center] (pup) at (10,-3.5){\textbf{Nymphes}\\(Immobile)\\\SI{7-10}{jours}};
\node[ellipse, draw=arcol, thick, fill=adultcol, minimum width=2.2cm, minimum height=1.2cm, align=center] (adult) at (5,-3.5){\textbf{Adultes}\\(Mouches \SI{15-20}{mm})\\\SI{5-8}{jours}\\\textit{Accouplement, Ponte}};

% Flèches
\draw[->, >=Stealth, thick, arcol] (egg.east) -- node[above, font=\tiny] {Éclosion} (larv.west);
\draw[->, >=Stealth, thick, arcol] (larv.east) -- node[above, font=\tiny] {Migration} (prep.west);
\draw[->, >=Stealth, thick, arcol] (prep.south) -- node[right, font=\tiny] {Nymphose} (pup.north);
\draw[->, >=Stealth, thick, arcol] (pup.west) -- node[below, font=\tiny] {Émergence} (adult.east);
\draw[->, >=Stealth, thick, arcol] (adult.north) to[out=90, in=-90, looseness=1.5] node[left, font=\tiny, align=center] {Ponte sur\\carton/appât} (egg.south);

% Conditions
\node[draw=popdark, fill=popcream, rounded corners, align=left, font=\tiny, anchor=north west] at (-2.5, 1.5) {
\textbf{Conditions optimales} \\
Température : \SI{27-30}{\celsius} \\
Hygrométrie : \SI{60-70}{\percent} \\
Photopériode : 12L:12D \\
Substrat : Déchets organiques \\ humidifiés \SI{65}{\percent}
};

% Légende interne
\node[font=\tiny, align=center, text=popmuted] at (5, -5.5) {Cycle complet : \SI{\sim 3}{semaines} à \SI{28}{\celsius}};
\end{tikzpicture}
\legende{Cycle biologique de la Mouche Soldat Noire (\emph{Hermetia illucens}) en entomoculture maîtrisée. Les flèches indiquent les transitions développementales ; les durées sont données pour \SI{28}{\celsius}. L'auto-récolte des pré-nymphes (migration) est l'innovation clé facilitant l'industrialisation.}
\end{popfigure}

\begin{popfigure}
\begin{tikzpicture}
\begin{axis}[
    ybar,
    bar width=18pt,
    width=\linewidth,
    height=7cm,
    enlarge x limits=0.25,
    ylabel={Masse d'aliment pour 30 g de protéines (g)},
    ylabel style={font=\small},
    symbolic x coords={Bœuf (frais), Termites (séchés), Chenilles (séchées), Farine MSN},
    xtick=data,
    xticklabel style={rotate=30, anchor=east, font=\small},
    ymin=0, ymax=180,
    ymajorgrids=true,
    grid style=dashed,
    nodes near coords,
    nodes near coords style={font=\small, anchor=south, yshift=2pt},
    title={Efficacité protéique : Masse nécessaire pour couvrir 50\% des besoins journaliers (30 g prot.)},
    title style={font=\small, yshift=4pt},
]
\addplot[fill=popblue!70] coordinates {(Bœuf (frais),150) (Termites (séchés),75) (Chenilles (séchées),60) (Farine MSN,67)};
\end{axis}
\end{tikzpicture}
\legende{Comparaison de la masse d'aliment nécessaire pour apporter \SI{30}{g} de protéines (moitié du besoin journalier d'un adolescent de 50 kg). Données : Bœuf frais \SI{20}{\percent} prot. ; Termites séchés \SI{40}{\percent} prot. MS ; Chenilles séchées \SI{50}{\percent} prot. MS ; Farine MSN \SI{45}{\percent} prot. MS. Les insectes séchés sont 2 à 2,5 fois plus concentrés que la viande fraîche.}
\end{popfigure}

\courssub{Communication, sensibilisation et acceptabilité sociale}

\begin{methode}[Démarche de conception d'un support de sensibilisation (affiche, sketch, dépliant, vidéo courte)]
\begin{enumerate}
\item \textbf{Définir la cible} (élèves lycée, ménagères marché, éleveurs avicoles, décideurs) et \textbf{l'objectif unique mesurable} (ex. : « Augmenter de 20\% l'intention de consommer de la farine de grillons enrichie chez les élèves »).
\item \textbf{Choisir le canal} adapté : affiche illustrée (marché, centre de santé), sketch radiophonique en langue locale (fulfulde, bassa, ewondo, pidgin), dépliant avec tableau comparatif (lycées, formations agricoles), vidéo TikTok/WhatsApp (jeunes urbains), dégustation commentée (foire agro-pastorale).
\item \textbf{Sélectionner 3 messages clés} fondés sur des données scientifiques vérifiées :
\begin{enumerate}
\item \textbf{Nutrition :} « \SI{100}{g} de chenilles séchées = \SI{150}{g} de bœuf en protéines + 3 fois plus de fer. »
\item \textbf{Santé / Économie :} « Élever des larves de mouche soldat noire sur déchets = farine 2 fois moins chère que la farine de poisson importée. »
\item \textbf{Environnement :} « 1 kg protéines insectes = \SI{100}{fois} moins d'eau, \SI{10}{fois} moins de surface, \SI{100}{fois} moins de $\ce{CO2}$ que 1 kg bœuf. »
\end{enumerate}
\item \textbf{Soigner la forme} : slogan court (\emph{« Protéines locales, planète gagnante »}), photos d'insectes familiers (pas d'insectes « effrayants »), tableau comparatif simple (icônes + chiffres), langage accessible (éviter jargon), respect des tabous culturels (ne pas imposer, valoriser l'existant).
\item \textbf{Prévoir l'évaluation d'impact} : questionnaire anonyme avant/après (connaissances, attitudes, intentions), comptage participants dégustation, suivi des demandes de formation/équipement entomoculture.
\end{enumerate}
\textbf{Réflexe :} Un bon outil ne dit jamais « Mangez des insectes ! » ; il \textbf{démontre, chiffres à l'appui}, pourquoi les insectes déjà consommés traditionnellement méritent une place \textbf{valorisée, sécurisée et scalable} dans l'alimentation humaine et animale.
\end{methode}

\begin{exoresolu}[Conception d'une affiche pour un club de santé scolaire]
\textbf{Énoncé.} Le club de santé de ton lycée veut sensibiliser les élèves sur la valeur nutritionnelle des chenilles comestibles. Propose un plan d'affiche en quatre rubriques, avec pour chacune le message scientifique exact à faire figurer, puis indique comment évaluer l'efficacité de la campagne.
\tcblower
\textbf{Solution détaillée.}

\textbf{Rubrique 1 — Titre accrocheur \& Visuel :} « Les chenilles : des protéines de champion à portée de main ! » — Photo appétissante de chenilles \emph{Imbrasia} grillées dans une assiette colorée (riz, légumes), pas d'insecte vivant isolé.

\textbf{Rubrique 2 — Message nutritionnel chiffré (Tableau comparatif) :}
\begin{center}
\begin{tabular}{|c|c|c|}\hline
\rowcolor{popblueL} \textbf{Nutriment (pour 100 g MS)} & \textbf{Chenilles séchées} & \textbf{Bœuf frais} \\ \hline
Protéines & \SI{\sim 53}{g} & \SI{\sim 20}{g} \\ \hline
Fer & \SI{\sim 35}{mg} & \SI{\sim 2,5}{mg} \\ \hline
Zinc & \SI{\sim 15}{mg} & \SI{\sim 4,5}{mg} \\ \hline
Calcium & \SI{\sim 200}{mg} & \SI{\sim 12}{mg} \\ \hline
\end{tabular}
\end{center}
\emph{Message :} « \SI{75}{g} de chenilles séchées = \SI{150}{g} de viande en protéines, mais 14 fois plus de fer et 3 fois plus de zinc ! »

\textbf{Rubrique 3 — Message santé, économie \& environnement :} « Collectées en forêt ou élevées sur feuilles de \emph{Triplochiton}, les chenilles génèrent des revenus pour les mamans du village (vente au kg), coûtent 3 fois moins cher que le bœuf au kg de protéines, et ne détruisent pas la forêt. »

\textbf{Rubrique 4 — Pratique \& Précautions :} « \textbf{Préparation sûre :} 1. Trier (ôter débris). 2. Blancher \SI{5}{min} eau bouillante. 3. Griller/Frire \SI{10}{min} (cuisson à cœur \SI{> 70}{\celsius}). 4. Consommer ou sécher au soleil \SI{2}{jours} (conservation 1 an). \textbf{Attention :} Bien cuire (destroy germes/parasites). Déconseillé si allergie aux crustacés/acariens (protéines tropomyosine croisées). »

\textbf{Évaluation de la campagne :}
\begin{enumerate}
\item Questionnaire anonyme \textbf{avant} exposition (J-1) : 5 QCM (teneur prot., fer, prix, prépa, envie).
\item Exposition affiche + dégustation encadrée (brochettes chenilles grillées) pendant la récréation (J0).
\item Questionnaire \textbf{après} (J+7) : mêmes items + « As-tu parlé des chenilles à ta famille ? ».
\item \textbf{Indicateurs de succès :} Hausse \SI{> 30}{\percent} réponses correctes nutrition ; \SI{> 50}{\percent} intentions de consommation positives ; \SI{> 20}{\percent} élèves ayant demandé la recette à la maison.
\end{enumerate}
\end{exoresolu}

\courssub{Synthèse et erreurs à éviter}

\begin{aretenir}[Points clés à maîtriser pour le Bac]
\begin{itemize}
\item \textbf{Entomophagie} = levier majeur pour la sécurité alimentaire (protéines, fer, zinc) et la durabilité (faible empreinte eau/surface/GES, économie circulaire).
\item \textbf{Insectes locaux stars :} Termites ailés (\emph{Macrotermes}), Vers blancs (\emph{Rhynchophorus}), Chenilles (\emph{Imbrasia}, \emph{Cirina}), Fourmis (\emph{Carebara}, reines \emph{Dorylus}), Abeilles (couvain).
\item \textbf{Produits ruche :} Miel (énergétique, cicatrisant $\ce{H2O2}$/osmolarité), Venin (mélittine anti-inflammatoire, apipuncture), Propolis (flavonoïdes antibactériens), Gelée royale (stimulante).
\item \textbf{Lutte biologique :} Coccinelles (pucerons), Libellules adultes (moustiques volants) / larves (larves moustiques aquatiques), Parasitoïdes (\emph{Trichogramma}, \emph{Cotesia}). Alternative aux pesticides (résidus, résistance, biodiversité).
\item \textbf{Entomoculture industrielle :} Mouche Soldat Noire (\emph{Hermetia illucens}) sur déchets $\rightarrow$ farine protéique pour volailles/poissons (substitution farine poisson). Cycle \SI{3}{semaines}, TCR \SI{\sim 4}{kg/kg}.
\item \textbf{Calculs nutritionnels :} \textbf{Toujours} préciser base (MS ou MF). Convertir MF $\leftrightarrow$ MS via teneur en eau. Comparer masses pour même apport protéique.
\item \textbf{Sensibilisation :} Cible + Objectif + 3 messages chiffrés + Canal adapté + Évaluation avant/après.
\end{itemize}
\end{aretenir}

\begin{attention}[Les 5 erreurs qui coûtent des points au Bac]
\begin{enumerate}
\item \textbf{Confondre matière fraîche (MF) et matière sèche (MS).} Les teneurs en protéines des insectes (\SI{40-60}{\percent}) sont exprimées en \emph{matière sèche} ; un insecte frais contient \SI{60-75}{\percent} d'eau. Oublier cette conversion fausse tous les calculs comparatifs (ex. : \SI{100}{g} termites frais $\neq$ \SI{40}{g} prot., mais \SI{\sim 10}{g} prot.).
\item \textbf{Donner un pourcentage sans référence ni unité.} Écrire « les termites contiennent 40 » est nul. Écrire : « environ \SI{40}{\percent} de protéines \textbf{en matière sèche} ». Toute donnée chiffrée exige son unité et sa base de référence.
\item \textbf{Attribuer aux auxiliaires un rôle contre les maladies.} La coccinelle et la libellule consomment des \emph{ravageurs} (pucerons, moustiques), pas des agents pathogènes (virus, bactéries, champignons). La lutte biologique ne remplace ni la vaccination, ni les traitements vétérinaires curatifs.
\item \textbf{Oublier les précautions d'emploi des produits de la ruche.} Citer le miel ou le venin sans mentionner le risque allergique (choc anaphylactique, test cutané préalable) ou la contre-indication formelle du miel chez le nourrisson \SI{< 12}{mois} (botulisme infantile) rend la réponse incomplète et dangereuse.
\item \textbf{Décrire une technique sans la justifier biologiquement.} En entomoculture comme en collecte, chaque étape doit être reliée à la biologie de l'insecte : essaimage lié aux premières pluies et crépuscule ; température \SI{27-30}{\celsius} pour croissance larvaire MSN ; jeûne avant abattage pour vider le tube digestif. Une liste de gestes sans justification scientifique ne vaut que la moitié des points.
\end{enumerate}
\end{attention}

\newpage

\coursec{Exercices}

\courssub{Je vérifie mes connaissances}

\begin{exercice}[QCM : insectes et leurs rôles]
Pour chaque affirmation, choisir la (ou les) bonne(s) réponse(s) en justifiant brièvement.
\begin{enumerate}
\item L'entomophagie désigne :
\begin{enumerate}
\item l'élevage industriel des insectes ;
\item la consommation d'insectes par l'Homme ou les animaux d'élevage ;
\item la lutte contre les insectes ravageurs ;
\item l'étude scientifique des insectes.
\end{enumerate}
\item Parmi les insectes suivants, lequel est un \cle{auxiliaire} de lutte biologique contre les pucerons ?
\begin{enumerate}
\item la coccinelle ; \item le criquet ; \item le ver blanc du palmier ; \item le moustique.
\end{enumerate}
\item Le miel est produit par l'abeille à partir :
\begin{enumerate}
\item du pollen seul ; \item du nectar des fleurs ; \item de la sève des arbres ; \item de la propolis.
\end{enumerate}
\item Le \cle{frass} est :
\begin{enumerate}
\item le venin des fourmis ; \item le substrat de culture des champignons ;
\item le mélange de déjections et de résidus de substrat produit par les larves, utilisable comme engrais organique ;
\item le nom scientifique du ver de palmier.
\end{enumerate}
\end{enumerate}
\end{exercice}

\begin{exercice}[Vrai ou faux justifié]
Dire si chacune des affirmations suivantes est vraie ou fausse, puis justifier la réponse en une à trois phrases.
\begin{enumerate}
\item Les termites ailés (\emph{Macrotermes} sp.) consommés au Cameroun sont plus riches en protéines (sur matière sèche) que la viande de bœuf.
\item Le venin d'abeille ne présente aucun intérêt thérapeutique connu.
\item La libellule est un insecte auxiliaire car ses larves et ses adultes consomment de grandes quantités de moustiques.
\item Les larves de \emph{Hermetia illucens} (mouche soldat noire) peuvent transformer des déchets organiques en biomasse riche en protéines pour l'alimentation des poulets.
\item L'élevage d'insectes émet autant de gaz à effet de serre par kilogramme de protéine produite que l'élevage bovin.
\end{enumerate}
\end{exercice}

\begin{exercice}[Définitions]
Définir en une phrase claire et scientifiquement correcte chacun des termes suivants :
\begin{enumerate}
\item Entomophagie ;
\item Entomoculture ;
\item Lutte biologique (par insectes auxiliaires) ;
\item Bioconversion ;
\item Propolis.
\end{enumerate}
\end{exercice}

\begin{exercice}[Calcul direct : apport protéique des chenilles]
Un adolescent de 15 ans a des apports journaliers recommandés (AJR) en protéines de \SI{52}{g/jour}. Les chenilles de \emph{Imbrasia} sp. séchées contiennent \SI{55}{g} de protéines pour \SI{100}{g} de matière sèche.
\begin{enumerate}
\item Calculer la masse de chenilles séchées nécessaire pour couvrir la totalité des AJR en protéines de cet adolescent.
\item La même quantité de protéines serait apportée par du bœuf maigre contenant \SI{20}{g} de protéines pour \SI{100}{g} de viande fraîche. Calculer la masse de bœuf correspondante.
\item Sachant que \SI{1}{kg} de chenilles séchées coûte \SI{3500}{FCFA} et \SI{1}{kg} de bœuf \SI{3000}{FCFA} au marché de Yaoundé, comparer le coût des deux sources pour couvrir ces AJR. Conclure.
\end{enumerate}
\end{exercice}

\courssub{Je m'entraîne}

\begin{exercice}[Analyse d'un tableau de composition]
Le tableau ci-dessous donne la composition centésimale (pour \SI{100}{g} de matière sèche) de quelques aliments.
\begin{center}
\begin{tabularx}{\linewidth}{|l|X|X|X|X|}\hline
\rowcolor{popblueL} \textbf{Aliment} & \textbf{Protéines (g)} & \textbf{Lipides (g)} & \textbf{Fer (mg)} & \textbf{Zinc (mg)} \\ \hline
Termites ailés séchés & 38 & 45 & 12 & 9 \\ \hline
Vers de palmier séchés & 25 & 60 & 6 & 8 \\ \hline
Bœuf (base sèche) & 65 & 25 & 8 & 12 \\ \hline
Sauterelles séchées & 60 & 8 & 10 & 14 \\ \hline
\end{tabularx}
\end{center}
\begin{enumerate}
\item Classer ces quatre aliments par ordre décroissant de richesse en protéines.
\item Quel aliment est le plus énergétique ? Justifier en utilisant les valeurs énergétiques : protéines \SI{17}{kJ/g}, lipides \SI{38}{kJ/g}.
\item Un enfant souffre d'anémie ferriprive. Quel aliment, parmi les quatre, conseilleriez-vous en priorité ? Justifier.
\item Expliquer pourquoi les comparaisons nutritionnelles doivent se faire sur la matière sèche et non sur la matière fraîche.
\end{enumerate}
\end{exercice}

\begin{exercice}[Produits de la ruche et leurs utilisations]
\begin{enumerate}
\item Citer trois produits issus de la ruche et, pour chacun, donner une utilisation alimentaire ou thérapeutique traditionnelle au Cameroun.
\item Le miel contient environ \SI{80}{\percent} de sucres (glucose, fructose) et possède une activité antibactérienne. Citer deux facteurs expliquant cette activité antibactérienne.
\item L'apithérapie utilise le venin d'abeille dans le traitement de certaines affections rhumatismales. Nommer le principal composé actif anti-inflammatoire du venin d'abeille et préciser une précaution d'emploi indispensable.
\end{enumerate}
\end{exercice}

\begin{exercice}[Lutte biologique en cacaoyère]
Dans une cacaoyère de la région du Centre, les cabosses sont attaquées par des mirides (punaises) et les feuilles par des chenilles défoliatrices. Un technicien agricole propose d'introduire des fourmis \emph{Oecophylla} (fourmis tisserandes) et de favoriser la présence de coccinelles et de libellules aux abords de la plantation.
\begin{enumerate}
\item Définir le terme \og insecte auxiliaire \fg{} et préciser le rôle de chacun des trois insectes cités.
\item Donner deux avantages de la lutte biologique par rapport aux insecticides chimiques de synthèse.
\item Donner une limite de cette méthode.
\end{enumerate}
\end{exercice}

\begin{exercice}[Conception d'un outil de sensibilisation]
Dans le cadre de la semaine de la nutrition de ton lycée, tu dois concevoir une affiche destinée aux élèves de seconde sur le thème : \og Manger des insectes : une solution d'avenir pour le Cameroun \fg.
\begin{enumerate}
\item Dégager trois messages clés à faire figurer sur l'affiche (un nutritionnel, un économique, un environnemental).
\item Proposer un slogan court et percutant.
\item Indiquer deux précautions d'hygiène à mentionner concernant la collecte et la préparation des insectes comestibles.
\end{enumerate}
\end{exercice}

\courssub{Je me prépare au Bac}

\begin{exercice}[Type Bac D : couverture des AJR par l'entomophagie]
\textbf{Document 1.} Composition (base matière sèche) de quelques sources de protéines disponibles au Cameroun.
\begin{center}
\begin{tabularx}{\linewidth}{|l|X|X|X|X|}\hline
\rowcolor{popblueL} \textbf{Aliment} & \textbf{Protéines (g/100 g MS)} & \textbf{Fer (mg/100 g MS)} & \textbf{Zinc (mg/100 g MS)} & \textbf{Prix (FCFA/kg)} \\ \hline
Termites séchés & 38 & 12 & 9 & 2\,500 \\ \hline
Vers de palmier séchés & 25 & 6 & 8 & 3\,000 \\ \hline
Sauterelles séchées & 60 & 10 & 14 & 2\,000 \\ \hline
Bœuf frais (MS = 30 \%) & 65 & 8 & 12 & 3\,000 \\ \hline
Poisson fumé & 55 & 5 & 6 & 4\,000 \\ \hline
\end{tabularx}
\end{center}
\textbf{Document 2.} AJR d'un adolescent de 15 ans : protéines \SI{52}{g/jour} ; fer \SI{12}{mg/jour} ; zinc \SI{11}{mg/jour}. On suppose que l'aliment considéré couvre seul la totalité de chaque AJR.

\textbf{Questions :}
\begin{enumerate}
\item Calculer la masse de matière sèche de chaque aliment nécessaire pour couvrir l'AJR en protéines de cet adolescent. Présenter les résultats dans un tableau.
\item Sachant que le bœuf frais contient \SI{70}{\percent} d'eau, calculer la masse de bœuf \emph{frais} correspondant à la valeur trouvée à la question 1.
\item Calculer la masse de sauterelles séchées nécessaire pour couvrir l'AJR en fer, puis celle nécessaire pour l'AJR en zinc. Quelle masse faut-il retenir pour couvrir simultanément les deux AJR ? Justifier.
\item En utilisant les prix du document 1, comparer le coût journalier de la couverture de l'AJR en protéines par les sauterelles et par le poisson fumé.
\item À partir de l'ensemble des résultats et de tes connaissances, discuter en quelques lignes la faisabilité économique et l'acceptabilité sociale de l'entomophagie comme solution à l'insuffisance protéique au Cameroun (deux arguments pour, deux limites).
\end{enumerate}
\end{exercice}

\begin{exercice}[Type Bac D : optimisation d'un élevage de Tenebrio molitor]
Un groupe d'élèves de Terminale D élève des larves de \emph{Tenebrio molitor} (vers de farine) destinées à l'alimentation des poulets, sur quatre régimes alimentaires à \SI{27}{\celsius}. Chaque lot comporte 100 larves au départ ; l'expérience dure 8 semaines.

\textbf{Document 1.} Résultats obtenus après 8 semaines.
\begin{center}
\begin{tabularx}{\linewidth}{|l|X|X|X|X|}\hline
\rowcolor{popblueL} \textbf{Régime} & \textbf{Poids moyen final (mg/larve)} & \textbf{Taux de survie (\%)} & \textbf{Durée avant nymphose (semaines)} & \textbf{Teneur en protéines des larves (g/100 g MS)} \\ \hline
R1 : son de blé seul & 95 & 78 & 8 & 48 \\ \hline
R2 : son + carotte râpée & 130 & 92 & 6 & 50 \\ \hline
R3 : son + drêche de brasserie & 118 & 85 & 7 & 55 \\ \hline
R4 : son + déchets ménagers & 102 & 70 & 8 & 46 \\ \hline
\end{tabularx}
\end{center}
\textbf{Document 2.} Composition approximative des suppléments : carotte (eau \SI{88}{\percent}, sucres, provitamine A) ; drêche de brasserie (protéines \SI{25}{g/100 g} MS, fibres, vitamines du groupe B) ; déchets ménagers (composition très variable, risque de contamination).

\textbf{Questions :}
\begin{enumerate}
\item Identifier, en justifiant par trois critères chiffrés du document 1, le régime le plus performant pour une production de larves destinée à l'alimentation animale.
\item Proposer une explication biochimique de la supériorité du régime R2 par rapport au régime R1 (rôle de l'eau, des vitamines et des glucides dans la croissance larvaire).
\item Expliquer pourquoi le régime R3 donne des larves plus riches en protéines que R1, alors que le poids final est inférieur à celui de R2.
\item Dégager deux raisons pour lesquelles le régime R4 doit être déconseillé malgré son faible coût.
\item Rédiger un protocole expérimental permettant de tester l'hypothèse : \og l'apport d'eau est le facteur limitant de la croissance des larves sur son seul \fg{} (préciser la variable testée, les variables contrôlées, le témoin et le critère de mesure).
\end{enumerate}
\end{exercice}

\begin{exercice}[Type Bac D : argumentation et schéma d'économie circulaire]
\textbf{Partie A. Argumentation.} Un article de presse affirme : \og L'entomoculture est la solution miracle à l'insécurité alimentaire au Cameroun. \fg{} En t'appuyant sur les données nutritionnelles, environnementales, économiques et sociologiques étudiées en cours, nuance cette affirmation en rédigeant un texte argumenté et structuré comportant trois arguments en faveur de cette affirmation et trois limites ou conditions de réussite.

\textbf{Partie B. Schéma annoté.} Une ferme intégrée de la périphérie de Bafoussam associe : un poulailler de 500 poulets, une unité de bioconversion par larves de mouche soldat noire (MSN) et une parcelle de maïs fourrager.
\begin{enumerate}
\item Représenter par un schéma fonctionnel annoté les flux de matière entre les trois compartiments (fientes et déchets du poulailler, larves produites, frass, grains de maïs, aliment du bétail). Utiliser des flèches légendées.
\item Indiquer sur le schéma les trois rôles joués par les insectes dans ce système : recycleur de déchets, source d'aliment protéique, producteur d'engrais organique.
\item Un éleveur voisin souhaite utiliser ses fientes de poules \emph{fraîches}, sans aucun traitement, comme substrat d'élevage des larves de MSN. Identifier trois risques sanitaires (microbiologiques ou chimiques) de cette pratique.
\item Proposer un protocole de prétraitement (hygiénisation) du substrat avant son utilisation, en précisant au moins deux étapes validées scientifiquement (compostage thermophile, séchage, contrôle de l'humidité) et leur justification.
\end{enumerate}
\end{exercice}

\newpage

\coursec{Corrigés des exercices}

\begin{corrige}[Exercice 1 — QCM : insectes et leurs rôles]
\begin{enumerate}
\item \textbf{Réponse b.} L'\cle{entomophagie} (du grec \emph{entomon}, insecte, et \emph{phagein}, manger) désigne la consommation d'insectes par l'Homme ou par les animaux d'élevage. L'élevage d'insectes est l'\cle{entomoculture} (a), la lutte contre les ravageurs est la lutte phytosanitaire ou biologique (c), et l'étude des insectes est l'\cle{entomologie} (d).
\item \textbf{Réponse a.} La coccinelle (\emph{Coccinellidae}), à l'état adulte comme larvaire, est un prédateur vorace des pucerons : une larve peut consommer jusqu'à \num{100} pucerons par jour. Le criquet (b) et le ver blanc du palmier (c) sont des insectes comestibles ou ravageurs, le moustique (d) est un vecteur de maladies.
\item \textbf{Réponse b.} Le miel est élaboré par les abeilles à partir du \cle{nectar} des fleurs, enrichi en enzymes (invertase, glucose-oxydase) et déshydraté dans la ruche. Le pollen sert surtout de source protéique et la propolis est une résine collectée sur les bourgeons.
\item \textbf{Réponse c.} Le \cle{frass} est le mélange de déjections larvaires et de résidus de substrat non consommés ; riche en azote, phosphore et potassium, il constitue un engrais organique valorisable en agriculture.
\end{enumerate}
\end{corrige}

\begin{corrige}[Exercice 2 — Vrai ou faux justifié]
\begin{enumerate}
\item \textbf{Vrai.} Sur matière sèche, les termites ailés contiennent environ \SIrange{35}{45}{g} de protéines pour \SI{100}{g} de matière sèche, tandis que le bœuf frais, riche en eau (environ \SI{70}{\percent}), n'apporte que \SIrange{18}{22}{g} de protéines pour \SI{100}{g} de viande fraîche. La comparaison doit toutefois être faite à base sèche égale pour être rigoureuse.
\item \textbf{Faux.} Le venin d'abeille (\cle{apitoxine}) contient la \cle{mélittine}, peptide à propriétés anti-inflammatoires et antibactériennes, utilisé en apithérapie contre certaines affections rhumatismales. Son emploi exige toutefois des précautions (risque d'allergie, choc anaphylactique).
\item \textbf{Vrai.} La libellule est un insecte \cle{auxiliaire} : ses larves aquatiques (naïades) dévorent les larves de moustiques, et l'adulte capture les moustiques en vol. Elle contribue ainsi à la limitation naturelle des vecteurs du paludisme.
\item \textbf{Vrai.} Les larves de \emph{Hermetia illucens} réalisent une \cle{bioconversion} : elles transforment des déchets organiques (fientes, résidus de cuisine) en biomasse larvaire contenant jusqu'à \SIrange{40}{45}{\percent} de protéines (base sèche), utilisables dans l'alimentation des volailles et des poissons.
\item \textbf{Faux.} L'élevage d'insectes émet beaucoup \emph{moins} de gaz à effet de serre que l'élevage bovin par kilogramme de protéine produite : les insectes sont ectothermes (pas de dépense énergétique de thermorégulation), ont un taux de conversion alimentaire élevé et ne produisent pas de méthane entérique comme les ruminants.
\end{enumerate}
\end{corrige}

\begin{corrige}[Exercice 3 — Définitions]
\begin{enumerate}
\item \textbf{Entomophagie} : pratique alimentaire consistant à consommer des insectes, par l'Homme ou par les animaux d'élevage, comme source de nutriments (protéines, lipides, minéraux).
\item \textbf{Entomoculture} : ensemble des techniques d'élevage contrôlé d'insectes (choix du substrat, température, humidité, densité) en vue de produire de la biomasse alimentaire, des produits dérivés ou des auxiliaires de lutte biologique.
\item \textbf{Lutte biologique par insectes auxiliaires} : méthode de protection des cultures consistant à introduire, favoriser ou conserver des insectes prédateurs ou parasitoïdes (coccinelles, libellules, fourmis tisserandes) qui limitent naturellement les populations de ravageurs.
\item \textbf{Bioconversion} : transformation, par des organismes vivants (ici des larves d'insectes), de déchets ou de substrats de faible valeur en biomasse riche en protéines et en lipides valorisables.
\item \textbf{Propolis} : substance résineuse collectée par les abeilles sur les bourgeons et l'écorce des arbres, utilisée dans la ruche comme agent d'étanchéité et d'asepsie, et par l'Homme pour ses propriétés antiseptiques et anti-inflammatoires.
\end{enumerate}
\end{corrige}

\begin{corrige}[Exercice 4 — Calcul direct : apport protéique des chenilles]
\begin{enumerate}
\item \SI{100}{g} de chenilles séchées apportent \SI{55}{g} de protéines. La masse $m$ nécessaire pour apporter \SI{52}{g} de protéines est :
\[ m = \frac{52 \times 100}{55} = \frac{5200}{55} \approx \SI{94,5}{g}. \]
Il faut donc environ \SI{94,5}{g} de chenilles séchées.
\item Pour le bœuf maigre à \SI{20}{g} de protéines pour \SI{100}{g} :
\[ m' = \frac{52 \times 100}{20} = \SI{260}{g}. \]
Il faut \SI{260}{g} de bœuf frais maigre.
\item Coût des chenilles : $\dfrac{\num{94,5}}{1000} \times 3500 \approx \SI{331}{FCFA}$. Coût du bœuf : $\dfrac{260}{1000} \times 3000 = \SI{780}{FCFA}$.
\textbf{Conclusion :} bien que le kilogramme de chenilles séchées soit plus cher que celui de bœuf, la couverture des AJR en protéines coûte environ \num{2,4} fois moins cher avec les chenilles (\SI{331}{FCFA} contre \SI{780}{FCFA}), car leur concentration en protéines (base sèche) est beaucoup plus élevée. L'entomophagie est donc économiquement compétitive pour l'apport protéique.
\end{enumerate}
\end{corrige}

\begin{corrige}[Exercice 5 — Analyse d'un tableau de composition]
\begin{enumerate}
\item Ordre décroissant de richesse en protéines (g/100 g MS) : \textbf{bœuf (65) $>$ sauterelles (60) $>$ termites (38) $>$ vers de palmier (25)}.
\item Calcul de l'énergie pour \SI{100}{g} de matière sèche ($E = 17\,m_P + 38\,m_L$) :
\begin{align*}
\text{Termites : } & 17 \times 38 + 38 \times 45 = 646 + 1710 = \SI{2356}{kJ} \\
\text{Vers de palmier : } & 17 \times 25 + 38 \times 60 = 425 + 2280 = \SI{2705}{kJ} \\
\text{Bœuf : } & 17 \times 65 + 38 \times 25 = 1105 + 950 = \SI{2055}{kJ} \\
\text{Sauterelles : } & 17 \times 60 + 38 \times 8 = 1020 + 304 = \SI{1324}{kJ}
\end{align*}
Les \textbf{vers de palmier} sont les plus énergétiques (\SI{2705}{kJ/100 g} MS), en raison de leur forte teneur en lipides (\SI{60}{g/100 g} MS), les lipides étant plus de deux fois plus énergétiques que les protéines (\SI{38}{kJ/g} contre \SI{17}{kJ/g}).
\item L'anémie ferriprive est due à une carence en fer : on conseillera les \textbf{termites ailés séchés}, qui présentent la plus forte teneur en fer (\SI{12}{mg/100 g} MS), devant les sauterelles (\SI{10}{mg}).
\item La teneur en eau varie énormément d'un aliment à l'autre (ex. \SI{70}{\percent} pour la viande fraîche, moins de \SI{10}{\percent} pour les insectes séchés). Comparer sur matière fraîche diluerait artificiellement les nutriments des aliments riches en eau. La comparaison sur \cle{matière sèche} élimine la variable « eau » et rend les teneurs réellement comparables.
\end{enumerate}
\end{corrige}

\begin{corrige}[Exercice 6 — Produits de la ruche et leurs utilisations]
\begin{enumerate}
\item Trois produits de la ruche :
\begin{itemize}
\item \textbf{Miel} : aliment énergétique, édulcorant naturel ; utilisé traditionnellement contre la toux, les maux de gorge et pour cicatriser les plaies.
\item \textbf{Propolis} : antiseptique naturel, employée en infusion ou en teinture contre les infections de la gorge et les affections cutanées.
\item \textbf{Gelée royale} (ou pollen) : tonifiant et complément nutritionnel traditionnellement utilisé contre la fatigue ; le pollen sert de complément protéique et vitaminique. (Le venin, utilisé en apithérapie antirhumatismale, est aussi acceptable.)
\end{itemize}
\item L'activité antibactérienne du miel s'explique notamment par :
\begin{itemize}
\item sa \textbf{forte concentration en sucres} (\SI{80}{\percent}), qui crée une pression osmotique élevée provoquant la plasmolyse (déshydratation) des bactéries ;
\item la présence de \textbf{peroxyde d'hydrogène} (\ce{H2O2}) produit par l'enzyme glucose-oxydase, à effet bactéricide. (On accepte aussi : acidité du miel, pH voisin de 4, défavorisant la croissance bactérienne.)
\end{itemize}
\item Le principal composé actif anti-inflammatoire du venin d'abeille est la \textbf{mélittine}. Précaution indispensable : vérifier l'absence d'\textbf{allergie} au venin avant tout traitement, car une piqûre ou une injection peut déclencher un \textbf{choc anaphylactique} potentiellement mortel ; l'apithérapie doit être pratiquée sous contrôle médical.
\end{enumerate}
\end{corrige}

\begin{corrige}[Exercice 7 — Lutte biologique en cacaoyère]
\begin{enumerate}
\item Un \textbf{insecte auxiliaire} est un insecte prédateur ou parasitoïde qui, en se nourrissant de ravageurs des cultures, limite naturellement leurs populations et les dégâts qu'ils causent.
\begin{itemize}
\item \textbf{Fourmis tisserandes \emph{Oecophylla}} : prédatrices généralistes, elles chassent et capturent les mirides (punaises des cabosses) et les chenilles défoliatrices présents sur le cacaoyer.
\item \textbf{Coccinelles} : prédatrices des pucerons et autres petits insectes phytophages aux abords de la plantation.
\item \textbf{Libellules} : adultes et larves consomment les moustiques et divers insectes volants, réduisant les populations de ravageurs et de vecteurs autour de la parcelle.
\end{itemize}
\item Avantages de la lutte biologique (deux parmi) :
\begin{itemize}
\item absence de \textbf{résidus chimiques toxiques} dans les récoltes, les sols et les nappes phréatiques (protection de l'environnement et de la santé du consommateur) ;
\item \textbf{coût réduit} à long terme (pas d'achat répété d'insecticides) et absence de phénomène de résistance des ravageurs ;
\item préservation de la \textbf{biodiversité} et des pollinisateurs utiles.
\end{itemize}
\item Limite (une parmi) : action \textbf{lente} et parfois incomplète (les auxiliaires n'éradient pas totalement le ravageur) ; efficacité dépendante des conditions climatiques ; difficulté d'élevage et de lâcher massif des auxiliaires ; risque que l'auxiliaire introduit devienne lui-même envahissant.
\end{enumerate}
\end{corrige}

\begin{corrige}[Exercice 8 — Conception d'un outil de sensibilisation]
\begin{enumerate}
\item Trois messages clés :
\begin{itemize}
\item \textbf{Nutritionnel} : \og Les insectes comestibles (termites, chenilles, sauterelles) sont aussi riches en protéines que la viande et apportent en plus du fer et du zinc, indispensables contre l'anémie. \fg
\item \textbf{Économique} : \og Disponibles localement et peu coûteux, les insectes séchés couvrent les besoins en protéines à moindre coût que le bœuf ou le poisson. \fg
\item \textbf{Environnemental} : \og Élever des insectes demande peu d'eau, peu d'espace et émet très peu de gaz à effet de serre : c'est l'élevage de demain. \fg
\end{itemize}
\item Slogan (exemple) : \og Petits mais costauds : les insectes nourrissent le Cameroun ! \fg{} ou \og Mangeons malin, mangeons insectes : protéines d'ici, santé pour tous. \fg
\item Précautions d'hygiène (deux parmi) :
\begin{itemize}
\item collecter les insectes sur des sites \textbf{non pollués} (loin des zones traitées aux pesticides et des décharges) et identifier correctement les espèces comestibles ;
\item \textbf{cuire ou griller} systématiquement les insectes avant consommation pour détruire les germes pathogènes et les parasites ;
\item se laver les mains et les ustensiles, conserver les insectes séchés à l'abri de l'humidité.
\end{itemize}
\end{enumerate}
\end{corrige}

\begin{corrige}[Exercice 9 — Type Bac D : couverture des AJR par l'entomophagie]
\textbf{Barème indicatif (sur 10 pts) :} Q1 : 3 pts (0,5 pt par calcul correct) ; Q2 : 1,5 pt ; Q3 : 2 pts ; Q4 : 1,5 pt ; Q5 : 2 pts.

\begin{enumerate}
\item Masse de matière sèche $m = \dfrac{52 \times 100}{t}$, où $t$ est la teneur en protéines (g/100 g MS) :
\begin{center}
\begin{tabularx}{\linewidth}{|l|X|X|}\hline
\rowcolor{popblueL} \textbf{Aliment} & \textbf{Calcul} & \textbf{Masse de MS (g)} \\ \hline
Termites séchés & $5200/38$ & 136,8 \\ \hline
Vers de palmier séchés & $5200/25$ & 208,0 \\ \hline
Sauterelles séchées & $5200/60$ & 86,7 \\ \hline
Bœuf (base sèche) & $5200/65$ & 80,0 \\ \hline
Poisson fumé & $5200/55$ & 94,5 \\ \hline
\end{tabularx}
\end{center}
\item Le bœuf frais contient \SI{70}{\percent} d'eau, donc \SI{30}{\percent} de matière sèche :
\[ m_{\text{frais}} = \frac{80}{0,30} \approx \SI{266,7}{g} \text{ de bœuf frais}. \]
\item Fer : les sauterelles apportent \SI{10}{mg} de fer pour \SI{100}{g} MS :
\[ m_{\text{Fe}} = \frac{12 \times 100}{10} = \SI{120}{g}. \]
Zinc : \SI{14}{mg} pour \SI{100}{g} MS :
\[ m_{\text{Zn}} = \frac{11 \times 100}{14} \approx \SI{78,6}{g}. \]
Pour couvrir \emph{simultanément} les deux AJR, il faut retenir la \textbf{masse la plus grande, soit \SI{120}{g}} : avec \SI{120}{g} de sauterelles, l'apport en zinc est de $120 \times 14/100 = \SI{16,8}{mg} > \SI{11}{mg}$, donc l'AJR en zinc est aussi couvert. La masse de \SI{78,6}{g} ne couvrirait pas l'AJR en fer.
\item Sauterelles : $\dfrac{\num{86,7}}{1000} \times 2000 \approx \SI{173}{FCFA/jour}$. Poisson fumé : $\dfrac{\num{94,5}}{1000} \times 4000 \approx \SI{378}{FCFA/jour}$. La couverture de l'AJR en protéines coûte environ \num{2,2} fois moins cher avec les sauterelles.
\item \textbf{Discussion.} \emph{Arguments pour :} (1) les insectes séchés sont très concentrés en protéines et en minéraux (fer, zinc), ce qui réduit les quantités à consommer ; (2) leur coût par gramme de protéine est inférieur à celui du poisson fumé et compétitif avec le bœuf ; (3) leur production est accessible aux ménages ruraux (collecte, petit élevage) et peu émettrice de gaz à effet de serre. \emph{Limites :} (1) l'acceptabilité sociale reste inégale (réticences culturelles ou religieuses, dégoût chez certaines populations urbaines) ; (2) la disponibilité saisonnière des insectes collectés, les risques sanitaires (contamination, allergènes) et l'absence de filière organisée freinent la généralisation. L'entomophagie est donc un \emph{complément} crédible, pas une solution unique.
\end{enumerate}
\textbf{Points de vigilance :} exiger la formule littérale avant l'application numérique ; ne pas oublier les unités ; à la question 3, justifier explicitement le choix de la masse maximale ; à la question 2, bien raisonner sur \SI{30}{\percent} de MS et non ajouter \SI{70}{\percent} à 80 g.
\end{corrige}

\begin{corrige}[Exercice 10 — Type Bac D : optimisation d'un élevage de \emph{Tenebrio molitor}]
\textbf{Barème indicatif (sur 10 pts) :} Q1 : 2 pts ; Q2 : 2 pts ; Q3 : 2 pts ; Q4 : 1,5 pt ; Q5 : 2,5 pts.

\begin{enumerate}
\item Le régime le plus performant est \textbf{R2 (son + carotte râpée)}, d'après trois critères chiffrés :
\begin{itemize}
\item poids moyen final le plus élevé : \SI{130}{mg/larve} (contre 95, 118 et \SI{102}{mg}) ;
\item taux de survie le plus élevé : \SI{92}{\percent} (contre 78, 85 et \SI{70}{\percent}) ;
\item cycle le plus court : nymphose en 6 semaines (contre 8, 7 et 8), donc rotation des lots plus rapide.
\end{itemize}
\item La carotte apporte environ \SI{88}{\percent} d'eau, des glucides et de la provitamine A. L'\textbf{eau} est indispensable aux réactions métaboliques (hydrolyses, milieu des réactions enzymatiques) et à l'excrétion ; le son seul est très sec, ce qui limite la croissance. Les \textbf{glucides} fournissent l'énergie (ATP par respiration cellulaire) nécessaire à la synthèse des protéines et à la croissance. Les \textbf{vitamines} (provitamine A, cofacteurs enzymatiques) permettent le bon fonctionnement des enzymes du métabolisme. Le régime R2 corrige donc les carences hydrique et vitaminique du régime R1.
\item La drêche de brasserie contient \SI{25}{g} de protéines pour \SI{100}{g} MS : le régime R3 est donc plus riche en \textbf{azote} que R1, ce qui permet aux larves de synthétiser et d'accumuler davantage de protéines corporelles (\SI{55}{g/100 g} MS contre 48). Mais la drêche, pauvre en eau et en glucides facilement utilisables comparée à la carotte, ne corrige pas le déficit hydrique et énergétique : la \textbf{croissance pondérale} (gain de masse totale) reste inférieure à celle de R2 (\SI{118}{mg} contre \SI{130}{mg}). La richesse en protéines (qualité) et le poids final (quantité) dépendent de nutriments différents.
\item R4 doit être déconseillé car :
\begin{itemize}
\item la composition des déchets ménagers est \textbf{très variable} et incontrôlable, ce qui rend la production irrégulière et les résultats non reproductibles (survie la plus faible : \SI{70}{\percent}) ;
\item ils présentent un \textbf{risque de contamination} microbienne (salmonelles, moisissures productrices de mycotoxines) et chimique (résidus de produits ménagers, métaux lourds), transmissible aux poulets puis au consommateur.
\end{itemize}
\item \textbf{Protocole expérimental :}
\begin{itemize}
\item \emph{Variable testée (facteur étudié) :} la teneur en eau du substrat — par exemple 4 lots de son de blé humidifiés à 0, 10, 20 et \SI{30}{\percent} d'eau.
\item \emph{Variables contrôlées (constantes) :} même souche et même âge des larves, même effectif (100 larves par lot), même température (\SI{27}{\celsius}), même humidité ambiante, même quantité de son, même durée (8 semaines), mêmes conditions de luminosité.
\item \emph{Lot témoin :} le lot sur son sec (0 \% d'eau ajoutée), équivalent au régime R1.
\item \emph{Critères de mesure :} poids moyen final des larves (mg), taux de survie (\%), durée avant nymphose ; répéter chaque condition au moins 3 fois et comparer les moyennes.
\item \emph{Résultat attendu si l'hypothèse est vraie :} la croissance augmente avec l'humidification jusqu'à un optimum, le lot témoin restant le moins performant.
\end{itemize}
\end{enumerate}
\textbf{Points de vigilance :} à la question 1, citer les \emph{valeurs chiffrées} et non des affirmations vagues ; au protocole, exiger explicitement le témoin, l'unicité de la variable testée et un critère de mesure quantifiable.
\end{corrige}

\begin{corrige}[Exercice 11 — Type Bac D : argumentation et schéma d'économie circulaire]
\textbf{Barème indicatif (sur 10 pts) :} Partie A : 4 pts ; Partie B Q1--Q2 : 3 pts ; Q3 : 1,5 pt ; Q4 : 1,5 pt.

\textbf{Partie A. Texte argumenté (corrigé type).}

L'affirmation \og l'entomoculture est la solution miracle à l'insécurité alimentaire au Cameroun \fg{} mérite d'être nuancée : elle repose sur des fondements réels, mais comporte des conditions de réussite.

\emph{Arguments en faveur.} (1) \textbf{Nutritionnel} : les insectes comestibles (termites, chenilles, sauterelles, larves de MSN) présentent des teneurs en protéines comparables ou supérieures à celles de la viande sur matière sèche (\SIrange{38}{60}{g/100 g} MS) et apportent fer et zinc, nutriments critiques contre l'anémie et les carences. (2) \textbf{Environnemental} : les insectes, ectothermes à fort taux de conversion alimentaire, exigent peu d'eau, peu de surface et émettent très peu de gaz à effet de serre ; ils valorisent en outre les déchets organiques par bioconversion. (3) \textbf{Économique} : l'entomoculture demande peu de capital de départ, peut être conduite à l'échelle familiale et produit à la fois un aliment protéique et un engrais organique (frass), créant un revenu complémentaire.

\emph{Limites et conditions.} (1) \textbf{Acceptabilité socioculturelle} : la consommation d'insectes se heurte à des réticences dans certaines communautés ; une éducation nutritionnelle prolongée est indispensable. (2) \textbf{Sécurité sanitaire} : sans hygiénisation des substrats ni maîtrise des espèces élevées, il existe des risques microbiologiques, chimiques et allergiques ; une réglementation et des contrôles qualité doivent encadrer la filière. (3) \textbf{Technicité et marché} : la production à grande échelle exige la maîtrise des paramètres d'élevage (température, humidité, alimentation) et l'existence de débouchés ; l'entomoculture ne peut donc être qu'un \emph{complément} des productions classiques, non une solution unique et miraculeuse.

\textbf{Partie B.}
\begin{enumerate}
\item Schéma fonctionnel des flux de matière :
\begin{popfigure}
\begin{center}
\begin{tikzpicture}[font=\small, >=Stealth, node distance=1.6cm]
\node[draw, rounded corners, fill=poporangeL, minimum width=3.4cm, minimum height=1.1cm, align=center] (poulet) {\textbf{Poulailler}\\ 500 poulets};
\node[draw, rounded corners, fill=popgreenL, minimum width=3.8cm, minimum height=1.1cm, align=center, right=3.6cm of poulet] (msn) {\textbf{Unité MSN}\\ bioconversion};
\node[draw, rounded corners, fill=popblueL, minimum width=3.4cm, minimum height=1.1cm, align=center, below=2.2cm of msn] (mais) {\textbf{Parcelle de}\\ maïs fourrager};
\draw[->, very thick, popdark] (poulet) -- node[above, align=center]{fientes et déchets\\ du poulailler (substrat)} (msn);
\draw[->, very thick, popdark] (msn.west) .. controls +(-1.5,0.8) and +(1.2,0.8) .. node[below, align=center]{larves séchées : aliment\\ protéique des poulets} (poulet.east);
\draw[->, very thick, popdark] (msn) -- node[right, align=center]{frass : engrais\\ organique (N, P, K)} (mais);
\draw[->, very thick, popdark] (mais.west) -| node[below left, align=center, pos=0.25]{grains de maïs :\\ aliment du bétail} (poulet.south);
\end{tikzpicture}
\end{center}
\legende{Schéma fonctionnel de la ferme intégrée de Bafoussam. Flèches : flux de matière entre les trois compartiments (poulailler, unité de bioconversion par larves de mouche soldat noire, parcelle de maïs fourrager).}
\end{popfigure}
\item Les trois rôles des insectes (à annoter sur le compartiment « Unité MSN » du schéma) :
\begin{itemize}
\item \textbf{Recycleur de déchets} : les larves transforment les fientes et déchets du poulailler en biomasse (bioconversion) ;
\item \textbf{Source d'aliment protéique} : les larves séchées nourrissent les poulets ;
\item \textbf{Producteur d'engrais organique} : le frass fertilise la parcelle de maïs.
\end{itemize}
\item Trois risques sanitaires de l'utilisation de fientes fraîches non traitées :
\begin{itemize}
\item \textbf{risque microbiologique} : présence de bactéries pathogènes (\emph{Salmonella}, \emph{Escherichia coli}) et de parasites transmissibles aux larves puis aux poulets ;
\item \textbf{risque chimique} : résidus d'antibiotiques, de coccidiostatiques ou de métaux lourds (cuivre, zinc des aliments médicamentés) accumulés dans les fientes ;
\item \textbf{risque de fermentation et de toxines} : développement de moisissures productrices de mycotoxines et dégagement d'ammoniac dans un substrat frais trop humide.
\end{itemize}
\item Protocole de prétraitement (hygiénisation) du substrat :
\begin{itemize}
\item \textbf{Étape 1 — Compostage thermophile} : composter les fientes 2 à 3 semaines en andains retournés ; la température atteinte (\SIrange{55}{65}{\celsius}) détruit la majorité des bactéries pathogènes et des œufs de parasites.
\item \textbf{Étape 2 — Séchage et contrôle de l'humidité} : sécher le substrat puis le réhumidifier à environ \SIrange{60}{70}{\percent} d'humidité ; le séchage réduit les germes et les moisissures, et une humidité maîtrisée évite les fermentations anaérobies et le dégagement d'ammoniac.
\item (\emph{Complément accepté} : tamisage pour éliminer les corps étrangers, et analyse ou délai d'attente après tout traitement médicamenteux du poulailler pour limiter les résidus chimiques.)
\end{itemize}
\end{enumerate}
\textbf{Points de vigilance :} en Partie A, exiger un texte structuré (thèse, arguments, limites, conclusion nuancée) et non une liste ; sur le schéma, chaque flèche doit porter une légende indiquant la nature du flux ; à la question 4, chaque étape doit être \emph{justifiée} scientifiquement (effet de la température, de l'humidité).
\end{corrige}

\newpage

\coursec{Fiche bilan}

\begin{popfigure}
\begin{tikzpicture}[mindmap, grow cyclic, every node/.style={concept}, text width=2.6cm, align=center,
  root concept/.append style={concept color=popblue, font=\bfseries, text width=3.2cm},
  level 1 concept/.append style={level distance=4.6cm, sibling angle=51.4, font=\small}]
\node[root concept]{Amélioration de la production animale par les insectes}
  child[concept color=poporange]{ node {Contexte : insuffisance protéique, entomophagie} }
  child[concept color=popgreen]{ node {Valeurs nutritionnelles : termites, vers blancs, chenilles} }
  child[concept color=poppurple]{ node {Vertus thérapeutiques des insectes} }
  child[concept color=popgold]{ node {Produits de la ruche : miel, propolis, venins} }
  child[concept color=poppink]{ node {Lutte biologique : coccinelle, libellule} }
  child[concept color=popblue!70]{ node {Entomoculture : collecte et élevage} }
  child[concept color=popgreen!70!popblue]{ node {Sensibilisation et acceptabilité sociale} };
\end{tikzpicture}
\legende{Carte mentale de la leçon : les sept idées forces de l'exploitation des insectes pour améliorer la production animale et la sécurité alimentaire.}
\end{popfigure}

\begin{bilan}
\textbf{1. Définitions clés}
\begin{itemize}
  \item \cle{Entomophagie} : consommation d'insectes par l'Homme (ou les animaux d'élevage) comme source de nutriments.
  \item \cle{Entomoculture} : élevage contrôlé d'insectes comestibles ou utiles (vers de farine, termites, grillons) à des fins alimentaires ou thérapeutiques.
  \item \cle{Lutte biologique} : utilisation d'ennemis naturels (insectes auxiliaires prédateurs ou parasitoïdes) pour limiter les ravageurs des cultures.
  \item \cle{Insecte auxiliaire} : insecte utile qui détruit les ravageurs (coccinelle contre les pucerons, libellule contre les moustiques).
  \item \cle{Produit de la ruche} : substance issue des abeilles (miel, propolis, gelée royale, venin) à usage alimentaire ou thérapeutique.
\end{itemize}

\textbf{2. Données à connaître par c\oe ur}
\begin{center}
\begin{tabularx}{\linewidth}{|l|X|X|}
\hline
\rowcolor{popblueL} \textbf{Insecte / produit} & \textbf{Valeur nutritionnelle} & \textbf{Vertu thérapeutique / usage} \\
\hline
Termites ailés & Riches en protéines ($\approx$ 35--45 \% de la matière sèche) et en lipides & Apport énergétique ; lutte contre la malnutrition \\
\hline
Vers blancs (larves de coléoptères) & Protéines, lipides, fer & Traditionnellement utilisés contre certaines affections \\
\hline
Chenilles comestibles & Protéines, calcium, fer & Prévention de l'anémie et des carences \\
\hline
Miel & Glucides (glucose, fructose), vitamines, minéraux & Antiseptique, cicatrisant, adoucissant des voies respiratoires \\
\hline
Venin d'abeille & Substances bioactives (mélittine) & Apithérapie : rhumatismes, inflammations \\
\hline
Propolis & Flavonoïdes & Antibactérienne, antifongique \\
\hline
\end{tabularx}
\end{center}

\textbf{3. Lutte biologique : exemples types}
\begin{itemize}
  \item \cle{Coccinelle} : prédatrice des pucerons ravageurs des cultures vivrières et fourragères.
  \item \cle{Libellule} : prédatrice des moustiques (adultes et larves) : lutte indirecte contre le paludisme.
  \item Avantage : réduit l'usage des pesticides chimiques, protège les cultures fourragères, donc la production animale.
\end{itemize}

\textbf{4. Méthodes express}
\begin{itemize}
  \item \textbf{Justifier l'entomophagie} : comparer la teneur en protéines des insectes à celle de la viande classique $\rightarrow$ coût faible, disponibilité, faible impact environnemental.
  \item \textbf{Techniques d'utilisation} : collecte (termittières, troncs pourris), séchage, grillage, mouture en farine incorporée à l'alimentation humaine ou animale.
  \item \textbf{Concevoir un outil de sensibilisation} : affiche ou dépliant = message clair + données chiffrées + images locales + bénéfices nutritionnels et économiques.
\end{itemize}
\end{bilan}

\begin{aretenir}[Checklist avant le Bac]
\begin{itemize}
  \item $\square$ Je sais définir l'entomophagie, l'entomoculture et la lutte biologique.
  \item $\square$ Je cite au moins quatre insectes comestibles locaux (termites, vers blancs, chenilles, fourmis).
  \item $\square$ Je connais la richesse en protéines des termites ($\approx$ 35--45 \% de matière sèche).
  \item $\square$ J'explique l'apport de l'entomophagie face à l'insuffisance des ressources protéiques classiques.
  \item $\square$ Je cite les produits de la ruche : miel, propolis, gelée royale, venin.
  \item $\square$ Je connais les vertus thérapeutiques du miel (antiseptique, cicatrisant) et du venin d'abeille (apithérapie).
  \item $\square$ Je cite deux insectes auxiliaires et leur cible : coccinelle/pucerons, libellule/moustiques.
  \item $\square$ J'explique l'intérêt de la lutte biologique par rapport aux pesticides chimiques.
  \item $\square$ Je décris les étapes d'utilisation des insectes comestibles : collecte, séchage, grillage, mouture.
  \item $\square$ Je sais concevoir un outil de sensibilisation (affiche, dépliant) sur la valeur des insectes.
  \item $\square$ Je relie l'exploitation des insectes à l'amélioration de la production animale (farines d'insectes pour le bétail et la volaille).
  \item $\square$ Je rédige mes réponses avec un vocabulaire scientifique précis et des données chiffrées avec unités.
\end{itemize}
\end{aretenir}

\findecours

\end{document}
